% GENERATED FILE -- DO NOT EDIT.
%
% Produced by harness/assemble_contract.py from module.yaml and claims/*.tex.
% Edit the claim files and regenerate; `--check` fails if this file is stale.
\documentclass[11pt,a4paper]{article}
\pdftrailerid{80e4e513ecbaae79ed31eb37b37bab5f}
%% --- 1. Mathematics (order matters: physics after braket wins \ket, \tr) ---
\usepackage{amsmath}
\usepackage{amssymb}
\usepackage{amsfonts}
\usepackage{amsthm}
\usepackage{mathtools}
\usepackage{dsfont}
\usepackage{braket}
\usepackage{physics}
%% --- 2. Page furniture ------------------------------------------------------
\usepackage{xcolor}
\usepackage{array}    % >{...} column prefixes in the wide tables of some claims
\usepackage{enumitem}
\usepackage{graphicx}
\usepackage{booktabs}
\usepackage{longtable}
\usepackage{xltabular}              % front- and back-matter tables across pages
\usepackage{needspace}
\usepackage{xurl}
\usepackage{tcolorbox}
\tcbuselibrary{skins,breakable}
%% --- 3. Links and navigation (hyperref late; bookmark, refcount after it) ---
\usepackage{xr-hyper}   % before hyperref: the contract's labels
\usepackage{hyperref}
\usepackage{bookmark}
\usepackage{refcount}
%% --- 4. Colours: every distinction also has a shape (rule, caps, size, weight)
\definecolor{nnGrey}{gray}{0.40}      % non-normative text
\definecolor{metaGrey}{gray}{0.30}    % the generated meta line
\definecolor{techGrey}{gray}{0.45}    % heading of a technical claim
\definecolor{mainRule}{RGB}{0,110,140}
\definecolor{reqAccent}{RGB}{165,70,0}
\definecolor{linkBlue}{RGB}{0,62,150}
\definecolor{extLink}{RGB}{120,40,140}  % links into another PDF
\hypersetup{colorlinks=true, linkcolor=linkBlue, citecolor=linkBlue,
  urlcolor=linkBlue, filecolor=extLink, bookmarksnumbered=true,
  bookmarksopen=true, bookmarksopenlevel=1, pdfdisplaydoctitle=true}
% A link into another PDF (\href to a .pdf, \hyperref{file}..., xr-hyper) is a
% relative URI, <file>#nameddest=<destination>, in the file-link colour:
% browser viewers (Edge, Chrome) ignore hyperref's default GoToR action.
\makeatletter
\renewcommand*{\hyper@linkfile}[3]{% anchor text, file name, destination
  \begingroup
    \ifx\\#3\\\def\CQ@target{#2}\else\def\CQ@target{#2\hyper@hash nameddest=#3}\fi
    \let\@urlcolor\@filecolor
    \hyper@linkurl{#1}{\CQ@target}%
  \endgroup}
\makeatother
%% --- 5. The mapped Unicode set (harness/contract_navigation.py, UNICODE_TEX)
\DeclareUnicodeCharacter{03B1}{\ensuremath{\alpha}}
\DeclareUnicodeCharacter{03B2}{\ensuremath{\beta}}
\DeclareUnicodeCharacter{03B3}{\ensuremath{\gamma}}
\DeclareUnicodeCharacter{03B4}{\ensuremath{\delta}}
\DeclareUnicodeCharacter{03B5}{\ensuremath{\varepsilon}}
\DeclareUnicodeCharacter{03F5}{\ensuremath{\epsilon}}
\DeclareUnicodeCharacter{03B6}{\ensuremath{\zeta}}
\DeclareUnicodeCharacter{03B7}{\ensuremath{\eta}}
\DeclareUnicodeCharacter{03B8}{\ensuremath{\theta}}
\DeclareUnicodeCharacter{03D1}{\ensuremath{\vartheta}}
\DeclareUnicodeCharacter{03B9}{\ensuremath{\iota}}
\DeclareUnicodeCharacter{03BA}{\ensuremath{\kappa}}
\DeclareUnicodeCharacter{03BB}{\ensuremath{\lambda}}
\DeclareUnicodeCharacter{03BC}{\ensuremath{\mu}}
\DeclareUnicodeCharacter{03BD}{\ensuremath{\nu}}
\DeclareUnicodeCharacter{03BE}{\ensuremath{\xi}}
\DeclareUnicodeCharacter{03C0}{\ensuremath{\pi}}
\DeclareUnicodeCharacter{03C1}{\ensuremath{\rho}}
\DeclareUnicodeCharacter{03F1}{\ensuremath{\varrho}}
\DeclareUnicodeCharacter{03C3}{\ensuremath{\sigma}}
\DeclareUnicodeCharacter{03C2}{\ensuremath{\varsigma}}
\DeclareUnicodeCharacter{03C4}{\ensuremath{\tau}}
\DeclareUnicodeCharacter{03C5}{\ensuremath{\upsilon}}
\DeclareUnicodeCharacter{03C6}{\ensuremath{\varphi}}
\DeclareUnicodeCharacter{03D5}{\ensuremath{\phi}}
\DeclareUnicodeCharacter{03C7}{\ensuremath{\chi}}
\DeclareUnicodeCharacter{03C8}{\ensuremath{\psi}}
\DeclareUnicodeCharacter{03C9}{\ensuremath{\omega}}
\DeclareUnicodeCharacter{0393}{\ensuremath{\Gamma}}
\DeclareUnicodeCharacter{0394}{\ensuremath{\Delta}}
\DeclareUnicodeCharacter{0398}{\ensuremath{\Theta}}
\DeclareUnicodeCharacter{039B}{\ensuremath{\Lambda}}
\DeclareUnicodeCharacter{039E}{\ensuremath{\Xi}}
\DeclareUnicodeCharacter{03A0}{\ensuremath{\Pi}}
\DeclareUnicodeCharacter{03A3}{\ensuremath{\Sigma}}
\DeclareUnicodeCharacter{03A5}{\ensuremath{\Upsilon}}
\DeclareUnicodeCharacter{03A6}{\ensuremath{\Phi}}
\DeclareUnicodeCharacter{03A8}{\ensuremath{\Psi}}
\DeclareUnicodeCharacter{03A9}{\ensuremath{\Omega}}
\DeclareUnicodeCharacter{00E1}{\'a}
\DeclareUnicodeCharacter{00E0}{\`a}
\DeclareUnicodeCharacter{00E2}{\^a}
\DeclareUnicodeCharacter{00E4}{\"a}
\DeclareUnicodeCharacter{00E5}{\aa{}}
\DeclareUnicodeCharacter{00E9}{\'e}
\DeclareUnicodeCharacter{00E8}{\`e}
\DeclareUnicodeCharacter{00EA}{\^e}
\DeclareUnicodeCharacter{00EB}{\"e}
\DeclareUnicodeCharacter{00ED}{\'{\i}}
\DeclareUnicodeCharacter{00EF}{\"{\i}}
\DeclareUnicodeCharacter{00F3}{\'o}
\DeclareUnicodeCharacter{00F2}{\`o}
\DeclareUnicodeCharacter{00F4}{\^o}
\DeclareUnicodeCharacter{00F6}{\"o}
\DeclareUnicodeCharacter{00F8}{\o{}}
\DeclareUnicodeCharacter{00FA}{\'u}
\DeclareUnicodeCharacter{00F9}{\`u}
\DeclareUnicodeCharacter{00FB}{\^u}
\DeclareUnicodeCharacter{00FC}{\"u}
\DeclareUnicodeCharacter{00F1}{\~n}
\DeclareUnicodeCharacter{00E7}{\c{c}}
\DeclareUnicodeCharacter{00DF}{\ss{}}
\DeclareUnicodeCharacter{00C4}{\"A}
\DeclareUnicodeCharacter{00C5}{\AA{}}
\DeclareUnicodeCharacter{00C9}{\'E}
\DeclareUnicodeCharacter{00D6}{\"O}
\DeclareUnicodeCharacter{00DC}{\"U}
\DeclareUnicodeCharacter{2013}{\textendash{}}
\DeclareUnicodeCharacter{2014}{\textemdash{}}
\DeclareUnicodeCharacter{2018}{`}
\DeclareUnicodeCharacter{2019}{'}
\DeclareUnicodeCharacter{201C}{``}
\DeclareUnicodeCharacter{201D}{''}
\DeclareUnicodeCharacter{2026}{\ldots{}}
\DeclareUnicodeCharacter{00A0}{~}
\DeclareUnicodeCharacter{00A7}{\S{}}
\DeclareUnicodeCharacter{2264}{\ensuremath{\le}}
\DeclareUnicodeCharacter{2265}{\ensuremath{\ge}}
\DeclareUnicodeCharacter{2260}{\ensuremath{\neq}}
\DeclareUnicodeCharacter{2248}{\ensuremath{\approx}}
\DeclareUnicodeCharacter{2261}{\ensuremath{\equiv}}
\DeclareUnicodeCharacter{223C}{\ensuremath{\sim}}
\DeclareUnicodeCharacter{2243}{\ensuremath{\simeq}}
\DeclareUnicodeCharacter{2245}{\ensuremath{\cong}}
\DeclareUnicodeCharacter{221D}{\ensuremath{\propto}}
\DeclareUnicodeCharacter{226A}{\ensuremath{\ll}}
\DeclareUnicodeCharacter{226B}{\ensuremath{\gg}}
\DeclareUnicodeCharacter{2192}{\ensuremath{\to}}
\DeclareUnicodeCharacter{2190}{\ensuremath{\leftarrow}}
\DeclareUnicodeCharacter{2194}{\ensuremath{\leftrightarrow}}
\DeclareUnicodeCharacter{21D2}{\ensuremath{\Rightarrow}}
\DeclareUnicodeCharacter{21D0}{\ensuremath{\Leftarrow}}
\DeclareUnicodeCharacter{21D4}{\ensuremath{\Leftrightarrow}}
\DeclareUnicodeCharacter{21A6}{\ensuremath{\mapsto}}
\DeclareUnicodeCharacter{2208}{\ensuremath{\in}}
\DeclareUnicodeCharacter{2209}{\ensuremath{\notin}}
\DeclareUnicodeCharacter{2286}{\ensuremath{\subseteq}}
\DeclareUnicodeCharacter{2282}{\ensuremath{\subset}}
\DeclareUnicodeCharacter{2287}{\ensuremath{\supseteq}}
\DeclareUnicodeCharacter{2283}{\ensuremath{\supset}}
\DeclareUnicodeCharacter{2229}{\ensuremath{\cap}}
\DeclareUnicodeCharacter{222A}{\ensuremath{\cup}}
\DeclareUnicodeCharacter{2216}{\ensuremath{\setminus}}
\DeclareUnicodeCharacter{2205}{\ensuremath{\emptyset}}
\DeclareUnicodeCharacter{2200}{\ensuremath{\forall}}
\DeclareUnicodeCharacter{2203}{\ensuremath{\exists}}
\DeclareUnicodeCharacter{2115}{\ensuremath{\mathbb{N}}}
\DeclareUnicodeCharacter{2124}{\ensuremath{\mathbb{Z}}}
\DeclareUnicodeCharacter{211A}{\ensuremath{\mathbb{Q}}}
\DeclareUnicodeCharacter{211D}{\ensuremath{\mathbb{R}}}
\DeclareUnicodeCharacter{2102}{\ensuremath{\mathbb{C}}}
\DeclareUnicodeCharacter{00B7}{\ensuremath{\cdot}}
\DeclareUnicodeCharacter{00D7}{\ensuremath{\times}}
\DeclareUnicodeCharacter{2212}{\ensuremath{-}}
\DeclareUnicodeCharacter{00B1}{\ensuremath{\pm}}
\DeclareUnicodeCharacter{2213}{\ensuremath{\mp}}
\DeclareUnicodeCharacter{221A}{\ensuremath{\surd}}
\DeclareUnicodeCharacter{2211}{\ensuremath{\sum}}
\DeclareUnicodeCharacter{220F}{\ensuremath{\prod}}
\DeclareUnicodeCharacter{222B}{\ensuremath{\int}}
\DeclareUnicodeCharacter{221E}{\ensuremath{\infty}}
\DeclareUnicodeCharacter{2202}{\ensuremath{\partial}}
\DeclareUnicodeCharacter{2207}{\ensuremath{\nabla}}
\DeclareUnicodeCharacter{2297}{\ensuremath{\otimes}}
\DeclareUnicodeCharacter{2295}{\ensuremath{\oplus}}
\DeclareUnicodeCharacter{2218}{\ensuremath{\circ}}
\DeclareUnicodeCharacter{2308}{\ensuremath{\lceil}}
\DeclareUnicodeCharacter{2309}{\ensuremath{\rceil}}
\DeclareUnicodeCharacter{230A}{\ensuremath{\lfloor}}
\DeclareUnicodeCharacter{230B}{\ensuremath{\rfloor}}
\DeclareUnicodeCharacter{27E8}{\ensuremath{\langle}}
\DeclareUnicodeCharacter{27E9}{\ensuremath{\rangle}}
\DeclareUnicodeCharacter{2016}{\ensuremath{\|}}
\DeclareUnicodeCharacter{2113}{\ensuremath{\ell}}
\DeclareUnicodeCharacter{00B0}{\ensuremath{^{\circ}}}
\DeclareUnicodeCharacter{00B2}{\ensuremath{^{2}}}
\DeclareUnicodeCharacter{00B3}{\ensuremath{^{3}}}
\DeclareUnicodeCharacter{00B9}{\ensuremath{^{1}}}
\DeclareUnicodeCharacter{2080}{\ensuremath{_{0}}}
\DeclareUnicodeCharacter{2081}{\ensuremath{_{1}}}
\DeclareUnicodeCharacter{2082}{\ensuremath{_{2}}}
\DeclareUnicodeCharacter{207F}{\ensuremath{^{n}}}
\DeclareUnicodeCharacter{2096}{\ensuremath{_{k}}}
%% --- 6. Numbering: one counter per section -----------------------
\numberwithin{equation}{section}
\setcounter{tocdepth}{1}
%% --- 7. The project macro set, then the claim environments -----------------
% spec/macros.tex, read relative to this module's directory (latexmk -cd).
% =============================================================================
%  spec/macros.tex -- the one macro set every generated contract inputs.
%
%  Two parts: part A, the macros spec/notation.md catalogues, and part B,
%  the macros local to one module, grouped by module.
%  harness/assemble_contract.py emits % =============================================================================
%  spec/macros.tex -- the one macro set every generated contract inputs.
%
%  Two parts: part A, the macros spec/notation.md catalogues, and part B,
%  the macros local to one module, grouped by module.
%  harness/assemble_contract.py emits % =============================================================================
%  spec/macros.tex -- the one macro set every generated contract inputs.
%
%  Two parts: part A, the macros spec/notation.md catalogues, and part B,
%  the macros local to one module, grouped by module.
%  harness/assemble_contract.py emits \input{../../macros} for it, so the
%  file is read from spec/modules/<id>/.
%
%  RULES
%    * A macro that spec/notation.md catalogues is a verbatim copy of it
%      (part A). Changing one is a notation change: bump spec/notation.md.
%    * Part B holds the module-local macros. A macro defined by two modules
%      with the same expansion appears once, under the first module (in the
%      order below) that defines it. Two macros with the same name and
%      DIFFERENT expansions are both kept under a module prefix, and the
%      claims write the prefixed names:
%        \Chain  -> \redChain (reductions, \mathfrak{C}: the reduction chain)
%                   \pcChain  (protocol-compiler, \mathfrak{R}: the two-segment
%                             chain; \Chsrc and \Chrec are built on it)
%        \Nmsm   -> \redNmsm  (reductions, \mathcal{N}: a POVM measurement
%                             channel)
%                   \pcNmsm   (\mathcal{N}_{\mathsf A}: Alice's measurement
%                             channel, written by
%                             transmitter/thm:certificate-sufficiency,
%                             transmitter/rem:xreg and
%                             reductions/lem:source-transfer)
%    * Deliberately omitted: a \provenance{tag}{details} macro. Provenance
%      is a header field of every claim file (spec/claim-format.md section
%      3); a claim body that calls such a macro fails to compile, which is
%      the intended signal.
%
%  REQUIRED PACKAGES, loaded by the assembler's preamble before this file:
%    amsmath, amssymb, amsfonts, mathtools, dsfont, braket, physics, xcolor
%    (braket before physics, so that physics wins for \ket, \bra, \tr, ...).
%
%  This file contains only definitions: no \usepackage, no document body.
% =============================================================================


% =============================================================================
%  PART A -- notation macros, the ones spec/notation.md catalogues, each a
%  verbatim copy of its catalogue row; grouped under numbered headings.
% =============================================================================

% -----------------------------------------------------------------------------
% 1. State spaces and operator sets (a "system label" -> spaces of operators)
% -----------------------------------------------------------------------------
\newcommand{\Hs}[1]{\mathcal{H}(#1)}        % Hilbert space of a system label
\newcommand{\Lin}[1]{\mathcal{L}(#1)}       % bounded linear operators on H(#1)
\newcommand{\Linn}[2]{\mathcal{L}(#1,#2)}   % linear maps H(#1) -> H(#2)
\newcommand{\Pos}[1]{\mathrm{Pos}(#1)}      % positive semidefinite operators on H(#1)
\newcommand{\Density}{\mathrm{D}}           % density operators: \Density(A) = D(A)
\newcommand{\Seq}[1]{S_{=}(#1)}             % normalized states (trace = 1)
\newcommand{\Sleq}[1]{S_{\leq}(#1)}         % subnormalized states (0 <= trace <= 1)
\newcommand{\id}{\mathds{1}}                % identity operator
\newcommand{\idon}[1]{\mathds{1}_{#1}}      % identity on a named system, e.g. \idon{C}

% DEPRECATED: applies \mathcal to both arguments itself, so it takes bare
% letters, not system labels. Prefer \Linn; do not use in new contracts.
\providecommand{\Lop}[2]{\mathcal{L}(\mathcal{#1}, \mathcal{#2})}


% -----------------------------------------------------------------------------
% 2. Composite systems
% -----------------------------------------------------------------------------
\newcommand{\dsum}{\oplus}                  % direct sum of system labels
% Tensor product: use \otimes directly.
% n-fold direct sum of an operator space: \Lin{A}^{\oplus n}


% -----------------------------------------------------------------------------
% 3. Events, conditioning and classical registers (for STATES)
%    Convention:  \wedge  -> subnormalized "and-event" component, e.g. \rho_{A\wedge x}
%                 |        -> normalized conditional state,        e.g. \rho_{A|x}
% (Write these directly in subscripts; no extra macros needed.)
% -----------------------------------------------------------------------------


% -----------------------------------------------------------------------------
% 4. Photon-number subspaces / direct-sum blocks
%    Convention:  operators carry the restriction "on top", in brackets;
%                 states carry it as a \wedge subscript (see Sec. 3).
% -----------------------------------------------------------------------------
\newcommand{\Pnum}[1]{P_{#1}}               % photon-number projector (system from context)
\newcommand{\Pnumon}[2]{P_{#1}^{#2}}        % with explicit system superscript, e.g. P_n^C
\newcommand{\blk}[2]{#1^{(#2)}}             % operator block, e.g. \blk{M}{\leqN} = M^{(<=N)}
\newcommand{\leqN}{\leq N}                  % use as \blk{M}{\leqN}
\newcommand{\gtN}{> N}                      % use as \blk{M}{\gtN}
\newcommand{\eqn}{= n}                      % use as \blk{M}{\eqn}


% -----------------------------------------------------------------------------
% 5. System labels used in the protocol
% -----------------------------------------------------------------------------
\newcommand{\Chan}{C}                       % quantum-channel (Eve) input to Bob
\newcommand{\Prb}{P}                        % probe input
\newcommand{\Flag}{{\mathrm{Flg}}}          % flag system (EXISTS ONLY AFTER squashing)
\newcommand{\Loss}{{\mathrm{Lss}}}          % loss system
\newcommand{\Vc}{{\mathrm{Vcm}}}            % vacuum system
\newcommand{\Nout}{N_{\mathrm{out}}}
\newcommand{\Nin}{N_{\mathrm{in}}}


% -----------------------------------------------------------------------------
% 6. Alphabets
% -----------------------------------------------------------------------------
\newcommand{\Yset}{\mathcal{Y}}             % Bob's (coarse-grained) detection alphabet
\newcommand{\Yraw}{\mathcal{Y}_{\mathrm{raw}}}% raw click patterns {0,1}^4
\newcommand{\Yscset}{\mathcal{Y}_{\singlec}}% single-click outcomes
\newcommand{\Ymcset}{\mathcal{Y}_{\mc}}     % multi-click outcomes
\newcommand{\Tset}{\mathcal{T}_{\mathrm{set}}}% private setting set (the probe settings in the optical instance)
\newcommand{\Xset}{\mathcal{X}}             % Alice's alphabet
\newcommand{\Dset}{\{H,V,+,-\}}             % polarization / detector labels
\newcommand{\pass}{\top}                    % pass sifting


% -----------------------------------------------------------------------------
% 7. Mode / outcome abbreviations
% -----------------------------------------------------------------------------
\newcommand{\on}{{\mathrm{on}}}
\newcommand{\off}{{\mathrm{off}}}
\newcommand{\dc}{{\mathrm{dc}}}             % dark count
\newcommand{\sq}{\mathrm{sq}}               % squashed
\newcommand{\mc}{\mathrm{mc}}               % multi-click
\newcommand{\nc}{\mathrm{nc}}               % no-click
\newcommand{\singlec}{\mathrm{sc}}          % single-click
\newcommand{\vac}{\mathrm{vac}}             % vacuum
\newcommand{\kout}{{k_\mathrm{out}}}
\newcommand{\kin}{{k_\mathrm{in}}}


% -----------------------------------------------------------------------------
% 8. Maps, channels and miscellaneous operators
% -----------------------------------------------------------------------------
\newcommand{\Prob}{\mathrm{Prob}}           % probability (use \Prob[\cdot] or \Prob(\cdot))
\newcommand{\lmin}{\lambda_{\min}}          % minimum eigenvalue
\newcommand{\Wflag}{W_{\Flag}}              % flag-space weight
\newcommand{\transpose}{\mathsf{T}}
\newcommand{\innerp}[2]{\left\langle #1 \vert #2 \right\rangle}
\renewcommand{\vec}[1]{\boldsymbol{#1}}     % NB: overrides the LaTeX arrow accent


% -----------------------------------------------------------------------------
% 9. Key-rate numerics
% -----------------------------------------------------------------------------
\newcommand{\cG}{\mathcal{G}}               % protocol map
\newcommand{\cZ}{\mathcal{Z}}               % pinching map
\newcommand{\Amap}{\mathcal{A}_{\mathrm{ann}}} % announcement + measurement map
\newcommand{\Amapnb}{\mathcal{A}_{\mathrm{ann}|\mathrm{nb}}}
\newcommand{\simpl}{}                       % empty suffix for simplified block maps
\newcommand{\acSet}{\mathcal{S}_{\mathrm{acc}}} % acceptance set
\newcommand{\testE}{\mathcal{E}_{\mathrm{test}}} % test-statistics map
\newcommand{\ecLeak}{\mathrm{leak}^{\mathrm{EC}}_{\mathrm{obs}}}
\newcommand{\pgen}{p_{\mathrm{gen}}}
\newcommand{\askeyr}{R_{\infty}}             % asymptotic key rate
\newcommand{\ppass}{p_{\mathrm{pass}}}       % sifting pass probability


% -----------------------------------------------------------------------------
% 10. Finite-size analysis / Renyi entropies
% -----------------------------------------------------------------------------

% --- Round types: test vs. key generation (register T, alphabet \Fset) --------
\newcommand{\tst}{\mathrm{test}}            % test round
\newcommand{\gen}{\mathrm{gen}}             % key-generation round
\newcommand{\Fset}{\mathcal{T}}             % round-type alphabet {test, gen} (not \Tset)
\newcommand{\Chat}{\widehat{C}}             % acceptance-test register
\newcommand{\Chatset}{\widehat{\mathcal{C}}}% its alphabet

% --- Basis / key extraction functions ----------------------------------------
\newcommand{\bA}{b_{A}}                     % Alice's basis-extraction function
\newcommand{\bB}{b_{B}}                     % Bob's basis-extraction function
\newcommand{\kB}{k_{B}}                     % Bob's key-extraction function

% --- Maps --------------------------------------------------------------------
\newcommand{\Mprot}{\mathcal{M}}            % single-round protocol map
\newcommand{\Phimap}{\boldsymbol{\Phi}}     % test-statistics map
\newcommand{\Pdeph}{\mathcal{P}}            % dephasing (pinching) map on \Blreg

% --- Isometries --------------------------------------------------------------
\newcommand{\Vmsm}{V_{\mathrm{msm}}}        % measurement isometry
\newcommand{\Vann}{V_{\mathrm{ann}}}        % announcement isometry
\newcommand{\VI}{V_{I}}                     % Eve's copy of the announcements
\newcommand{\VChat}{V_{\Chat}}              % acceptance-data isometry
\newcommand{\VR}{V_{R}}                     % sifting + key map -> raw key R
\newcommand{\VZ}{V_{Z}}                     % raw key -> key register Z
\newcommand{\Wprot}{W}                      % full protocol isometry

% --- Renyi entropies ---------------------------------------------------------
% SECURITY-CRITICAL: see spec/notation.md §4. Never change an arrow to
% make an inequality work out; never convert silently between the sandwiched
% (tilde) and Petz (overline) families.
\newcommand{\Hsandup}{\widetilde{H}^{\uparrow}_{\alpha}}   % sandwiched, up
\newcommand{\Hsanddown}{\widetilde{H}^{\downarrow}_{\alpha}}% sandwiched, down
\newcommand{\Hpetzup}{\overline{H}^{\uparrow}_{\alpha}}    % Petz, up
\newcommand{\Hpetzdown}{\overline{H}^{\downarrow}_{\alpha}}% Petz, down
\newcommand{\hup}{h^{\uparrow}_{\alpha}}                   % single-round rate
\newcommand{\hqkd}{h^{\uparrow,\mathrm{QKD}}_{\alpha}}     % single-round QKD rate
\newcommand{\hlow}{\underline{h}_{\alpha}}                 % lower bound on \hqkd
% --- spec/notation.md 3.5 and 4.5 -------------------------------------------
% \hlb  is the first-order linearization lower bound on \hlow: the headline
%       key rate (the `first-order` tier of foundations/def:tiers). NOT \hlow
%       itself.
% \tnum is the instance's one declared numerical tolerance. Its
%       near-collision with the declared marginal \tau_A is disclosed in
%       spec/notation.md 3.5; bare \tau is reserved and denotes neither.
\newcommand{\hlb}{\underline{h}^{\mathrm{FO}}_{\alpha}}    % first-order lower bound on \hlow
\newcommand{\tnum}{\tau_{\mathrm{num}}}                    % declared numerical tolerance


% =============================================================================
%  PART B -- module-local macros, grouped by module. Deprecated aliases,
%  defined so that a text written with them compiles, used in no new contract: \Lop,
%  \testE, \Pdeph, \pass, \ecLeak (spec/notation.md section 3.9).
% =============================================================================

% --- meat-core -----------------------------------------------------------------
\newcommand{\SecSet}{\mathbb{S}}            % sector index set of Q
\newcommand{\Qsys}{Q}                       % declared sector-carrying measured system
\newcommand{\Vset}{\mathcal{V}}             % declared joint outcome alphabet of the Q-POVM
\newcommand{\CAset}{\mathcal{C}_{A}}        % Alice announcement alphabet
\newcommand{\CBset}{\mathcal{C}_{B}}        % Bob announcement alphabet
\newcommand{\Party}{\mathfrak{P}}           % declared party partition
\newcommand{\Oset}[1]{\mathcal{O}_{#1}}     % readout alphabet of party #1
\newcommand{\rd}[1]{\pi_{#1}}               % readout map of party #1
\newcommand{\Sadm}{\mathcal{S}_{A}}         % admissible Alice subspace
\newcommand{\nomarg}{\mathsf{none}}         % "no trusted marginal declared"
\newcommand{\Zset}{\mathcal{Z}_{\perp}}     % raw-key alphabet incl. \perp
\newcommand{\Zval}{\mathcal{Z}}             % key alphabet without \perp
\newcommand{\Etil}{\widetilde{E}}           % single-round purifying register
\newcommand{\Ehat}{\widehat{E}}             % Eve's n-round register
\newcommand{\OmAT}{\Omega_{\mathrm{AT}}}    % acceptance event
\newcommand{\epa}{\varepsilon_{\mathrm{PA}}}
\newcommand{\eev}{\varepsilon_{\mathrm{EV}}}
\newcommand{\lEC}{\lambda_{\mathrm{EC}}}
\newcommand{\Decl}{\mathfrak{D}}            % protocol declaration
\newcommand{\CList}{\mathfrak{L}}           % declared constraint list
\newcommand{\PinS}{\mathcal{P}_{\SecSet}}   % sector pinching on Q
\newcommand{\Feas}{\mathcal{W}}             % feasible set of block states
\newcommand{\Chkey}{\Chatset_{\mathrm{key}}}% key-live acceptance values
\newcommand{\Kann}{\mathcal{K}}             % announcement-class alphabet
\newcommand{\hold}{\mathsf{h}}              % declared key holder (A or B)
\newcommand{\sfD}{D}                        % classical KL divergence (base 2)
\newcommand{\Dtwo}{D}                       % quantum relative entropy (base 2)
\newcommand{\Dsand}{\widetilde{D}_{\alpha}} % sandwiched Renyi divergence
\newcommand{\idmap}{\mathrm{id}}            % identity channel
\newcommand{\Aev}{\mathcal{A}}              % Eve's per-round attack channel
\newcommand{\AttSet}{\mathfrak{A}}          % licensed attack-channel set
\newcommand{\Chacc}{\Chatset_{\mathrm{acc}}}% acceptance support of \acSet
\newcommand{\epsstar}{\varepsilon^{\star}}      % distance to an exact POVM
\newcommand{\dfH}{\Delta_{\mathrm{herm}}}       % Hermiticity defect
\newcommand{\dfP}{\Delta_{\mathrm{pos}}}        % negative part
\newcommand{\dfC}{\Delta_{\mathrm{cpl}}}        % completeness defect
\newcommand{\lmoq}{g}                       % acceptance-side gradient block
\newcommand{\lmoJ}{G}                       % feasible-side gradient block
\newcommand{\bacc}{\beta_{\mathrm{acc}}}    % certified acceptance-side minimum
\newcommand{\Itset}{\mathfrak{I}}           % finite index set of iterates (C22)

% --- reductions ----------------------------------------------------------------
\newcommand{\Bsys}{B}                        % physical (pre-squashing) system
\newcommand{\Tcl}[1]{\mathcal{T}(#1)}        % trace-class operators on H(#1)
\newcommand{\Ins}{\mathcal{I}}               % generalized-instrument branch map
\newcommand{\redNmsm}{\mathcal{N}}              % POVM measurement channel
\newcommand{\Eprot}{\mathcal{E}}             % n-round protocol channel
\newcommand{\Grnd}{\mathsf{G}}               % one round map of a protocol
\newcommand{\Kfin}{\mathsf{K}}               % end-of-protocol processing
\newcommand{\Crit}{\mathsf{sec}}             % value of a security criterion
\newcommand{\Pflag}{\Pi_{\Flag}}             % flag-sector projector on Q
\newcommand{\Pleq}{\Pi_{\leqN}}              % preserved-block projector
\newcommand{\Pgt}{\Pi_{\gtN}}                % above-cutoff projector
\newcommand{\Sph}{S^{\mathrm{ph}}}           % family of physical inputs
\newcommand{\Srch}{S^{\mathrm{rch}}}         % reachable squashed states
\newcommand{\Ssq}{S^{\sq}}                   % squashed admissible set
\newcommand{\Ffss}{F^{\mathrm{FSS}}}         % flag-state-squasher target
\newcommand{\Fwp}{F^{\mathrm{WP}}}           % multi-bound WPFSS target
\newcommand{\Mwp}{M^{\mathrm{WP}}}           % renormalized above-cutoff POVM
\newcommand{\Mdec}{{M'}^{\mathrm{ph}}}       % declared physical family on B
\newcommand{\Ppvm}{P^{\Flag}}                % flag PVM
\newcommand{\Lmp}{{\Lambda_{\mathrm{mp}}}}   % measure-and-prepare channel
\newcommand{\Lwp}{{\Lambda^{\mathrm{WP}}}}   % multi-bound WPFSS channel
\newcommand{\Lfss}{{\Lambda^{\mathrm{FSS}}}} % flag-state-squasher channel
\newcommand{\lmax}{\lambda_{\max}}           % maximum eigenvalue
\newcommand{\Lnoise}{\widetilde{\Lambda}}    % flag-noise channel
\newcommand{\cg}{\varpi}                     % declared coarse-graining map
\newcommand{\MCK}{\mathsf{MC}}               % multiclick raw-pattern set
\newcommand{\redChain}{\mathfrak{C}}            % reduction chain
\newcommand{\Lnk}{\ell}                      % a single typed link
\newcommand{\dsq}{\preceq}                      % the squashing relation
\newcommand{\opn}[1]{\left\lVert #1\right\rVert}           % operator (spectral) norm
\newcommand{\mxn}[1]{\left\lVert #1\right\rVert_{\max}}    % entrywise max modulus
\newcommand{\Fld}{\mathbb{F}}                   % exact ordered substrate field
\newcommand{\Fldi}{\mathbb{F}_{i}}              % its Gaussian extension F + iF
\newcommand{\lat}[1]{\big\lfloor #1\big\rfloor_{\kappa}} % round DOWN onto the lattice

% --- protocol-compiler ---------------------------------------------------------
\newcommand{\Spec}{\mathfrak{S}}             % protocol specification (compiler input)
\newcommand{\pcChain}{\mathfrak{R}}            % reduction chain (two typed segments)
\newcommand{\Chsrc}{\pcChain_{\mathrm{src}}}   % its Alice-side source-map segment
\newcommand{\Chrec}{\pcChain_{\mathrm{rec}}}   % its receiver-side squashing segment
\newcommand{\rvirt}{\rho^{\mathrm{v}}}       % virtual (PRE-source-map) signal state
\newcommand{\Fph}{F^{\mathrm{ph}}}           % physical (pre-reduction) POVM element
\newcommand{\Vph}{\mathcal{V}^{\mathrm{ph}}} % physical outcome alphabet
\newcommand{\Tpub}{\mathcal{T}_{\mathrm{pub}}}% public-choice alphabet
\newcommand{\Uset}{\mathcal{U}}              % residual outcome alphabet
\newcommand{\idxm}{\mathsf{idx}}             % label-to-index map
\newcommand{\Asig}{A'}                       % transmitted (signal) system -- the REAL one
\newcommand{\Avirt}{A''}                     % virtual transmitted system (source descriptor of Xi_1)
\newcommand{\Abar}{\bar{A}}                  % Alice's measured register
\newcommand{\Ahat}{\hat{A}}                  % Alice's shield register
\newcommand{\pcNmsm}{\mathcal{N}_{\mathsf A}}  % Alice's measurement channel (see the RULES block)
\newcommand{\Phip}{\Phi^{+}}                 % maximally entangled vector label
\newcommand{\Smap}{\mathfrak{s}}             % the emitted sector map (fraktur: a compile-time object)

% --- receiver ------------------------------------------------------------------
\newcommand{\Modes}{\mathbb{M}}              % declared ordered mode set
\newcommand{\Inp}{\mathbb{I}}                % declared input modes
\newcommand{\InpC}{\mathbb{I}_{\Chan}}       % channel input modes
\newcommand{\InpS}{\mathbb{I}_{\Prb}}        % probe (source) input modes
\newcommand{\Env}{\mathbb{E}}                % dilation (environment) modes
\newcommand{\Port}[1]{\mathbb{P}_{#1}}       % mode set of detector port #1
\newcommand{\Emd}[1]{E_{#1}}                 % coordinate projector on modes #1
\newcommand{\Gmat}{G}                        % instrument (transfer) matrix
\newcommand{\Fock}{\mathfrak{F}}             % Fock / second-quantization functor
\newcommand{\Symn}[1]{\mathrm{Sym}^{#1}}     % symmetric tensor power
\newcommand{\normord}[1]{{:}\mspace{2mu}#1\mspace{2mu}{:}}   % normal ordering
\newcommand{\rk}{\operatorname{rank}}        % rank (physics.sty owns \rank)
\newcommand{\SCK}{\mathsf{SC}}               % single-click pattern set
\newcommand{\Recv}{\mathfrak{R}}             % receiver instrument tuple
\newcommand{\Cat}{\mathfrak{C}}              % component catalogue
\newcommand{\Grph}{\mathfrak{E}}             % declared element graph
\newcommand{\defic}{\mathbb{L}}              % loss deficiency operator
\newcommand{\Pvac}[1]{\Pi_{\vac}(#1)}        % joint no-click element on ports #1
\newcommand{\pmc}{p_{\mc}}                   % uniform high-photon mc weight
\newcommand{\Bnd}{\beta}                     % the scalar bound function beta(n)
\newcommand{\Rot}{R}                         % planar rotation matrix
\newcommand{\Wplate}{W}                      % retarder (wave-plate) matrix
\newcommand{\Bsp}{B}                         % beamsplitter core matrix

% --- statistics ----------------------------------------------------------------
\newcommand{\honest}{p^{\mathrm{hon}}}          % honest per-round distribution
\newcommand{\ecomp}{\varepsilon_{\mathrm{comp}}}% completeness budget
\newcommand{\simplex}[1]{\Delta(#1)}            % probability simplex
\newcommand{\Chatcon}{\Chatset_{\mathrm{con}}}  % constrained outcomes
\newcommand{\Ireg}[3]{\mathcal{I}_{#1}(#2,#3)}  % regularized incomplete beta
\newcommand{\Ncount}[1]{\mathsf{N}_{#1}}        % per-outcome count
\newcommand{\qlo}{\mathbf{q}^{-}}               % box lower bounds
\newcommand{\qhi}{\mathbf{q}^{+}}               % box upper bounds
 for it, so the
%  file is read from spec/modules/<id>/.
%
%  RULES
%    * A macro that spec/notation.md catalogues is a verbatim copy of it
%      (part A). Changing one is a notation change: bump spec/notation.md.
%    * Part B holds the module-local macros. A macro defined by two modules
%      with the same expansion appears once, under the first module (in the
%      order below) that defines it. Two macros with the same name and
%      DIFFERENT expansions are both kept under a module prefix, and the
%      claims write the prefixed names:
%        \Chain  -> \redChain (reductions, \mathfrak{C}: the reduction chain)
%                   \pcChain  (protocol-compiler, \mathfrak{R}: the two-segment
%                             chain; \Chsrc and \Chrec are built on it)
%        \Nmsm   -> \redNmsm  (reductions, \mathcal{N}: a POVM measurement
%                             channel)
%                   \pcNmsm   (\mathcal{N}_{\mathsf A}: Alice's measurement
%                             channel, written by
%                             transmitter/thm:certificate-sufficiency,
%                             transmitter/rem:xreg and
%                             reductions/lem:source-transfer)
%    * Deliberately omitted: a \provenance{tag}{details} macro. Provenance
%      is a header field of every claim file (spec/claim-format.md section
%      3); a claim body that calls such a macro fails to compile, which is
%      the intended signal.
%
%  REQUIRED PACKAGES, loaded by the assembler's preamble before this file:
%    amsmath, amssymb, amsfonts, mathtools, dsfont, braket, physics, xcolor
%    (braket before physics, so that physics wins for \ket, \bra, \tr, ...).
%
%  This file contains only definitions: no \usepackage, no document body.
% =============================================================================


% =============================================================================
%  PART A -- notation macros, the ones spec/notation.md catalogues, each a
%  verbatim copy of its catalogue row; grouped under numbered headings.
% =============================================================================

% -----------------------------------------------------------------------------
% 1. State spaces and operator sets (a "system label" -> spaces of operators)
% -----------------------------------------------------------------------------
\newcommand{\Hs}[1]{\mathcal{H}(#1)}        % Hilbert space of a system label
\newcommand{\Lin}[1]{\mathcal{L}(#1)}       % bounded linear operators on H(#1)
\newcommand{\Linn}[2]{\mathcal{L}(#1,#2)}   % linear maps H(#1) -> H(#2)
\newcommand{\Pos}[1]{\mathrm{Pos}(#1)}      % positive semidefinite operators on H(#1)
\newcommand{\Density}{\mathrm{D}}           % density operators: \Density(A) = D(A)
\newcommand{\Seq}[1]{S_{=}(#1)}             % normalized states (trace = 1)
\newcommand{\Sleq}[1]{S_{\leq}(#1)}         % subnormalized states (0 <= trace <= 1)
\newcommand{\id}{\mathds{1}}                % identity operator
\newcommand{\idon}[1]{\mathds{1}_{#1}}      % identity on a named system, e.g. \idon{C}

% DEPRECATED: applies \mathcal to both arguments itself, so it takes bare
% letters, not system labels. Prefer \Linn; do not use in new contracts.
\providecommand{\Lop}[2]{\mathcal{L}(\mathcal{#1}, \mathcal{#2})}


% -----------------------------------------------------------------------------
% 2. Composite systems
% -----------------------------------------------------------------------------
\newcommand{\dsum}{\oplus}                  % direct sum of system labels
% Tensor product: use \otimes directly.
% n-fold direct sum of an operator space: \Lin{A}^{\oplus n}


% -----------------------------------------------------------------------------
% 3. Events, conditioning and classical registers (for STATES)
%    Convention:  \wedge  -> subnormalized "and-event" component, e.g. \rho_{A\wedge x}
%                 |        -> normalized conditional state,        e.g. \rho_{A|x}
% (Write these directly in subscripts; no extra macros needed.)
% -----------------------------------------------------------------------------


% -----------------------------------------------------------------------------
% 4. Photon-number subspaces / direct-sum blocks
%    Convention:  operators carry the restriction "on top", in brackets;
%                 states carry it as a \wedge subscript (see Sec. 3).
% -----------------------------------------------------------------------------
\newcommand{\Pnum}[1]{P_{#1}}               % photon-number projector (system from context)
\newcommand{\Pnumon}[2]{P_{#1}^{#2}}        % with explicit system superscript, e.g. P_n^C
\newcommand{\blk}[2]{#1^{(#2)}}             % operator block, e.g. \blk{M}{\leqN} = M^{(<=N)}
\newcommand{\leqN}{\leq N}                  % use as \blk{M}{\leqN}
\newcommand{\gtN}{> N}                      % use as \blk{M}{\gtN}
\newcommand{\eqn}{= n}                      % use as \blk{M}{\eqn}


% -----------------------------------------------------------------------------
% 5. System labels used in the protocol
% -----------------------------------------------------------------------------
\newcommand{\Chan}{C}                       % quantum-channel (Eve) input to Bob
\newcommand{\Prb}{P}                        % probe input
\newcommand{\Flag}{{\mathrm{Flg}}}          % flag system (EXISTS ONLY AFTER squashing)
\newcommand{\Loss}{{\mathrm{Lss}}}          % loss system
\newcommand{\Vc}{{\mathrm{Vcm}}}            % vacuum system
\newcommand{\Nout}{N_{\mathrm{out}}}
\newcommand{\Nin}{N_{\mathrm{in}}}


% -----------------------------------------------------------------------------
% 6. Alphabets
% -----------------------------------------------------------------------------
\newcommand{\Yset}{\mathcal{Y}}             % Bob's (coarse-grained) detection alphabet
\newcommand{\Yraw}{\mathcal{Y}_{\mathrm{raw}}}% raw click patterns {0,1}^4
\newcommand{\Yscset}{\mathcal{Y}_{\singlec}}% single-click outcomes
\newcommand{\Ymcset}{\mathcal{Y}_{\mc}}     % multi-click outcomes
\newcommand{\Tset}{\mathcal{T}_{\mathrm{set}}}% private setting set (the probe settings in the optical instance)
\newcommand{\Xset}{\mathcal{X}}             % Alice's alphabet
\newcommand{\Dset}{\{H,V,+,-\}}             % polarization / detector labels
\newcommand{\pass}{\top}                    % pass sifting


% -----------------------------------------------------------------------------
% 7. Mode / outcome abbreviations
% -----------------------------------------------------------------------------
\newcommand{\on}{{\mathrm{on}}}
\newcommand{\off}{{\mathrm{off}}}
\newcommand{\dc}{{\mathrm{dc}}}             % dark count
\newcommand{\sq}{\mathrm{sq}}               % squashed
\newcommand{\mc}{\mathrm{mc}}               % multi-click
\newcommand{\nc}{\mathrm{nc}}               % no-click
\newcommand{\singlec}{\mathrm{sc}}          % single-click
\newcommand{\vac}{\mathrm{vac}}             % vacuum
\newcommand{\kout}{{k_\mathrm{out}}}
\newcommand{\kin}{{k_\mathrm{in}}}


% -----------------------------------------------------------------------------
% 8. Maps, channels and miscellaneous operators
% -----------------------------------------------------------------------------
\newcommand{\Prob}{\mathrm{Prob}}           % probability (use \Prob[\cdot] or \Prob(\cdot))
\newcommand{\lmin}{\lambda_{\min}}          % minimum eigenvalue
\newcommand{\Wflag}{W_{\Flag}}              % flag-space weight
\newcommand{\transpose}{\mathsf{T}}
\newcommand{\innerp}[2]{\left\langle #1 \vert #2 \right\rangle}
\renewcommand{\vec}[1]{\boldsymbol{#1}}     % NB: overrides the LaTeX arrow accent


% -----------------------------------------------------------------------------
% 9. Key-rate numerics
% -----------------------------------------------------------------------------
\newcommand{\cG}{\mathcal{G}}               % protocol map
\newcommand{\cZ}{\mathcal{Z}}               % pinching map
\newcommand{\Amap}{\mathcal{A}_{\mathrm{ann}}} % announcement + measurement map
\newcommand{\Amapnb}{\mathcal{A}_{\mathrm{ann}|\mathrm{nb}}}
\newcommand{\simpl}{}                       % empty suffix for simplified block maps
\newcommand{\acSet}{\mathcal{S}_{\mathrm{acc}}} % acceptance set
\newcommand{\testE}{\mathcal{E}_{\mathrm{test}}} % test-statistics map
\newcommand{\ecLeak}{\mathrm{leak}^{\mathrm{EC}}_{\mathrm{obs}}}
\newcommand{\pgen}{p_{\mathrm{gen}}}
\newcommand{\askeyr}{R_{\infty}}             % asymptotic key rate
\newcommand{\ppass}{p_{\mathrm{pass}}}       % sifting pass probability


% -----------------------------------------------------------------------------
% 10. Finite-size analysis / Renyi entropies
% -----------------------------------------------------------------------------

% --- Round types: test vs. key generation (register T, alphabet \Fset) --------
\newcommand{\tst}{\mathrm{test}}            % test round
\newcommand{\gen}{\mathrm{gen}}             % key-generation round
\newcommand{\Fset}{\mathcal{T}}             % round-type alphabet {test, gen} (not \Tset)
\newcommand{\Chat}{\widehat{C}}             % acceptance-test register
\newcommand{\Chatset}{\widehat{\mathcal{C}}}% its alphabet

% --- Basis / key extraction functions ----------------------------------------
\newcommand{\bA}{b_{A}}                     % Alice's basis-extraction function
\newcommand{\bB}{b_{B}}                     % Bob's basis-extraction function
\newcommand{\kB}{k_{B}}                     % Bob's key-extraction function

% --- Maps --------------------------------------------------------------------
\newcommand{\Mprot}{\mathcal{M}}            % single-round protocol map
\newcommand{\Phimap}{\boldsymbol{\Phi}}     % test-statistics map
\newcommand{\Pdeph}{\mathcal{P}}            % dephasing (pinching) map on \Blreg

% --- Isometries --------------------------------------------------------------
\newcommand{\Vmsm}{V_{\mathrm{msm}}}        % measurement isometry
\newcommand{\Vann}{V_{\mathrm{ann}}}        % announcement isometry
\newcommand{\VI}{V_{I}}                     % Eve's copy of the announcements
\newcommand{\VChat}{V_{\Chat}}              % acceptance-data isometry
\newcommand{\VR}{V_{R}}                     % sifting + key map -> raw key R
\newcommand{\VZ}{V_{Z}}                     % raw key -> key register Z
\newcommand{\Wprot}{W}                      % full protocol isometry

% --- Renyi entropies ---------------------------------------------------------
% SECURITY-CRITICAL: see spec/notation.md §4. Never change an arrow to
% make an inequality work out; never convert silently between the sandwiched
% (tilde) and Petz (overline) families.
\newcommand{\Hsandup}{\widetilde{H}^{\uparrow}_{\alpha}}   % sandwiched, up
\newcommand{\Hsanddown}{\widetilde{H}^{\downarrow}_{\alpha}}% sandwiched, down
\newcommand{\Hpetzup}{\overline{H}^{\uparrow}_{\alpha}}    % Petz, up
\newcommand{\Hpetzdown}{\overline{H}^{\downarrow}_{\alpha}}% Petz, down
\newcommand{\hup}{h^{\uparrow}_{\alpha}}                   % single-round rate
\newcommand{\hqkd}{h^{\uparrow,\mathrm{QKD}}_{\alpha}}     % single-round QKD rate
\newcommand{\hlow}{\underline{h}_{\alpha}}                 % lower bound on \hqkd
% --- spec/notation.md 3.5 and 4.5 -------------------------------------------
% \hlb  is the first-order linearization lower bound on \hlow: the headline
%       key rate (the `first-order` tier of foundations/def:tiers). NOT \hlow
%       itself.
% \tnum is the instance's one declared numerical tolerance. Its
%       near-collision with the declared marginal \tau_A is disclosed in
%       spec/notation.md 3.5; bare \tau is reserved and denotes neither.
\newcommand{\hlb}{\underline{h}^{\mathrm{FO}}_{\alpha}}    % first-order lower bound on \hlow
\newcommand{\tnum}{\tau_{\mathrm{num}}}                    % declared numerical tolerance


% =============================================================================
%  PART B -- module-local macros, grouped by module. Deprecated aliases,
%  defined so that a text written with them compiles, used in no new contract: \Lop,
%  \testE, \Pdeph, \pass, \ecLeak (spec/notation.md section 3.9).
% =============================================================================

% --- meat-core -----------------------------------------------------------------
\newcommand{\SecSet}{\mathbb{S}}            % sector index set of Q
\newcommand{\Qsys}{Q}                       % declared sector-carrying measured system
\newcommand{\Vset}{\mathcal{V}}             % declared joint outcome alphabet of the Q-POVM
\newcommand{\CAset}{\mathcal{C}_{A}}        % Alice announcement alphabet
\newcommand{\CBset}{\mathcal{C}_{B}}        % Bob announcement alphabet
\newcommand{\Party}{\mathfrak{P}}           % declared party partition
\newcommand{\Oset}[1]{\mathcal{O}_{#1}}     % readout alphabet of party #1
\newcommand{\rd}[1]{\pi_{#1}}               % readout map of party #1
\newcommand{\Sadm}{\mathcal{S}_{A}}         % admissible Alice subspace
\newcommand{\nomarg}{\mathsf{none}}         % "no trusted marginal declared"
\newcommand{\Zset}{\mathcal{Z}_{\perp}}     % raw-key alphabet incl. \perp
\newcommand{\Zval}{\mathcal{Z}}             % key alphabet without \perp
\newcommand{\Etil}{\widetilde{E}}           % single-round purifying register
\newcommand{\Ehat}{\widehat{E}}             % Eve's n-round register
\newcommand{\OmAT}{\Omega_{\mathrm{AT}}}    % acceptance event
\newcommand{\epa}{\varepsilon_{\mathrm{PA}}}
\newcommand{\eev}{\varepsilon_{\mathrm{EV}}}
\newcommand{\lEC}{\lambda_{\mathrm{EC}}}
\newcommand{\Decl}{\mathfrak{D}}            % protocol declaration
\newcommand{\CList}{\mathfrak{L}}           % declared constraint list
\newcommand{\PinS}{\mathcal{P}_{\SecSet}}   % sector pinching on Q
\newcommand{\Feas}{\mathcal{W}}             % feasible set of block states
\newcommand{\Chkey}{\Chatset_{\mathrm{key}}}% key-live acceptance values
\newcommand{\Kann}{\mathcal{K}}             % announcement-class alphabet
\newcommand{\hold}{\mathsf{h}}              % declared key holder (A or B)
\newcommand{\sfD}{D}                        % classical KL divergence (base 2)
\newcommand{\Dtwo}{D}                       % quantum relative entropy (base 2)
\newcommand{\Dsand}{\widetilde{D}_{\alpha}} % sandwiched Renyi divergence
\newcommand{\idmap}{\mathrm{id}}            % identity channel
\newcommand{\Aev}{\mathcal{A}}              % Eve's per-round attack channel
\newcommand{\AttSet}{\mathfrak{A}}          % licensed attack-channel set
\newcommand{\Chacc}{\Chatset_{\mathrm{acc}}}% acceptance support of \acSet
\newcommand{\epsstar}{\varepsilon^{\star}}      % distance to an exact POVM
\newcommand{\dfH}{\Delta_{\mathrm{herm}}}       % Hermiticity defect
\newcommand{\dfP}{\Delta_{\mathrm{pos}}}        % negative part
\newcommand{\dfC}{\Delta_{\mathrm{cpl}}}        % completeness defect
\newcommand{\lmoq}{g}                       % acceptance-side gradient block
\newcommand{\lmoJ}{G}                       % feasible-side gradient block
\newcommand{\bacc}{\beta_{\mathrm{acc}}}    % certified acceptance-side minimum
\newcommand{\Itset}{\mathfrak{I}}           % finite index set of iterates (C22)

% --- reductions ----------------------------------------------------------------
\newcommand{\Bsys}{B}                        % physical (pre-squashing) system
\newcommand{\Tcl}[1]{\mathcal{T}(#1)}        % trace-class operators on H(#1)
\newcommand{\Ins}{\mathcal{I}}               % generalized-instrument branch map
\newcommand{\redNmsm}{\mathcal{N}}              % POVM measurement channel
\newcommand{\Eprot}{\mathcal{E}}             % n-round protocol channel
\newcommand{\Grnd}{\mathsf{G}}               % one round map of a protocol
\newcommand{\Kfin}{\mathsf{K}}               % end-of-protocol processing
\newcommand{\Crit}{\mathsf{sec}}             % value of a security criterion
\newcommand{\Pflag}{\Pi_{\Flag}}             % flag-sector projector on Q
\newcommand{\Pleq}{\Pi_{\leqN}}              % preserved-block projector
\newcommand{\Pgt}{\Pi_{\gtN}}                % above-cutoff projector
\newcommand{\Sph}{S^{\mathrm{ph}}}           % family of physical inputs
\newcommand{\Srch}{S^{\mathrm{rch}}}         % reachable squashed states
\newcommand{\Ssq}{S^{\sq}}                   % squashed admissible set
\newcommand{\Ffss}{F^{\mathrm{FSS}}}         % flag-state-squasher target
\newcommand{\Fwp}{F^{\mathrm{WP}}}           % multi-bound WPFSS target
\newcommand{\Mwp}{M^{\mathrm{WP}}}           % renormalized above-cutoff POVM
\newcommand{\Mdec}{{M'}^{\mathrm{ph}}}       % declared physical family on B
\newcommand{\Ppvm}{P^{\Flag}}                % flag PVM
\newcommand{\Lmp}{{\Lambda_{\mathrm{mp}}}}   % measure-and-prepare channel
\newcommand{\Lwp}{{\Lambda^{\mathrm{WP}}}}   % multi-bound WPFSS channel
\newcommand{\Lfss}{{\Lambda^{\mathrm{FSS}}}} % flag-state-squasher channel
\newcommand{\lmax}{\lambda_{\max}}           % maximum eigenvalue
\newcommand{\Lnoise}{\widetilde{\Lambda}}    % flag-noise channel
\newcommand{\cg}{\varpi}                     % declared coarse-graining map
\newcommand{\MCK}{\mathsf{MC}}               % multiclick raw-pattern set
\newcommand{\redChain}{\mathfrak{C}}            % reduction chain
\newcommand{\Lnk}{\ell}                      % a single typed link
\newcommand{\dsq}{\preceq}                      % the squashing relation
\newcommand{\opn}[1]{\left\lVert #1\right\rVert}           % operator (spectral) norm
\newcommand{\mxn}[1]{\left\lVert #1\right\rVert_{\max}}    % entrywise max modulus
\newcommand{\Fld}{\mathbb{F}}                   % exact ordered substrate field
\newcommand{\Fldi}{\mathbb{F}_{i}}              % its Gaussian extension F + iF
\newcommand{\lat}[1]{\big\lfloor #1\big\rfloor_{\kappa}} % round DOWN onto the lattice

% --- protocol-compiler ---------------------------------------------------------
\newcommand{\Spec}{\mathfrak{S}}             % protocol specification (compiler input)
\newcommand{\pcChain}{\mathfrak{R}}            % reduction chain (two typed segments)
\newcommand{\Chsrc}{\pcChain_{\mathrm{src}}}   % its Alice-side source-map segment
\newcommand{\Chrec}{\pcChain_{\mathrm{rec}}}   % its receiver-side squashing segment
\newcommand{\rvirt}{\rho^{\mathrm{v}}}       % virtual (PRE-source-map) signal state
\newcommand{\Fph}{F^{\mathrm{ph}}}           % physical (pre-reduction) POVM element
\newcommand{\Vph}{\mathcal{V}^{\mathrm{ph}}} % physical outcome alphabet
\newcommand{\Tpub}{\mathcal{T}_{\mathrm{pub}}}% public-choice alphabet
\newcommand{\Uset}{\mathcal{U}}              % residual outcome alphabet
\newcommand{\idxm}{\mathsf{idx}}             % label-to-index map
\newcommand{\Asig}{A'}                       % transmitted (signal) system -- the REAL one
\newcommand{\Avirt}{A''}                     % virtual transmitted system (source descriptor of Xi_1)
\newcommand{\Abar}{\bar{A}}                  % Alice's measured register
\newcommand{\Ahat}{\hat{A}}                  % Alice's shield register
\newcommand{\pcNmsm}{\mathcal{N}_{\mathsf A}}  % Alice's measurement channel (see the RULES block)
\newcommand{\Phip}{\Phi^{+}}                 % maximally entangled vector label
\newcommand{\Smap}{\mathfrak{s}}             % the emitted sector map (fraktur: a compile-time object)

% --- receiver ------------------------------------------------------------------
\newcommand{\Modes}{\mathbb{M}}              % declared ordered mode set
\newcommand{\Inp}{\mathbb{I}}                % declared input modes
\newcommand{\InpC}{\mathbb{I}_{\Chan}}       % channel input modes
\newcommand{\InpS}{\mathbb{I}_{\Prb}}        % probe (source) input modes
\newcommand{\Env}{\mathbb{E}}                % dilation (environment) modes
\newcommand{\Port}[1]{\mathbb{P}_{#1}}       % mode set of detector port #1
\newcommand{\Emd}[1]{E_{#1}}                 % coordinate projector on modes #1
\newcommand{\Gmat}{G}                        % instrument (transfer) matrix
\newcommand{\Fock}{\mathfrak{F}}             % Fock / second-quantization functor
\newcommand{\Symn}[1]{\mathrm{Sym}^{#1}}     % symmetric tensor power
\newcommand{\normord}[1]{{:}\mspace{2mu}#1\mspace{2mu}{:}}   % normal ordering
\newcommand{\rk}{\operatorname{rank}}        % rank (physics.sty owns \rank)
\newcommand{\SCK}{\mathsf{SC}}               % single-click pattern set
\newcommand{\Recv}{\mathfrak{R}}             % receiver instrument tuple
\newcommand{\Cat}{\mathfrak{C}}              % component catalogue
\newcommand{\Grph}{\mathfrak{E}}             % declared element graph
\newcommand{\defic}{\mathbb{L}}              % loss deficiency operator
\newcommand{\Pvac}[1]{\Pi_{\vac}(#1)}        % joint no-click element on ports #1
\newcommand{\pmc}{p_{\mc}}                   % uniform high-photon mc weight
\newcommand{\Bnd}{\beta}                     % the scalar bound function beta(n)
\newcommand{\Rot}{R}                         % planar rotation matrix
\newcommand{\Wplate}{W}                      % retarder (wave-plate) matrix
\newcommand{\Bsp}{B}                         % beamsplitter core matrix

% --- statistics ----------------------------------------------------------------
\newcommand{\honest}{p^{\mathrm{hon}}}          % honest per-round distribution
\newcommand{\ecomp}{\varepsilon_{\mathrm{comp}}}% completeness budget
\newcommand{\simplex}[1]{\Delta(#1)}            % probability simplex
\newcommand{\Chatcon}{\Chatset_{\mathrm{con}}}  % constrained outcomes
\newcommand{\Ireg}[3]{\mathcal{I}_{#1}(#2,#3)}  % regularized incomplete beta
\newcommand{\Ncount}[1]{\mathsf{N}_{#1}}        % per-outcome count
\newcommand{\qlo}{\mathbf{q}^{-}}               % box lower bounds
\newcommand{\qhi}{\mathbf{q}^{+}}               % box upper bounds
 for it, so the
%  file is read from spec/modules/<id>/.
%
%  RULES
%    * A macro that spec/notation.md catalogues is a verbatim copy of it
%      (part A). Changing one is a notation change: bump spec/notation.md.
%    * Part B holds the module-local macros. A macro defined by two modules
%      with the same expansion appears once, under the first module (in the
%      order below) that defines it. Two macros with the same name and
%      DIFFERENT expansions are both kept under a module prefix, and the
%      claims write the prefixed names:
%        \Chain  -> \redChain (reductions, \mathfrak{C}: the reduction chain)
%                   \pcChain  (protocol-compiler, \mathfrak{R}: the two-segment
%                             chain; \Chsrc and \Chrec are built on it)
%        \Nmsm   -> \redNmsm  (reductions, \mathcal{N}: a POVM measurement
%                             channel)
%                   \pcNmsm   (\mathcal{N}_{\mathsf A}: Alice's measurement
%                             channel, written by
%                             transmitter/thm:certificate-sufficiency,
%                             transmitter/rem:xreg and
%                             reductions/lem:source-transfer)
%    * Deliberately omitted: a \provenance{tag}{details} macro. Provenance
%      is a header field of every claim file (spec/claim-format.md section
%      3); a claim body that calls such a macro fails to compile, which is
%      the intended signal.
%
%  REQUIRED PACKAGES, loaded by the assembler's preamble before this file:
%    amsmath, amssymb, amsfonts, mathtools, dsfont, braket, physics, xcolor
%    (braket before physics, so that physics wins for \ket, \bra, \tr, ...).
%
%  This file contains only definitions: no \usepackage, no document body.
% =============================================================================


% =============================================================================
%  PART A -- notation macros, the ones spec/notation.md catalogues, each a
%  verbatim copy of its catalogue row; grouped under numbered headings.
% =============================================================================

% -----------------------------------------------------------------------------
% 1. State spaces and operator sets (a "system label" -> spaces of operators)
% -----------------------------------------------------------------------------
\newcommand{\Hs}[1]{\mathcal{H}(#1)}        % Hilbert space of a system label
\newcommand{\Lin}[1]{\mathcal{L}(#1)}       % bounded linear operators on H(#1)
\newcommand{\Linn}[2]{\mathcal{L}(#1,#2)}   % linear maps H(#1) -> H(#2)
\newcommand{\Pos}[1]{\mathrm{Pos}(#1)}      % positive semidefinite operators on H(#1)
\newcommand{\Density}{\mathrm{D}}           % density operators: \Density(A) = D(A)
\newcommand{\Seq}[1]{S_{=}(#1)}             % normalized states (trace = 1)
\newcommand{\Sleq}[1]{S_{\leq}(#1)}         % subnormalized states (0 <= trace <= 1)
\newcommand{\id}{\mathds{1}}                % identity operator
\newcommand{\idon}[1]{\mathds{1}_{#1}}      % identity on a named system, e.g. \idon{C}

% DEPRECATED: applies \mathcal to both arguments itself, so it takes bare
% letters, not system labels. Prefer \Linn; do not use in new contracts.
\providecommand{\Lop}[2]{\mathcal{L}(\mathcal{#1}, \mathcal{#2})}


% -----------------------------------------------------------------------------
% 2. Composite systems
% -----------------------------------------------------------------------------
\newcommand{\dsum}{\oplus}                  % direct sum of system labels
% Tensor product: use \otimes directly.
% n-fold direct sum of an operator space: \Lin{A}^{\oplus n}


% -----------------------------------------------------------------------------
% 3. Events, conditioning and classical registers (for STATES)
%    Convention:  \wedge  -> subnormalized "and-event" component, e.g. \rho_{A\wedge x}
%                 |        -> normalized conditional state,        e.g. \rho_{A|x}
% (Write these directly in subscripts; no extra macros needed.)
% -----------------------------------------------------------------------------


% -----------------------------------------------------------------------------
% 4. Photon-number subspaces / direct-sum blocks
%    Convention:  operators carry the restriction "on top", in brackets;
%                 states carry it as a \wedge subscript (see Sec. 3).
% -----------------------------------------------------------------------------
\newcommand{\Pnum}[1]{P_{#1}}               % photon-number projector (system from context)
\newcommand{\Pnumon}[2]{P_{#1}^{#2}}        % with explicit system superscript, e.g. P_n^C
\newcommand{\blk}[2]{#1^{(#2)}}             % operator block, e.g. \blk{M}{\leqN} = M^{(<=N)}
\newcommand{\leqN}{\leq N}                  % use as \blk{M}{\leqN}
\newcommand{\gtN}{> N}                      % use as \blk{M}{\gtN}
\newcommand{\eqn}{= n}                      % use as \blk{M}{\eqn}


% -----------------------------------------------------------------------------
% 5. System labels used in the protocol
% -----------------------------------------------------------------------------
\newcommand{\Chan}{C}                       % quantum-channel (Eve) input to Bob
\newcommand{\Prb}{P}                        % probe input
\newcommand{\Flag}{{\mathrm{Flg}}}          % flag system (EXISTS ONLY AFTER squashing)
\newcommand{\Loss}{{\mathrm{Lss}}}          % loss system
\newcommand{\Vc}{{\mathrm{Vcm}}}            % vacuum system
\newcommand{\Nout}{N_{\mathrm{out}}}
\newcommand{\Nin}{N_{\mathrm{in}}}


% -----------------------------------------------------------------------------
% 6. Alphabets
% -----------------------------------------------------------------------------
\newcommand{\Yset}{\mathcal{Y}}             % Bob's (coarse-grained) detection alphabet
\newcommand{\Yraw}{\mathcal{Y}_{\mathrm{raw}}}% raw click patterns {0,1}^4
\newcommand{\Yscset}{\mathcal{Y}_{\singlec}}% single-click outcomes
\newcommand{\Ymcset}{\mathcal{Y}_{\mc}}     % multi-click outcomes
\newcommand{\Tset}{\mathcal{T}_{\mathrm{set}}}% private setting set (the probe settings in the optical instance)
\newcommand{\Xset}{\mathcal{X}}             % Alice's alphabet
\newcommand{\Dset}{\{H,V,+,-\}}             % polarization / detector labels
\newcommand{\pass}{\top}                    % pass sifting


% -----------------------------------------------------------------------------
% 7. Mode / outcome abbreviations
% -----------------------------------------------------------------------------
\newcommand{\on}{{\mathrm{on}}}
\newcommand{\off}{{\mathrm{off}}}
\newcommand{\dc}{{\mathrm{dc}}}             % dark count
\newcommand{\sq}{\mathrm{sq}}               % squashed
\newcommand{\mc}{\mathrm{mc}}               % multi-click
\newcommand{\nc}{\mathrm{nc}}               % no-click
\newcommand{\singlec}{\mathrm{sc}}          % single-click
\newcommand{\vac}{\mathrm{vac}}             % vacuum
\newcommand{\kout}{{k_\mathrm{out}}}
\newcommand{\kin}{{k_\mathrm{in}}}


% -----------------------------------------------------------------------------
% 8. Maps, channels and miscellaneous operators
% -----------------------------------------------------------------------------
\newcommand{\Prob}{\mathrm{Prob}}           % probability (use \Prob[\cdot] or \Prob(\cdot))
\newcommand{\lmin}{\lambda_{\min}}          % minimum eigenvalue
\newcommand{\Wflag}{W_{\Flag}}              % flag-space weight
\newcommand{\transpose}{\mathsf{T}}
\newcommand{\innerp}[2]{\left\langle #1 \vert #2 \right\rangle}
\renewcommand{\vec}[1]{\boldsymbol{#1}}     % NB: overrides the LaTeX arrow accent


% -----------------------------------------------------------------------------
% 9. Key-rate numerics
% -----------------------------------------------------------------------------
\newcommand{\cG}{\mathcal{G}}               % protocol map
\newcommand{\cZ}{\mathcal{Z}}               % pinching map
\newcommand{\Amap}{\mathcal{A}_{\mathrm{ann}}} % announcement + measurement map
\newcommand{\Amapnb}{\mathcal{A}_{\mathrm{ann}|\mathrm{nb}}}
\newcommand{\simpl}{}                       % empty suffix for simplified block maps
\newcommand{\acSet}{\mathcal{S}_{\mathrm{acc}}} % acceptance set
\newcommand{\testE}{\mathcal{E}_{\mathrm{test}}} % test-statistics map
\newcommand{\ecLeak}{\mathrm{leak}^{\mathrm{EC}}_{\mathrm{obs}}}
\newcommand{\pgen}{p_{\mathrm{gen}}}
\newcommand{\askeyr}{R_{\infty}}             % asymptotic key rate
\newcommand{\ppass}{p_{\mathrm{pass}}}       % sifting pass probability


% -----------------------------------------------------------------------------
% 10. Finite-size analysis / Renyi entropies
% -----------------------------------------------------------------------------

% --- Round types: test vs. key generation (register T, alphabet \Fset) --------
\newcommand{\tst}{\mathrm{test}}            % test round
\newcommand{\gen}{\mathrm{gen}}             % key-generation round
\newcommand{\Fset}{\mathcal{T}}             % round-type alphabet {test, gen} (not \Tset)
\newcommand{\Chat}{\widehat{C}}             % acceptance-test register
\newcommand{\Chatset}{\widehat{\mathcal{C}}}% its alphabet

% --- Basis / key extraction functions ----------------------------------------
\newcommand{\bA}{b_{A}}                     % Alice's basis-extraction function
\newcommand{\bB}{b_{B}}                     % Bob's basis-extraction function
\newcommand{\kB}{k_{B}}                     % Bob's key-extraction function

% --- Maps --------------------------------------------------------------------
\newcommand{\Mprot}{\mathcal{M}}            % single-round protocol map
\newcommand{\Phimap}{\boldsymbol{\Phi}}     % test-statistics map
\newcommand{\Pdeph}{\mathcal{P}}            % dephasing (pinching) map on \Blreg

% --- Isometries --------------------------------------------------------------
\newcommand{\Vmsm}{V_{\mathrm{msm}}}        % measurement isometry
\newcommand{\Vann}{V_{\mathrm{ann}}}        % announcement isometry
\newcommand{\VI}{V_{I}}                     % Eve's copy of the announcements
\newcommand{\VChat}{V_{\Chat}}              % acceptance-data isometry
\newcommand{\VR}{V_{R}}                     % sifting + key map -> raw key R
\newcommand{\VZ}{V_{Z}}                     % raw key -> key register Z
\newcommand{\Wprot}{W}                      % full protocol isometry

% --- Renyi entropies ---------------------------------------------------------
% SECURITY-CRITICAL: see spec/notation.md §4. Never change an arrow to
% make an inequality work out; never convert silently between the sandwiched
% (tilde) and Petz (overline) families.
\newcommand{\Hsandup}{\widetilde{H}^{\uparrow}_{\alpha}}   % sandwiched, up
\newcommand{\Hsanddown}{\widetilde{H}^{\downarrow}_{\alpha}}% sandwiched, down
\newcommand{\Hpetzup}{\overline{H}^{\uparrow}_{\alpha}}    % Petz, up
\newcommand{\Hpetzdown}{\overline{H}^{\downarrow}_{\alpha}}% Petz, down
\newcommand{\hup}{h^{\uparrow}_{\alpha}}                   % single-round rate
\newcommand{\hqkd}{h^{\uparrow,\mathrm{QKD}}_{\alpha}}     % single-round QKD rate
\newcommand{\hlow}{\underline{h}_{\alpha}}                 % lower bound on \hqkd
% --- spec/notation.md 3.5 and 4.5 -------------------------------------------
% \hlb  is the first-order linearization lower bound on \hlow: the headline
%       key rate (the `first-order` tier of foundations/def:tiers). NOT \hlow
%       itself.
% \tnum is the instance's one declared numerical tolerance. Its
%       near-collision with the declared marginal \tau_A is disclosed in
%       spec/notation.md 3.5; bare \tau is reserved and denotes neither.
\newcommand{\hlb}{\underline{h}^{\mathrm{FO}}_{\alpha}}    % first-order lower bound on \hlow
\newcommand{\tnum}{\tau_{\mathrm{num}}}                    % declared numerical tolerance


% =============================================================================
%  PART B -- module-local macros, grouped by module. Deprecated aliases,
%  defined so that a text written with them compiles, used in no new contract: \Lop,
%  \testE, \Pdeph, \pass, \ecLeak (spec/notation.md section 3.9).
% =============================================================================

% --- meat-core -----------------------------------------------------------------
\newcommand{\SecSet}{\mathbb{S}}            % sector index set of Q
\newcommand{\Qsys}{Q}                       % declared sector-carrying measured system
\newcommand{\Vset}{\mathcal{V}}             % declared joint outcome alphabet of the Q-POVM
\newcommand{\CAset}{\mathcal{C}_{A}}        % Alice announcement alphabet
\newcommand{\CBset}{\mathcal{C}_{B}}        % Bob announcement alphabet
\newcommand{\Party}{\mathfrak{P}}           % declared party partition
\newcommand{\Oset}[1]{\mathcal{O}_{#1}}     % readout alphabet of party #1
\newcommand{\rd}[1]{\pi_{#1}}               % readout map of party #1
\newcommand{\Sadm}{\mathcal{S}_{A}}         % admissible Alice subspace
\newcommand{\nomarg}{\mathsf{none}}         % "no trusted marginal declared"
\newcommand{\Zset}{\mathcal{Z}_{\perp}}     % raw-key alphabet incl. \perp
\newcommand{\Zval}{\mathcal{Z}}             % key alphabet without \perp
\newcommand{\Etil}{\widetilde{E}}           % single-round purifying register
\newcommand{\Ehat}{\widehat{E}}             % Eve's n-round register
\newcommand{\OmAT}{\Omega_{\mathrm{AT}}}    % acceptance event
\newcommand{\epa}{\varepsilon_{\mathrm{PA}}}
\newcommand{\eev}{\varepsilon_{\mathrm{EV}}}
\newcommand{\lEC}{\lambda_{\mathrm{EC}}}
\newcommand{\Decl}{\mathfrak{D}}            % protocol declaration
\newcommand{\CList}{\mathfrak{L}}           % declared constraint list
\newcommand{\PinS}{\mathcal{P}_{\SecSet}}   % sector pinching on Q
\newcommand{\Feas}{\mathcal{W}}             % feasible set of block states
\newcommand{\Chkey}{\Chatset_{\mathrm{key}}}% key-live acceptance values
\newcommand{\Kann}{\mathcal{K}}             % announcement-class alphabet
\newcommand{\hold}{\mathsf{h}}              % declared key holder (A or B)
\newcommand{\sfD}{D}                        % classical KL divergence (base 2)
\newcommand{\Dtwo}{D}                       % quantum relative entropy (base 2)
\newcommand{\Dsand}{\widetilde{D}_{\alpha}} % sandwiched Renyi divergence
\newcommand{\idmap}{\mathrm{id}}            % identity channel
\newcommand{\Aev}{\mathcal{A}}              % Eve's per-round attack channel
\newcommand{\AttSet}{\mathfrak{A}}          % licensed attack-channel set
\newcommand{\Chacc}{\Chatset_{\mathrm{acc}}}% acceptance support of \acSet
\newcommand{\epsstar}{\varepsilon^{\star}}      % distance to an exact POVM
\newcommand{\dfH}{\Delta_{\mathrm{herm}}}       % Hermiticity defect
\newcommand{\dfP}{\Delta_{\mathrm{pos}}}        % negative part
\newcommand{\dfC}{\Delta_{\mathrm{cpl}}}        % completeness defect
\newcommand{\lmoq}{g}                       % acceptance-side gradient block
\newcommand{\lmoJ}{G}                       % feasible-side gradient block
\newcommand{\bacc}{\beta_{\mathrm{acc}}}    % certified acceptance-side minimum
\newcommand{\Itset}{\mathfrak{I}}           % finite index set of iterates (C22)

% --- reductions ----------------------------------------------------------------
\newcommand{\Bsys}{B}                        % physical (pre-squashing) system
\newcommand{\Tcl}[1]{\mathcal{T}(#1)}        % trace-class operators on H(#1)
\newcommand{\Ins}{\mathcal{I}}               % generalized-instrument branch map
\newcommand{\redNmsm}{\mathcal{N}}              % POVM measurement channel
\newcommand{\Eprot}{\mathcal{E}}             % n-round protocol channel
\newcommand{\Grnd}{\mathsf{G}}               % one round map of a protocol
\newcommand{\Kfin}{\mathsf{K}}               % end-of-protocol processing
\newcommand{\Crit}{\mathsf{sec}}             % value of a security criterion
\newcommand{\Pflag}{\Pi_{\Flag}}             % flag-sector projector on Q
\newcommand{\Pleq}{\Pi_{\leqN}}              % preserved-block projector
\newcommand{\Pgt}{\Pi_{\gtN}}                % above-cutoff projector
\newcommand{\Sph}{S^{\mathrm{ph}}}           % family of physical inputs
\newcommand{\Srch}{S^{\mathrm{rch}}}         % reachable squashed states
\newcommand{\Ssq}{S^{\sq}}                   % squashed admissible set
\newcommand{\Ffss}{F^{\mathrm{FSS}}}         % flag-state-squasher target
\newcommand{\Fwp}{F^{\mathrm{WP}}}           % multi-bound WPFSS target
\newcommand{\Mwp}{M^{\mathrm{WP}}}           % renormalized above-cutoff POVM
\newcommand{\Mdec}{{M'}^{\mathrm{ph}}}       % declared physical family on B
\newcommand{\Ppvm}{P^{\Flag}}                % flag PVM
\newcommand{\Lmp}{{\Lambda_{\mathrm{mp}}}}   % measure-and-prepare channel
\newcommand{\Lwp}{{\Lambda^{\mathrm{WP}}}}   % multi-bound WPFSS channel
\newcommand{\Lfss}{{\Lambda^{\mathrm{FSS}}}} % flag-state-squasher channel
\newcommand{\lmax}{\lambda_{\max}}           % maximum eigenvalue
\newcommand{\Lnoise}{\widetilde{\Lambda}}    % flag-noise channel
\newcommand{\cg}{\varpi}                     % declared coarse-graining map
\newcommand{\MCK}{\mathsf{MC}}               % multiclick raw-pattern set
\newcommand{\redChain}{\mathfrak{C}}            % reduction chain
\newcommand{\Lnk}{\ell}                      % a single typed link
\newcommand{\dsq}{\preceq}                      % the squashing relation
\newcommand{\opn}[1]{\left\lVert #1\right\rVert}           % operator (spectral) norm
\newcommand{\mxn}[1]{\left\lVert #1\right\rVert_{\max}}    % entrywise max modulus
\newcommand{\Fld}{\mathbb{F}}                   % exact ordered substrate field
\newcommand{\Fldi}{\mathbb{F}_{i}}              % its Gaussian extension F + iF
\newcommand{\lat}[1]{\big\lfloor #1\big\rfloor_{\kappa}} % round DOWN onto the lattice

% --- protocol-compiler ---------------------------------------------------------
\newcommand{\Spec}{\mathfrak{S}}             % protocol specification (compiler input)
\newcommand{\pcChain}{\mathfrak{R}}            % reduction chain (two typed segments)
\newcommand{\Chsrc}{\pcChain_{\mathrm{src}}}   % its Alice-side source-map segment
\newcommand{\Chrec}{\pcChain_{\mathrm{rec}}}   % its receiver-side squashing segment
\newcommand{\rvirt}{\rho^{\mathrm{v}}}       % virtual (PRE-source-map) signal state
\newcommand{\Fph}{F^{\mathrm{ph}}}           % physical (pre-reduction) POVM element
\newcommand{\Vph}{\mathcal{V}^{\mathrm{ph}}} % physical outcome alphabet
\newcommand{\Tpub}{\mathcal{T}_{\mathrm{pub}}}% public-choice alphabet
\newcommand{\Uset}{\mathcal{U}}              % residual outcome alphabet
\newcommand{\idxm}{\mathsf{idx}}             % label-to-index map
\newcommand{\Asig}{A'}                       % transmitted (signal) system -- the REAL one
\newcommand{\Avirt}{A''}                     % virtual transmitted system (source descriptor of Xi_1)
\newcommand{\Abar}{\bar{A}}                  % Alice's measured register
\newcommand{\Ahat}{\hat{A}}                  % Alice's shield register
\newcommand{\pcNmsm}{\mathcal{N}_{\mathsf A}}  % Alice's measurement channel (see the RULES block)
\newcommand{\Phip}{\Phi^{+}}                 % maximally entangled vector label
\newcommand{\Smap}{\mathfrak{s}}             % the emitted sector map (fraktur: a compile-time object)

% --- receiver ------------------------------------------------------------------
\newcommand{\Modes}{\mathbb{M}}              % declared ordered mode set
\newcommand{\Inp}{\mathbb{I}}                % declared input modes
\newcommand{\InpC}{\mathbb{I}_{\Chan}}       % channel input modes
\newcommand{\InpS}{\mathbb{I}_{\Prb}}        % probe (source) input modes
\newcommand{\Env}{\mathbb{E}}                % dilation (environment) modes
\newcommand{\Port}[1]{\mathbb{P}_{#1}}       % mode set of detector port #1
\newcommand{\Emd}[1]{E_{#1}}                 % coordinate projector on modes #1
\newcommand{\Gmat}{G}                        % instrument (transfer) matrix
\newcommand{\Fock}{\mathfrak{F}}             % Fock / second-quantization functor
\newcommand{\Symn}[1]{\mathrm{Sym}^{#1}}     % symmetric tensor power
\newcommand{\normord}[1]{{:}\mspace{2mu}#1\mspace{2mu}{:}}   % normal ordering
\newcommand{\rk}{\operatorname{rank}}        % rank (physics.sty owns \rank)
\newcommand{\SCK}{\mathsf{SC}}               % single-click pattern set
\newcommand{\Recv}{\mathfrak{R}}             % receiver instrument tuple
\newcommand{\Cat}{\mathfrak{C}}              % component catalogue
\newcommand{\Grph}{\mathfrak{E}}             % declared element graph
\newcommand{\defic}{\mathbb{L}}              % loss deficiency operator
\newcommand{\Pvac}[1]{\Pi_{\vac}(#1)}        % joint no-click element on ports #1
\newcommand{\pmc}{p_{\mc}}                   % uniform high-photon mc weight
\newcommand{\Bnd}{\beta}                     % the scalar bound function beta(n)
\newcommand{\Rot}{R}                         % planar rotation matrix
\newcommand{\Wplate}{W}                      % retarder (wave-plate) matrix
\newcommand{\Bsp}{B}                         % beamsplitter core matrix

% --- statistics ----------------------------------------------------------------
\newcommand{\honest}{p^{\mathrm{hon}}}          % honest per-round distribution
\newcommand{\ecomp}{\varepsilon_{\mathrm{comp}}}% completeness budget
\newcommand{\simplex}[1]{\Delta(#1)}            % probability simplex
\newcommand{\Chatcon}{\Chatset_{\mathrm{con}}}  % constrained outcomes
\newcommand{\Ireg}[3]{\mathcal{I}_{#1}(#2,#3)}  % regularized incomplete beta
\newcommand{\Ncount}[1]{\mathsf{N}_{#1}}        % per-outcome count
\newcommand{\qlo}{\mathbf{q}^{-}}               % box lower bounds
\newcommand{\qhi}{\mathbf{q}^{+}}               % box upper bounds

\makeatletter
\newcommand*{\ClaimHeadFont}{\bfseries}   % switched by prominence
\newcommand*{\ClaimNoteFont}{\mdseries}
\newcommand*{\ClaimShortId}{}
\newcommand*{\SetShortId}[1]{\renewcommand*{\ClaimShortId}{#1}}
\newcommand*{\RestatedName}{}
% Upright bodies throughout: italic is reserved for non-normative text.
% \@endparpenalty\@M keeps the meta line on the statement's page.
\newtheoremstyle{claim}{\topsep}{\topsep}{\normalfont\@endparpenalty\@M}{}%
  {\ClaimHeadFont}{.}{.5em}{\thmname{#1}\thmnumber{ #2}\thmnote{ {\ClaimNoteFont(#3)}}}
\newtheoremstyle{claimass}{\topsep}{\topsep}{\normalfont\@endparpenalty\@M}{}%
  {\ClaimHeadFont}{.}{.5em}{\thmname{#1}\thmnumber{ #2}%
  \ifx\ClaimShortId\@empty\else\ [\textsf{\ClaimShortId}]\fi%
  \thmnote{ {\ClaimNoteFont(#3)}}}
\newtheoremstyle{claimremark}{.5\topsep}{.5\topsep}{\small\@endparpenalty\@M}{}%
  {\itshape}{.}{.5em}{\thmname{#1}\thmnumber{ #2}\thmnote{ {\upshape(#3)}}}
\newtheoremstyle{claimreq}{\topsep}{\topsep}{\normalfont\@endparpenalty\@M}{}%
  {\color{reqAccent}\scshape}{.}{.5em}%
  {\thmname{#1}\thmnumber{ #2}\thmnote{ {\upshape(#3)}}}
\makeatother
\theoremstyle{claim}
\newtheorem{theorem}{Theorem}[section]
\newtheorem{lemma}[theorem]{Lemma}
\newtheorem{proposition}[theorem]{Proposition}
\newtheorem{corollary}[theorem]{Corollary}
\newtheorem{fact}[theorem]{Fact}
\newtheorem{definition}[theorem]{Definition}
\newtheorem*{restatedclaim}{\RestatedName}
\theoremstyle{claimass}
\newtheorem{assumption}[theorem]{Assumption}
\newtheorem*{restatedassumption}{\RestatedName}
\theoremstyle{claimremark}
\newtheorem{remark}[theorem]{Remark}
\newtheorem{example}[theorem]{Example}
\theoremstyle{claimreq}
\newtheorem{requirement}[theorem]{Requirement}
\newtheorem*{restatedrequirement}{\RestatedName}
%% --- 8. Prominence: main results a margin rule and a bold heading ----------
\newtcolorbox{mainclaim}{enhanced jigsaw, breakable, blanker,
  borderline west={1.5pt}{-8pt}{mainRule}, left=0pt, right=0pt, top=0pt, bottom=0pt,
  boxsep=0pt, before skip=\topsep, after skip=\topsep, parbox=false,
  before upper={\renewcommand*{\ClaimNoteFont}{\bfseries}}}
\newenvironment{technicalclaim}{\renewcommand*{\ClaimHeadFont}{\mdseries\color{techGrey}}%
  \renewcommand*{\ClaimNoteFont}{}}{}
%% --- 9. Non-normative text, the meta line, navigation ------------------------
\makeatletter
% Tie what follows to this paragraph: no page break, no list penalty before
% the next theorem head (the \@nobreak route a section heading uses).
\newcommand{\keepwithnext}{\par\nobreak\vspace{3pt}\@nobreaktrue
  \everypar{\if@nobreak\@nobreakfalse\clubpenalty\@M\else\clubpenalty\@clubpenalty\everypar{}\fi}}
\newcommand{\claimsummary}[2]{% #1 claim id (named destination), #2 summary
  \par\addvspace{\bigskipamount}\Needspace{6\baselineskip}%
  {\small\itshape\color{nnGrey}\noindent\hypertarget{claim-#1}{}#2\par}\keepwithnext}
\newcommand{\claimanchor}[1]{% a claim without a summary: the destination alone
  \par\addvspace{\bigskipamount}\Needspace{6\baselineskip}%
  \phantomsection\hypertarget{claim-#1}{}\keepwithnext}
\newcommand{\claimleadin}[1]{{\color{nnGrey}\noindent#1\par}\keepwithnext}
\newcommand{\claimmeta}[1]{\par\nobreak\vspace{2pt}%
  {\footnotesize\color{metaGrey}\raggedright\interlinepenalty\@M\noindent#1\par}%
  \addvspace{\medskipamount}}
% After an assumption: the sentence generated from its verification field; after
% any claim with one: its plain-text reading (user_statement). Both
% non-normative, both grey.
\newcommand{\claimverification}[1]{\par\nobreak\vspace{2pt}%
  {\footnotesize\color{metaGrey}\noindent#1\par}\nobreak}
\newcommand{\claimreading}[1]{\par\nobreak\vspace{2pt}%
  {\footnotesize\color{nnGrey}\noindent\emph{Reading:}\enspace#1\par}\nobreak}
\makeatother
\newcommand{\claimbookmark}[4]{% #1 id, #2 kind word, #3 label, #4 plain title
  \bookmark[dest=claim-#1,level=2]{#2 \getrefnumber{#3} #4}}
\newcommand{\metasep}{\unskip\enspace\textperiodcentered\enspace\ignorespaces}
\newcommand{\metalabel}[1]{\emph{#1}\enspace}
\newcommand{\sectionintro}[1]{{\color{nnGrey}\noindent#1\par}\smallskip}
\newcommand{\sectionimports}[1]{{\footnotesize\color{metaGrey}\raggedright\noindent\metalabel{Imports}#1\par}\smallskip}
\newenvironment{overview}{\par\color{nnGrey}\itshape\noindent\ignorespaces}{\par}
% References: the whole "Def. 2.1" is the link; ids are the fallback text of
% every link into another PDF (a relative URI to a named destination).
\newcommand{\mref}[2]{\hyperref[#2]{#1~\ref*{#2}}}
\newcommand{\aref}[2]{\hyperref[#2]{\textsf{#1}}}
\DeclareUrlCommand{\claimid}{\urlstyle{sf}\def\UrlBreaks{\do\/\do\:\do\-\do\.\do\_}%
  \def\UrlBigBreaks{}\def\UrlNoBreaks{}}
\newcommand{\extref}[2]{\hyperref{#1}{}{claim-#2}{\claimid{#2}}}
\newcommand{\extaref}[3]{\hyperref{#1}{}{claim-#2}{\textsf{#3}}}
\newcommand{\coderef}[2]{\href{#1}{\claimid{#2}}}
% Claim-body commands (spec/claim-format.md section 4): \xref{<full id>} prints
% the kind, title and id of a claim of another contract, linked to its stable
% destination in that contract's PDF (one \DeclareXref per cited claim below);
% \term[gloss]{phrase} marks a defined phrase for the index (the gloss feeds
% the Terms table); \tref[phrase]{text} is a use of a defined term. The
% assembler rewrites both in the body into \termdef (the definition's named
% destination term-<slug>) and \termref / \exttermref (a link to it, in this
% or another contract's PDF); the fallbacks below serve text it copies
% elsewhere, such as the digest. \iface is an interface key.
\makeatletter
\newcommand{\DeclareXref}[4]{% id, kind word, title, contract PDF
  \@namedef{xref@#1}{\hyperref{#4}{}{claim-#1}{#2 ``#3'' ({\small\claimid{#1}})}}}
\newcommand{\xref}[1]{\@ifundefined{xref@#1}%
  {\textbf{[unresolved \nolinkurl{#1}]}\@latex@warning{Unresolved xref #1}}%
  {\@nameuse{xref@#1}}}
\makeatother
\DeclareUrlCommand{\iface}{\urlstyle{tt}}
\newcommand{\term}[2][]{\emph{#2}\index{#2}}
\newcommand{\tref}[2][]{#2}
\newcommand{\termdef}[2]{\hypertarget{#1}{}\emph{#2}\index{#2}}
\newcommand{\termref}[2]{\hyperlink{#1}{#2}}
\newcommand{\exttermref}[3]{\hyperref{#1}{}{#2}{#3}}
\hypersetup{pdftitle={Receiver POVM (receiver) -- statements digest 0.11.0}}

% The digest: the contract's labels under the prefix C-, its numbers checked.
\externaldocument[C-]{contract}
\theoremstyle{claim}\newtheorem*{digestclaim}{\RestatedName}
\theoremstyle{claimass}\newtheorem*{digestassumption}{\RestatedName}
\theoremstyle{claimreq}\newtheorem*{digestrequirement}{\RestatedName}
\makeatletter
% #1 label, #2 the number the assembler computed: a local label carrying it, so
% that a reference between two digest items stays in the digest; the build
% stops when the contract's own number differs.
\newcommand{\digestlabel}[2]{\phantomsection
  \edef\digest@theirs{\getrefnumber{C-#1}}\def\digest@ours{#2}%
  \ifx\digest@theirs\digest@ours\else
    \PackageError{digest}{#1 is #2 here but \digest@theirs\space in the contract}{}\fi
  \def\@currentlabel{#2}\label{#1}}
\makeatother
% A display's local label (amsmath records the \tag, i.e. the contract's number).
\newcommand{\digesteqlabel}[1]{\label{#1}}

% Claims of other contracts cited here (\xref, [[...]]).
\DeclareXref{foundations/ass:snapped-measurement}{Assumption}{The device measures the snapped family where the run reads the snap}{../foundations/contract.pdf}
\DeclareXref{foundations/cor:least-eigenvalue-over-enclosure}{Corollary}{Least eigenvalue and rank over an enclosure}{../foundations/contract.pdf}
\DeclareXref{foundations/cor:rank-over-factor-enclosure}{Corollary}{Rank over a factor enclosure}{../foundations/contract.pdf}
\DeclareXref{foundations/def:binary64}{Definition}{Binary64 numbers, declared values and dyadic lattices}{../foundations/contract.pdf}
\DeclareXref{foundations/def:declared-data}{Definition}{Declared data and the parties to a module boundary}{../foundations/contract.pdf}
\DeclareXref{foundations/def:evaluated-off-snap}{Definition}{A reference handed beside a snap; a block evaluated off its snap}{../foundations/contract.pdf}
\DeclareXref{foundations/def:exact-island}{Definition}{The exact island}{../foundations/contract.pdf}
\DeclareXref{foundations/def:exact-reference}{Definition}{Reference of a measurement family, exact or enclosed}{../foundations/contract.pdf}
\DeclareXref{foundations/def:factor-enclosure}{Definition}{Factor enclosure of a block}{../foundations/contract.pdf}
\DeclareXref{foundations/def:round-homogeneity}{Definition}{Round homogeneity of a per-round family}{../foundations/contract.pdf}
\DeclareXref{foundations/def:sector-map}{Definition}{Sector maps and block-list operators}{../foundations/contract.pdf}
\DeclareXref{foundations/def:snap}{Definition}{Snapping declared data onto a lattice}{../foundations/contract.pdf}
\DeclareXref{foundations/def:tiers}{Definition}{The tiers of a reported key rate}{../foundations/contract.pdf}
\DeclareXref{foundations/lem:enclosure-psd}{Lemma}{Positive semidefiniteness from an entrywise enclosure}{../foundations/contract.pdf}
\DeclareXref{foundations/lem:least-eigenvalue-bound}{Lemma}{Least-eigenvalue lower bound from an approximate eigenbasis}{../foundations/contract.pdf}
\DeclareXref{foundations/lem:rational-obstruction}{Lemma}{Rationality obstruction to the domination of a lattice family}{../foundations/contract.pdf}
\DeclareXref{foundations/rem:noise-floor}{Remark}{What domination controls, and why it needs a noise floor}{../foundations/contract.pdf}
\DeclareXref{foundations/req:consumer-decides-exactly}{Requirement}{Exact decision by every consumer}{../foundations/contract.pdf}
\DeclareXref{foundations/req:rounding-direction}{Requirement}{Rounding direction of emitted bounds}{../foundations/contract.pdf}
\DeclareXref{foundations/req:snap-reference}{Requirement}{Snaps of a measurement family: kept, or moved and disclosed}{../foundations/contract.pdf}
\DeclareXref{foundations/req:support-safe-snap}{Requirement}{Support-safe snaps of a measurement family}{../foundations/contract.pdf}
\DeclareXref{foundations/req:tiers}{Requirement}{Reporting of the tiers}{../foundations/contract.pdf}
\DeclareXref{foundations/req:tolerance-scope}{Requirement}{Scope of the declared tolerance}{../foundations/contract.pdf}
\DeclareXref{foundations/req:who-snaps}{Requirement}{Snapping by the declarer only}{../foundations/contract.pdf}
\DeclareXref{meat-core/frank-wolfe/rem:fo-arithmetic}{Remark}{Exact real arithmetic in the first-order bound}{../meat-core/contract.pdf}
\DeclareXref{meat-core/frank-wolfe/req:bracket-on-reference}{Requirement}{The first-order bound evaluated on the reference companion}{../meat-core/contract.pdf}
\DeclareXref{meat-core/frank-wolfe/req:certified-evaluation}{Requirement}{The certified evaluation of the first-order bound}{../meat-core/contract.pdf}
\DeclareXref{reductions/def:joint-family}{Definition}{Joint alphabet, classes and the joint family of a private setting}{../reductions/contract.pdf}
\DeclareXref{reductions/lem:row-relaxation}{Lemma}{A rounded constraint entry with its constant relaxed}{../reductions/contract.pdf}

\begin{document}
\begin{center}
{\LARGE\bfseries Receiver POVM (receiver)}\\[4pt]
{\large Statements digest, version 0.11.0}
\end{center}
\noindent{\small Every definition, fact, assumption, result and requirement of the \href{contract.pdf}{contract}, with the contract's numbers; remarks, examples, technical steps and all proofs are left out. Each heading opens its statement in the contract.}

\part*{Mathematics}\pdfbookmark[0]{Mathematics}{digest-part-1}

\section*{1\quad Modes, Fock space and the mode algebra}\pdfbookmark[1]{1 Modes, Fock space and the mode algebra}{digest-group-receiver/def:modes}

% claim receiver/def:modes
\digestlabel{def:modes}{1.1}\pdfbookmark[2]{Definition 1.1 Declared mode space}{digest-receiver/def:modes}
\renewcommand*{\RestatedName}{\hyperref{contract.pdf}{}{claim-receiver/def:modes}{Definition~1.1}}%
\begin{digestclaim}[Declared mode space]
A \emph{declared mode space} is a finite totally ordered set $\Modes$, each of
whose elements is an \emph{optical mode}, together with the Hilbert space
$\mathbb{C}^{\Modes}$ carrying the orthonormal basis $\{\ket{m}\}_{m\in\Modes}$
indexed in the declared order. A mode is a label $m=(p,\pi,\tau)$ recording a
spatial path $p$, a polarization $\pi\in\{H,V\}$ and a temporal/spectral index
$\tau$; the order is lexicographic in $(p,\tau,\pi)$ with $H$ before $V$. The
index $\tau$ is \emph{reserved}: it is carried in the label and in the emitted
artifact, it may take only a single declared value in this contract, and no
catalogue entry (Definition~\ref{def:catalogue-entries}) reads or writes it.
For $\Sigma\subseteq\Modes$ we write
$\Emd{\Sigma}=\sum_{m\in\Sigma}\dyad{m}$ for the coordinate projector, and
$\Modes$ is identified with $\{0,\dots,|\Modes|-1\}$ when matrix indices are
needed.
\end{digestclaim}

% claim receiver/def:fock
\digestlabel{def:fock}{1.2}\pdfbookmark[2]{Definition 1.2 Fock space and the second-quantization functor}{digest-receiver/def:fock}
\renewcommand*{\RestatedName}{\hyperref{contract.pdf}{}{claim-receiver/def:fock}{Definition~1.2}}%
\begin{digestclaim}[Fock space and the second-quantization functor]
For a finite-dimensional Hilbert space $\mathcal{K}_0$ the \emph{Fock space} is
$\mathcal{F}(\mathcal{K}_0)=\bigoplus_{n\ge0}\Symn{n}(\mathcal{K}_0)$, where
$\Symn{n}(\mathcal{K}_0)\subseteq\mathcal{K}_0^{\otimes n}$ is the subspace
invariant under every permutation of the $n$ tensor factors and
$\Symn{0}(\mathcal{K}_0)=\mathbb{C}\ket{\vac}$. For a linear operator $X$ on
$\mathcal{K}_0$ put
\begin{equation*}\tag{\getrefnumber{C-eq:fockfunctor}}\digesteqlabel{eq:fockfunctor}
  \Fock(X)\coloneqq\bigoplus_{n\ge0}\big(X^{\otimes n}\big)\big|_{\Symn{n}},
  \qquad \Fock(X)\big|_{\Symn{0}}=1 .
\end{equation*}
This is well defined because $X^{\otimes n}$ commutes with every permutation of
the factors and therefore preserves $\Symn{n}(\mathcal{K}_0)$
(Lemma~\ref{lem:fock-functor}). The same formula defines
$\Fock(X):\mathcal{F}(\mathcal{K}_0)\to\mathcal{F}(\mathcal{K}_1)$ for a linear
map $X:\mathcal{K}_0\to\mathcal{K}_1$ between two spaces, because
$X^{\otimes n}$ intertwines the permutation actions and therefore carries
$\Symn{n}(\mathcal{K}_0)$ into $\Symn{n}(\mathcal{K}_1)$. Occupation vectors
$\vec m\in\mathbb{N}_0^{\dim\mathcal{K}_0}$ with $|\vec m|=n$ index the
orthonormal basis $\{\ket{\vec m}\}$ of $\Symn{n}(\mathcal{K}_0)$ obtained by
symmetrizing and normalizing $\ket{e_1}^{\otimes m_1}\otimes\cdots$ in a fixed
orthonormal basis $\{\ket{e_j}\}$ of $\mathcal{K}_0$; $a_j$ and $a_j^\dagger$
denote the annihilation and creation operators of $\ket{e_j}$.
\end{digestclaim}

% claim receiver/def:embed
\digestlabel{def:embed}{1.3}\pdfbookmark[2]{Definition 1.3 Local element and its embedding}{digest-receiver/def:embed}
\renewcommand*{\RestatedName}{\hyperref{contract.pdf}{}{claim-receiver/def:embed}{Definition~1.3}}%
\begin{digestclaim}[Local element and its embedding]
A \emph{local element} on a mode subset $\Sigma\subseteq\Modes$ is a matrix
$\Gmat_\Sigma\in\mathbb{C}^{\Sigma\times\Sigma}$ with
$\Gmat_\Sigma^\dagger\Gmat_\Sigma\preceq\id_\Sigma$; $\Sigma$ is its
\emph{support}. Its \emph{global embedding} is
\begin{equation*}\tag{\getrefnumber{C-eq:embed}}\digesteqlabel{eq:embed}
  \hat\Gmat\coloneqq\Gmat_\Sigma\oplus\id_{\Modes\setminus\Sigma}
  \ \in\ \mathbb{C}^{\Modes\times\Modes},
\end{equation*}
that is, $\hat\Gmat\,\Emd{\Sigma}=\Gmat_\Sigma\Emd{\Sigma}$ and
$\hat\Gmat\,(\id-\Emd{\Sigma})=\id-\Emd{\Sigma}$. A local element with
$\Gmat_\Sigma^\dagger\Gmat_\Sigma=\id_\Sigma$ is a \emph{lossless element}.
\end{digestclaim}

% claim receiver/prop:mode-algebra
\digestlabel{prop:mode-algebra}{1.4}\pdfbookmark[2]{Proposition 1.4 Algebra of embedded local elements}{digest-receiver/prop:mode-algebra}
\renewcommand*{\RestatedName}{\hyperref{contract.pdf}{}{claim-receiver/prop:mode-algebra}{Proposition~1.4}}%
\begin{digestclaim}[Algebra of embedded local elements]
Let $\hat\Gmat,\hat\Gmat'$ be global embeddings of
local elements with supports $\Sigma,\Sigma'$.
\begin{enumerate}[label=(\roman*),leftmargin=3.0em]
\item \emph{Embedding preserves contractivity:}
  $\hat\Gmat^\dagger\hat\Gmat\preceq\id_{\Modes}$, with equality exactly when
  the element is lossless.
\item \emph{Products of contractions are contractions:} if
  $X^\dagger X\preceq\id$ and $Y^\dagger Y\preceq\id$ then
  $(XY)^\dagger(XY)\preceq\id$, with equality for all vectors exactly when
  $Y$ is an isometry and $X$ is isometric on $\operatorname{ran}Y$.
\item \emph{Disjoint supports commute:} if $\Sigma\cap\Sigma'=\varnothing$ then
  $\hat\Gmat\hat\Gmat'=\hat\Gmat'\hat\Gmat$.
\item \emph{Adjoints and unitaries:} $\hat\Gmat^\dagger$ is the embedding of
  $\Gmat_\Sigma^\dagger$, and $\hat\Gmat$ is unitary exactly when
  $\Gmat_\Sigma$ is.
\end{enumerate}
\end{digestclaim}

% claim receiver/lem:fock-functor
\digestlabel{lem:fock-functor}{1.6}\pdfbookmark[2]{Lemma 1.6 Functoriality of the Fock lift}{digest-receiver/lem:fock-functor}
\renewcommand*{\RestatedName}{\hyperref{contract.pdf}{}{claim-receiver/lem:fock-functor}{Lemma~1.6}}%
\begin{digestclaim}[Functoriality of the Fock lift]
Let $\mathcal{K}_0$ be finite-dimensional and let $X,Y$ be operators on
$\mathcal{K}_0$. Then, with $\Fock$ as in \eqref{eq:fockfunctor}:
\begin{enumerate}[label=(\roman*),leftmargin=3.0em]
\item $X^{\otimes n}$ preserves $\Symn{n}(\mathcal{K}_0)$, so $\Fock(X)$ is
  well defined, and it is block diagonal with respect to the photon-number
  decomposition: $\Fock(X)\Pnum{n}=\Pnum{n}\Fock(X)$ for every $n$;
\item $\Fock(XY)=\Fock(X)\Fock(Y)$, $\Fock(\id)=\id$, and
  $\Fock(X)^\dagger=\Fock(X^\dagger)$;
\item if $U$ is unitary then $\Fock(U)$ is unitary; if
  $X^\dagger X\preceq\id$ then $\lVert\Fock(X)\rVert\le1$, so $\Fock(X)$ is a
  bounded operator on $\mathcal{F}(\mathcal{K}_0)$;
\item if $P$ is an orthogonal projector then $\Fock(P)$ is the orthogonal
  projector onto $\mathcal{F}(\operatorname{ran}P)\subseteq
  \mathcal{F}(\mathcal{K}_0)$ --- equivalently, with the occupation vectors
  formed in an orthonormal basis of $\mathcal{K}_0$ adapted to
  $\operatorname{ran}P\oplus\ker P$, onto the closed span of those carrying
  no photon outside $\operatorname{ran}P$;
\item if $U$ is unitary then $\Fock(U)$ implements the mode transformation on
  the annihilation operators,
  $\Fock(U)^\dagger a_i\Fock(U)=\sum_jU_{ij}a_j$.
\end{enumerate}
\end{digestclaim}

% claim receiver/lem:fock-block-spectrum
\digestlabel{lem:fock-block-spectrum}{1.9}\pdfbookmark[2]{Lemma 1.9 Spectrum of a lifted no-click operator on a photon-number block}{digest-receiver/lem:fock-block-spectrum}
\renewcommand*{\RestatedName}{\hyperref{contract.pdf}{}{claim-receiver/lem:fock-block-spectrum}{Lemma~1.9}}%
\begin{digestclaim}[Spectrum of a lifted no-click operator on a photon-number block]
Let $A$ be Hermitian on $\mathcal{K}_0$ with $0\preceq A\preceq\id$ and
eigenvalues $\lambda_1,\dots,\lambda_d$ ($d=\dim\mathcal{K}_0$). Then
$\Fock(\id-A)$ is a positive semidefinite contraction, its $n$-photon block
$\Fock(\id-A)^{(n)}=\big((\id-A)^{\otimes n}\big)\big|_{\Symn{n}}$ has spectrum
\begin{equation*}\tag{\getrefnumber{C-eq:spectrum}}\digesteqlabel{eq:spectrum}
  \Big\{\textstyle\prod_{j=1}^{d}(1-\lambda_j)^{m_j}\ :\
  \vec m\in\mathbb{N}_0^{d},\ |\vec m|=n\Big\},
\end{equation*}
and consequently
\begin{equation*}\tag{\getrefnumber{C-eq:blockbounds}}\digesteqlabel{eq:blockbounds}
  \big(1-\lmax(A)\big)^{n}\,\id\ \preceq\ \Fock(\id-A)^{(n)}\ \preceq\
  \big(1-\lmin(A)\big)^{n}\,\id ,
\end{equation*}
both bounds being attained as eigenvalues.
\end{digestclaim}

% claim receiver/lem:fock-monotone
\digestlabel{lem:fock-monotone}{1.10}\pdfbookmark[2]{Lemma 1.10 Order monotonicity of the Fock lift}{digest-receiver/lem:fock-monotone}
\renewcommand*{\RestatedName}{\hyperref{contract.pdf}{}{claim-receiver/lem:fock-monotone}{Lemma~1.10}}%
\begin{digestclaim}[Order monotonicity of the Fock lift]
If $X\succeq Y\succeq0$ on $\mathcal{K}_0$ then $\Fock(X)\succeq\Fock(Y)$.
Consequently, if $0\preceq A\preceq A'\preceq\id$ then
$\Fock(\id-A)\succeq\Fock(\id-A')$.
\end{digestclaim}

% claim receiver/prop:normord
\digestlabel{prop:normord}{1.11}\pdfbookmark[2]{Proposition 1.11 Normal-ordered form of a lifted no-click operator}{digest-receiver/prop:normord}
\renewcommand*{\RestatedName}{\hyperref{contract.pdf}{}{claim-receiver/prop:normord}{Proposition~1.11}}%
\begin{digestclaim}[Normal-ordered form of a lifted no-click operator]
For $0\preceq A\preceq\id$ on $\mathcal{K}_0$,
\begin{equation*}\tag{\getrefnumber{C-eq:normordform}}\digesteqlabel{eq:normordform}
  \Fock(\id-A)=\normord{\exp\!\big(-\vec a^\dagger A\vec a\big)}
  \coloneqq\sum_{p\ge0}\frac{(-1)^p}{p!}
  \sum_{i_1\dots i_p}\sum_{j_1\dots j_p}
  \Big(\prod_{r=1}^{p}A_{i_rj_r}\Big)
  a_{i_1}^\dagger\cdots a_{i_p}^\dagger a_{j_1}\cdots a_{j_p},
\end{equation*}
the series acting on each photon-number block as a finite sum.
\end{digestclaim}

\section*{2\quad The receiver instrument and its physical model}\pdfbookmark[1]{2 The receiver instrument and its physical model}{digest-group-receiver/def:instrument}

% claim receiver/def:instrument
\digestlabel{def:instrument}{2.1}\pdfbookmark[2]{Definition 2.1 Receiver instrument}{digest-receiver/def:instrument}
\renewcommand*{\RestatedName}{\hyperref{contract.pdf}{}{claim-receiver/def:instrument}{Definition~2.1}}%
\begin{digestclaim}[Receiver instrument]
A \emph{receiver instrument} is a tuple
\begin{equation*}\tag{\getrefnumber{C-eq:instrument}}\digesteqlabel{eq:instrument}
  \Recv=\Big(\Modes;\ \Inp=\InpC\sqcup\InpS;\
  \{\Port{k}\}_{k=1}^{K};\ \Gmat\Big)
\end{equation*}
in which
\begin{enumerate}[label=(\roman*),leftmargin=3.0em]
\item $\Modes$ is a declared mode space (Definition~\ref{def:modes});
\item $\Inp\subseteq\Modes$ is the set of \emph{declared input modes}, split
  into the \emph{channel modes} $\InpC\neq\varnothing$ (the adversary's access) and the
  \emph{declared source modes} $\InpS$ (the probe); every mode of
  $\Modes\setminus\Inp$ carries vacuum at the input
  (Assumption~\ref{ass:vac});
\item $\Port{1},\dots,\Port{K}$ are pairwise disjoint subsets of $\Modes$, the
  \emph{detector ports}, with $K\ge1$; a mode in no port is undetected;
\item $\Gmat\in\mathbb{C}^{\Modes\times\Modes}$ is the \emph{instrument matrix},
  indexed rows-by-outputs and columns-by-inputs, with
  $\Gmat^\dagger\Gmat\preceq\id$.
\end{enumerate}
The basis $\{\ket{m}\}_{m\in\Inp}$ of $\mathbb{C}^{\Inp}$ is ordered with the
channel modes $\InpC$ before the declared source modes $\InpS$, each part in
the declared order of $\Modes$; matrices on $\mathbb{C}^{\Inp}$ and the
columns of a matrix with column set $\Inp$ are indexed in this order, which
differs from the declared order of $\Modes$ exactly when some declared source
mode precedes some channel mode in it.
We write $\Gmat_{\Inp}\in\mathbb{C}^{\Modes\times\Inp}$ for $\Gmat$ restricted
to the columns $\Inp$, and $\Emd{J}\coloneqq\Emd{\bigcup_{k\in J}\Port{k}}$ for
$J\subseteq[K]$, where $[K]\coloneqq\{1,\dots,K\}$.
The channel-input system of the shared notation is the Fock space over
the channel modes: $\Hs{\Chan}=\bigoplus_n\Hs{\Chan_n}$ with
$\Hs{\Chan_n}=\Symn{n}(\mathbb{C}^{\InpC})$ and photon-number projectors
$\Pnumon{n}{\Chan}$.
\end{digestclaim}

% claim receiver/def:gram
\digestlabel{def:gram}{2.2}\pdfbookmark[2]{Definition 2.2 Port Grams and the loss deficiency}{digest-receiver/def:gram}
\renewcommand*{\RestatedName}{\hyperref{contract.pdf}{}{claim-receiver/def:gram}{Definition~2.2}}%
\begin{digestclaim}[Port Grams and the loss deficiency]
The \emph{port Gram matrices} of $\Recv$ are the operators on
$\mathbb{C}^{\Inp}$
\begin{equation*}\tag{\getrefnumber{C-eq:gramdef}}\digesteqlabel{eq:gramdef}
  A_k\coloneqq\Gmat_{\Inp}^\dagger\,\Emd{\{k\}}\,\Gmat_{\Inp},
  \qquad k=1,\dots,K ,
\end{equation*}
and for $J\subseteq[K]$ we write $A_J\coloneqq\sum_{k\in J}A_k$ and
$A_\Sigma\coloneqq A_{[K]}$. The \emph{loss deficiency} is
\begin{equation*}\tag{\getrefnumber{C-eq:deficiency}}\digesteqlabel{eq:deficiency}
  \defic\coloneqq\id_{\Inp}-A_\Sigma .
\end{equation*}
A family $\{A_k\}_{k=1}^{K}$ of operators on a finite-dimensional space is an
\emph{admissible Gram family} when $A_k\succeq0$ for every $k$ and
$\sum_kA_k\preceq\id$. This is an exact condition; on a \emph{declared} family
it is \emph{decided} in floating point at the
receiver tolerance
$\varepsilon_2$, by Definition~\ref{def:tolmap}\,(i).
\end{digestclaim}

% claim receiver/def:probe
\digestlabel{def:probe}{2.3}\pdfbookmark[2]{Definition 2.3 Probe declaration}{digest-receiver/def:probe}
\renewcommand*{\RestatedName}{\hyperref{contract.pdf}{}{claim-receiver/def:probe}{Definition~2.3}}%
\begin{digestclaim}[Probe declaration]
A \emph{probe declaration} is a pair $\big(\Tset,\{\rho^{(t)}\}_{t\in\Tset}\big)$
in which $\Tset$ is a finite non-empty alphabet of probe settings (including
the setting $\off$, whose state is vacuum) and, for each $t$, $\rho^{(t)}$ is a
declared state on $\mathcal{F}(\mathbb{C}^{\InpS})$ that is
\emph{occupation-diagonal}: diagonal in the occupation basis of
$\mathcal{F}(\mathbb{C}^{\InpS})$ (diagonal in the Fock basis). In the
\emph{Gaussian tier} each $\rho^{(t)}$ is a zero-mean Gaussian state with
diagonal normally-ordered second-moment matrix
$\Gamma^{(t)}=\operatorname{diag}(\bar n^{(t)}_m)_{m\in\InpS}\succeq0$, that is,
a product of single-mode thermal states with mean photon numbers
$\bar n^{(t)}_m\ge0$; $\Gamma^{(t)}=0$ is vacuum. A state of the Gaussian tier
is occupation-diagonal by construction, a thermal state being diagonal in the
occupation basis. When $\InpS=\varnothing$ the
declaration is $\Tset=\{\off\}$ and every statement about conditioning is
vacuous.
\end{digestclaim}

% claim receiver/lem:dilation
\digestlabel{lem:dilation}{2.4}\pdfbookmark[2]{Lemma 2.4 Unitary dilation of a contraction}{digest-receiver/lem:dilation}
\renewcommand*{\RestatedName}{\hyperref{contract.pdf}{}{claim-receiver/lem:dilation}{Lemma~2.4}}%
\begin{digestclaim}[Unitary dilation of a contraction]
Let $\Gmat$ be a contraction on $\mathbb{C}^{\Modes}$ and let
$\Env$ be a disjoint copy of $\Modes$. With the defect operators
$D_{\mathrm{out}}\coloneqq(\id-\Gmat\Gmat^\dagger)^{1/2}$ and
$D_{\mathrm{in}}\coloneqq(\id-\Gmat^\dagger\Gmat)^{1/2}$, the matrix
\begin{equation*}\tag{\getrefnumber{C-eq:dilationmatrix}}\digesteqlabel{eq:dilationmatrix}
  U_{\Gmat}=\begin{pmatrix}\Gmat & D_{\mathrm{out}}\\
                           D_{\mathrm{in}} & -\Gmat^\dagger\end{pmatrix}
  \quad\text{on }\ \mathbb{C}^{\Modes}\oplus\mathbb{C}^{\Env}
\end{equation*}
is unitary, and its block from the inputs $\Modes$ to the outputs $\Modes$ is
$\Gmat$.
\end{digestclaim}

% claim receiver/def:physmodel
\digestlabel{def:physmodel}{2.5}\pdfbookmark[2]{Definition 2.5 The physical measurement model}{digest-receiver/def:physmodel}
\renewcommand*{\RestatedName}{\hyperref{contract.pdf}{}{claim-receiver/def:physmodel}{Definition~2.5}}%
\begin{digestclaim}[The physical measurement model]
Fix a receiver instrument $\Recv$ and a unitary extension $U$ of $\Gmat$
on $\mathbb{C}^{\Modes\sqcup\Env}$ (Lemma~\ref{lem:dilation} supplies one; any
other may be used, and by Theorem~\ref{thm:port-grams}\,(v) the choice does not
matter). Let
$\iota:\mathbb{C}^{\Inp}\hookrightarrow\mathbb{C}^{\Modes\sqcup\Env}$ be the
inclusion of the declared input modes and
$\Fock(\iota):\mathcal{F}(\mathbb{C}^{\Inp})\to
\mathcal{F}(\mathbb{C}^{\Modes\sqcup\Env})$ the induced map, which under the
factorization
$\mathcal{F}(\mathbb{C}^{\Modes\sqcup\Env})\cong
 \mathcal{F}(\mathbb{C}^{\Inp})\otimes
 \mathcal{F}(\mathbb{C}^{(\Modes\sqcup\Env)\setminus\Inp})$, the occupation
identification of Lemma~\ref{C-lem:fock-maps}\,(ii),
is $\ket{\psi}\mapsto\ket{\psi}\otimes\ket{\vac}$: it prepares vacuum on every
undeclared input mode and on every dilation mode
(Assumption~\ref{ass:vac}). The \emph{measurement isometry} of the receiver is
\begin{equation*}\tag{\getrefnumber{C-eq:measisometry}}\digesteqlabel{eq:measisometry}
  V\coloneqq\Fock(U)\,\Fock(\iota)
  \ :\ \mathcal{F}(\mathbb{C}^{\Inp})\longrightarrow
        \mathcal{F}(\mathbb{C}^{\Modes\sqcup\Env}) .
\end{equation*}
The detectors measure, on the output side, the $K$ commuting binary projective
measurements $\{\Pi^{0}_{k},\id-\Pi^{0}_{k}\}$ where
$\Pi^0_k\coloneqq\Fock(\id-\Emd{\{k\}})$, with the projector $\Emd{\{k\}}$ of
$\mathbb{C}^{\Modes}$ extended by zero to $\mathbb{C}^{\Modes\sqcup\Env}$, is
the projector onto ``no photon in the modes of $\Port{k}$''
(Lemma~\ref{lem:fock-functor}\,(iv), Assumption~\ref{ass:port}). For $J\subseteq[K]$ put
$\Pi^0_J\coloneqq\prod_{k\in J}\Pi^0_k$.
\end{digestclaim}

% claim receiver/lem:vacelim
\digestlabel{lem:vacelim}{2.6}\pdfbookmark[2]{Lemma 2.6 Exactness of vacuum elimination}{digest-receiver/lem:vacelim}
\renewcommand*{\RestatedName}{\hyperref{contract.pdf}{}{claim-receiver/lem:vacelim}{Lemma~2.6}}%
\begin{digestclaim}[Exactness of vacuum elimination]
Let $S\subseteq T$ be finite sets of modes and
$\iota:\mathbb{C}^{S}\hookrightarrow\mathbb{C}^{T}$ the coordinate inclusion.
Under the factorization
$\mathcal{F}(\mathbb{C}^{T})\cong\mathcal{F}(\mathbb{C}^{S})\otimes
\mathcal{F}(\mathbb{C}^{T\setminus S})$ of occupation bases, the occupation
identification of Lemma~\ref{C-lem:fock-maps}\,(ii), $\Fock(\iota)$ is
$\ket{\psi}\mapsto\ket{\psi}\otimes\ket{\vac}$; and for any operator $X$ on
$\mathbb{C}^{T}$,
\begin{equation*}\tag{\getrefnumber{C-eq:vacelim}}\digesteqlabel{eq:vacelim}
  \Fock(\iota)^\dagger\,\Fock(X)\,\Fock(\iota)
  \;=\;\Fock\big(\iota^\dagger X\iota\big),
\end{equation*}
where $\iota^\dagger X\iota$ is the $S\times S$ block of $X$. In
particular $\Fock(\iota)$ is an isometry. The inclusion $\iota$ of
Definition~\ref{def:physmodel}, with $S=\Inp$ and $T=\Modes\sqcup\Env$, is
one instance, and for it $V$ of \eqref{eq:measisometry} is an isometry as
well.
\end{digestclaim}

\section*{3\quad Component graphs and compilation}\pdfbookmark[1]{3 Component graphs and compilation}{digest-group-receiver/def:catalogue}

% claim receiver/def:catalogue
\digestlabel{def:catalogue}{3.1}\pdfbookmark[2]{Definition 3.1 Component catalogue}{digest-receiver/def:catalogue}
\renewcommand*{\RestatedName}{\hyperref{contract.pdf}{}{claim-receiver/def:catalogue}{Definition~3.1}}%
\begin{digestclaim}[Component catalogue]
A \emph{component catalogue} $\Cat$ is a finite set of entry types. An
\emph{entry type} fixes a support shape (which modes of $\Modes$ it acts on), a
parameter domain, and a rule assigning to every parameter tuple in that domain
a local element in the sense of Definition~\ref{def:embed}. The
catalogue of this contract is the one of Definition~\ref{def:catalogue-entries}
(Table~\ref{C-tab:catalogue}) in the component-catalogue submodule, and
Theorem~\ref{thm:catalogue-admissibility} there states that every entry type maps its whole
declared domain into local elements.
\end{digestclaim}

% claim receiver/def:graph
\digestlabel{def:graph}{3.2}\pdfbookmark[2]{Definition 3.2 Declared element graph}{digest-receiver/def:graph}
\renewcommand*{\RestatedName}{\hyperref{contract.pdf}{}{claim-receiver/def:graph}{Definition~3.2}}%
\begin{digestclaim}[Declared element graph]
A \emph{declared element graph} is a finite directed graph $\Grph=(V,E)$ whose
vertices $e\in V$ are \emph{elements}, each carrying a catalogue entry type, a
parameter tuple, and a support $\Sigma_e\subseteq\Modes$. Write $\preceq_\Grph$
for the reachability relation: $e\preceq_\Grph f$ when $e=f$ or there is a
directed path from $e$ to $f$. The graph is \emph{well formed} when
\begin{enumerate}[label=(E\arabic*),leftmargin=3.0em]
\item every element's parameter tuple lies in the declared domain of its entry
  type, so that, by Definition~\ref{def:catalogue}, $\Gmat_{\Sigma_e}$ is a
  local element: defined, with
  $\Gmat_{\Sigma_e}^\dagger\Gmat_{\Sigma_e}\preceq\id$;
\item $\Grph$ is acyclic, so that $\preceq_\Grph$ is a partial order;
\item (\emph{support comparability}) any two elements whose supports intersect
  are comparable: $\Sigma_e\cap\Sigma_f\neq\varnothing\implies
  e\preceq_\Grph f$ or $f\preceq_\Grph e$.
\end{enumerate}
A \emph{linear extension} of $\Grph$ is an enumeration
$e_1,\dots,e_L$ of $V$ with $e_i\preceq_\Grph e_j\implies i\le j$.
\end{digestclaim}

% claim receiver/def:compile
\digestlabel{def:compile}{3.3}\pdfbookmark[2]{Definition 3.3 Compilation of a declared element graph}{digest-receiver/def:compile}
\renewcommand*{\RestatedName}{\hyperref{contract.pdf}{}{claim-receiver/def:compile}{Definition~3.3}}%
\begin{digestclaim}[Compilation of a declared element graph]
Let $\Grph$ be a well-formed graph; it is in particular acyclic
(condition (E2) of Definition~\ref{def:graph}), so that a linear extension exists. For a
linear extension $\sigma=(e_1,\dots,e_L)$ of $\Grph$ put
\begin{equation*}\tag{\getrefnumber{C-eq:compile}}\digesteqlabel{eq:compile}
  \Gmat_{\Grph}^{\sigma}\coloneqq
  \hat\Gmat_{e_L}\,\hat\Gmat_{e_{L-1}}\cdots\hat\Gmat_{e_1}
  \ \in\ \mathbb{C}^{\Modes\times\Modes},
\end{equation*}
later elements multiplying from the left. The
\emph{compiled instrument matrix} of the receiver is
$\Gmat\coloneqq D_\eta\,\Gmat^{\sigma}_{\Grph}$, where
$D_\eta=\sum_k\sqrt{\eta_k}\,\Emd{\{k\}}+\Emd{\Modes\setminus\bigcup_k
\Port{k}}$ is the \emph{terminal detector-efficiency layer}, applied after every
element, and $\eta_k\in[0,1]$ is the declared efficiency of port $k$ (the
effect of that layer on the port Grams is Theorem~\ref{thm:efficiency} of the
detector-imperfections submodule). The detector ports $\Port{k}$ and the
port-embedding abbreviation $\Emd{\{k\}}$ written here are those of a
receiver instrument tuple: clause (iii) of
Definition~\ref{def:instrument} and its closing notation.
\end{digestclaim}

% claim receiver/thm:compilation-well-defined
\digestlabel{thm:compilation-well-defined}{3.4}\pdfbookmark[2]{Theorem 3.4 Well-definedness of compilation}{digest-receiver/thm:compilation-well-defined}
\renewcommand*{\RestatedName}{\hyperref{contract.pdf}{}{claim-receiver/thm:compilation-well-defined}{Theorem~3.4}}%
\begin{digestclaim}[Well-definedness of compilation]
In items (i)--(iii), let $\Grph$ be a well-formed declared element graph
(Definition~\ref{def:graph}); in particular $\Grph$ is acyclic
(condition (E2) of Definition~\ref{def:graph}), so that $\preceq_\Grph$ is a partial order and
linear extensions of $\Grph$ and of each of its induced subgraphs exist.
Then:
\begin{enumerate}[label=(\roman*),leftmargin=3.0em]
\item \emph{Independence of the sort.} $\Gmat^{\sigma}_{\Grph}$ is the same
  matrix for every linear extension $\sigma$ of $\Grph$; write
  $\Gmat_{\Grph}$ for it. In particular, if $V$ is partitioned into
  \emph{layers} $L_1,\dots,L_r$ such that elements of one layer have pairwise
  disjoint supports and $e\preceq_\Grph f$ implies
  $\operatorname{layer}(e)\le\operatorname{layer}(f)$, then
  $\Gmat_\Grph=\big(\prod_{e\in L_r}\hat\Gmat_e\big)\cdots
  \big(\prod_{e\in L_1}\hat\Gmat_e\big)$ with each intra-layer product
  order-independent.
\item \emph{Admissibility by construction.}
  $\Gmat_\Grph^\dagger\Gmat_\Grph\preceq\id$, and hence also
  $\Gmat^\dagger\Gmat\preceq\id$ for the compiled instrument matrix of
  Definition~\ref{def:compile}.
\item \emph{Losslessness.} $\Gmat_\Grph$ is unitary if and only if every
  element of $\Grph$ is lossless on its support. If in addition every
  $\eta_k=1$ then $\Gmat$ is unitary.
\item \emph{Decidability.} For an arbitrary declared element graph, not
  necessarily well formed, conditions (E1)--(E3) are decidable by finite
  inspection: (E1) by evaluating the declared domain predicate of each
  element, (E2) by a topological sort, (E3) by testing
  $\Sigma_e\cap\Sigma_f\neq\varnothing$ against the transitive closure of $E$
  for each of the $\binom{|V|}{2}$ pairs.
\end{enumerate}
\end{digestclaim}

% claim receiver/cor:compiled-instrument
\digestlabel{cor:compiled-instrument}{3.5}\pdfbookmark[2]{Corollary 3.5 Receiver instrument from a well-formed graph declaration}{digest-receiver/cor:compiled-instrument}
\renewcommand*{\RestatedName}{\hyperref{contract.pdf}{}{claim-receiver/cor:compiled-instrument}{Corollary~3.5}}%
\begin{digestclaim}[Receiver instrument from a well-formed graph declaration]
Let a route-(iii) declaration supply a declared mode space $\Modes$; a
well-formed declared element graph $\Grph$ on $\Modes$
(Definition~\ref{def:graph}), in particular acyclic
(condition (E2) of Definition~\ref{def:graph}); a set
$\Inp\subseteq\Modes$ of declared input modes with a split
$\Inp=\InpC\sqcup\InpS$, $\InpC\neq\varnothing$, every mode of $\Modes\setminus\Inp$ carrying vacuum
at the input (Assumption~\ref{ass:vac}); and a port map
$\Port{1},\dots,\Port{K}$, pairwise disjoint subsets of $\Modes$ with
$K\ge1$, carrying detector efficiencies $\eta_k$ in the domain $[0,1]$
declared for them in Definition~\ref{def:compile}. Let $\sigma$ be a linear
extension of $\Grph$ and let
$\Gmat\coloneqq D_\eta\,\Gmat^{\sigma}_{\Grph}$ be the compiled instrument
matrix of Definition~\ref{def:compile}. Then
\begin{enumerate}[label=(\roman*),leftmargin=3.0em]
\item $\Gmat$ does not depend on $\sigma$, so it is a function of the
  declaration alone; and
\item the tuple
  $\Recv=\big(\Modes;\ \Inp=\InpC\sqcup\InpS;\ \{\Port{k}\}_{k=1}^{K};\
  \Gmat\big)$ is a receiver instrument in the sense of
  Definition~\ref{def:instrument}: it has the shape \eqref{eq:instrument} and
  satisfies its four clauses.
\end{enumerate}
\end{digestclaim}

\section*{4\quad Port Grams and the click-pattern POVM}\pdfbookmark[1]{4 Port Grams and the click-pattern POVM}{digest-group-receiver/thm:port-grams}

% claim receiver/thm:port-grams
\digestlabel{thm:port-grams}{4.1}\pdfbookmark[2]{Theorem 4.1 Port Grams and dilation independence}{digest-receiver/thm:port-grams}
\renewcommand*{\RestatedName}{\hyperref{contract.pdf}{}{claim-receiver/thm:port-grams}{Theorem~4.1}}%
\begin{digestclaim}[Port Grams and dilation independence]
Let $\Recv$ be a receiver instrument with
port Grams $\{A_k\}$ of
Definition~\ref{def:gram}.
\begin{enumerate}[label=(\roman*),leftmargin=3.0em]
\item \emph{Admissibility.} $A_k\succeq0$ for every $k$;
  $\sum_{k}A_k=\Gmat_{\Inp}^\dagger\Emd{[K]}\Gmat_{\Inp}\preceq
  \Gmat_{\Inp}^\dagger\Gmat_{\Inp}\preceq\id_{\Inp}$, so
  $\{A_k\}$ is admissible and the loss deficiency
  $\defic=\id_{\Inp}-A_\Sigma$ is positive semidefinite.
\item \emph{Rank.} $\rk A_k\le\min\big(|\Inp|,|\Port{k}|\big)$.
\item \emph{Pullback.} For every $J\subseteq[K]$,
  \begin{equation*}\tag{\getrefnumber{C-eq:pullback}}\digesteqlabel{eq:pullback}
    V^\dagger\,\Pi^0_J\,V=\Fock\big(\id_{\Inp}-A_J\big).
  \end{equation*}
\item \emph{Vacuum elimination is exact.} The right-hand side of
  \eqref{eq:pullback} depends on the instrument only through
  $\Gmat_{\Inp}$, the columns of $\Gmat$ indexed by the
  declared input modes. Dropping the other columns is therefore an identity, not an
  approximation.
\item \emph{Dilation independence.} \eqref{eq:pullback} holds for
  \emph{every} unitary $U$ on $\mathbb{C}^{\Modes\sqcup\Env'}$, with any
  finite mode set $\Env'$, whose block from the inputs $\Inp$ to the outputs
  $\bigcup_k\Port{k}$ equals the corresponding block of $\Gmat$. In particular
  the sub-unitary description (a contraction $\Gmat$ on the declared modes) and
  the unitary-dilation description (a unitary on system plus environment, with
  vacuum environment and the environment discarded) induce the \emph{same}
  measurement on $\mathcal{F}(\mathbb{C}^{\Inp})$.
\end{enumerate}
\end{digestclaim}

% claim receiver/lem:deficiency
\digestlabel{lem:deficiency}{4.2}\pdfbookmark[2]{Lemma 4.2 Decomposition of the loss deficiency}{digest-receiver/lem:deficiency}
\renewcommand*{\RestatedName}{\hyperref{contract.pdf}{}{claim-receiver/lem:deficiency}{Lemma~4.2}}%
\begin{digestclaim}[Decomposition of the loss deficiency]
Let $\Recv$ be a receiver instrument with port Grams $\{A_k\}$ and let
$\Emd{\mathrm{un}}$ be the projector onto the declared modes lying in no
port. Then
\begin{equation*}\tag{\getrefnumber{C-eq:deficiencyid}}\digesteqlabel{eq:deficiencyid}
  \defic\;=\;\id_{\Inp}-A_\Sigma
  \;=\;\underbrace{\big(\id_{\Inp}
    -\Gmat_{\Inp}^\dagger\Gmat_{\Inp}\big)}_{\text{left through the
    dilation}}
  \;+\;\underbrace{\Gmat_{\Inp}^\dagger\,\Emd{\mathrm{un}}\,
    \Gmat_{\Inp}}_{\text{reached an unmonitored declared mode}} ,
\end{equation*}
and both summands are positive semidefinite.
\end{digestclaim}

% claim receiver/def:clickalphabet
\digestlabel{def:clickalphabet}{4.4}\pdfbookmark[2]{Definition 4.4 Click-pattern alphabet}{digest-receiver/def:clickalphabet}
\renewcommand*{\RestatedName}{\hyperref{contract.pdf}{}{claim-receiver/def:clickalphabet}{Definition~4.4}}%
\begin{digestclaim}[Click-pattern alphabet]
An \emph{outcome} is the set $S\subseteq[K]$ of ports that click; the
\emph{raw alphabet} is $2^{[K]}$, in bijection with $\{0,1\}^{K}$ by indicator
vectors. We write
\begin{equation*}\tag{\getrefnumber{C-eq:mcsc}}\digesteqlabel{eq:mcsc}
  \SCK\coloneqq\big\{\{k\}:k\in[K]\big\},\qquad
  \MCK\coloneqq\big\{S\subseteq[K]:|S|\ge2\big\}
\end{equation*}
for the \emph{single-click set} and the \emph{multiclick set}, so that
$2^{[K]}=\{\varnothing\}\sqcup\SCK\sqcup\MCK$, and $\varnothing$ is the
\emph{no-click outcome}.
\end{digestclaim}

% claim receiver/thm:click-povm
\digestlabel{thm:click-povm}{4.5}\pdfbookmark[2]{Theorem 4.5 Click-pattern POVM}{digest-receiver/thm:click-povm}
\begin{mainclaim}
\renewcommand*{\RestatedName}{\hyperref{contract.pdf}{}{claim-receiver/thm:click-povm}{Theorem~4.5}}%
\begin{digestclaim}[Click-pattern POVM]
Let $\Recv$ be a receiver instrument with measurement isometry $V$
(Definition~\ref{def:physmodel}) and port Grams $\{A_k\}$. Define
\begin{equation*}\tag{\getrefnumber{C-eq:MSdef}}\digesteqlabel{eq:MSdef}
  M_S\coloneqq V^\dagger\Big[\prod_{k\in S}\big(\id-\Pi^0_k\big)
  \prod_{k\notin S}\Pi^0_k\Big]V ,\qquad S\subseteq[K].
\end{equation*}
Then:
\begin{enumerate}[label=(\roman*),leftmargin=3.0em]
\item \emph{$\{M_S\}$ is a POVM.} $M_S\succeq0$ for every $S$ and
  $\sum_{S\subseteq[K]}M_S=\id$. Moreover
  $\Pvac{J}\coloneqq V^\dagger\Pi^0_JV=\sum_{S\subseteq[K]\setminus J}M_S$ for
  every $J\subseteq[K]$.
\item \emph{Closed form.} $\Pvac{J}=\Fock(\id_{\Inp}-A_J)$ and
  \begin{equation*}\tag{\getrefnumber{C-eq:incexc}}\digesteqlabel{eq:incexc}
    M_S=\sum_{T\subseteq S}(-1)^{|T|}\,
        \Fock\big(\id_{\Inp}-A_{([K]\setminus S)\cup T}\big).
  \end{equation*}
\item \emph{Photon-number block structure.} Every $M_S$ commutes with the total
  photon number of the declared inputs, so
  $M_S=\bigoplus_{n\ge0}\blk{M}{\eqn}_S$ with
  $\blk{M}{\eqn}_S$ acting on $\Symn{n}(\mathbb{C}^{\Inp})$, and
  \begin{equation*}\tag{\getrefnumber{C-eq:blockMS}}\digesteqlabel{eq:blockMS}
    \blk{M}{\eqn}_S=\sum_{T\subseteq S}(-1)^{|T|}
    \Big(\big(\id-A_{([K]\setminus S)\cup T}\big)^{\otimes n}\Big)
    \Big|_{\Symn{n}} .
  \end{equation*}
  In particular $\blk{M}{(=0)}_S=\delta_{S,\varnothing}$: no photon, no click.
\item \emph{The coarse elements.} With $A_\Sigma=A_{[K]}$,
  \begin{align*}
    M_{\nc}&\coloneqq M_{\varnothing}=\Fock(\id-A_\Sigma),
      \tag{\getrefnumber{C-eq:Mnc}}\digesteqlabel{eq:Mnc}\\
    M_{\singlec_k}&\coloneqq M_{\{k\}}
      =\Fock\big(\id-A_{[K]\setminus\{k\}}\big)-\Fock(\id-A_\Sigma),
      \tag{\getrefnumber{C-eq:Msc}}\digesteqlabel{eq:Msc}\\
    M_{\mc}&\coloneqq\sum_{S\in\MCK}M_S
      =\id+(K-1)\,\Fock(\id-A_\Sigma)
        -\sum_{k=1}^{K}\Fock\big(\id-A_{[K]\setminus\{k\}}\big).
      \tag{\getrefnumber{C-eq:Mmc}}\digesteqlabel{eq:Mmc}
  \end{align*}
\end{enumerate}
\end{digestclaim}
\end{mainclaim}

\section*{5\quad Representation independence and coarse-graining}\pdfbookmark[1]{5 Representation independence and coarse-graining}{digest-group-receiver/thm:representation-independence}

% claim receiver/thm:representation-independence
\digestlabel{thm:representation-independence}{5.1}\pdfbookmark[2]{Theorem 5.1 Representation independence of the click-pattern POVM}{digest-receiver/thm:representation-independence}
\renewcommand*{\RestatedName}{\hyperref{contract.pdf}{}{claim-receiver/thm:representation-independence}{Theorem~5.1}}%
\begin{digestclaim}[Representation independence of the click-pattern POVM]
Let $\Recv,\Recv'$ be receiver instruments on the same declared input space
$\mathbb{C}^{\Inp}$ with the same number of ports $K$.
\begin{enumerate}[label=(\roman*),leftmargin=3.0em]
\item \emph{Determination, in both directions.} $\Recv$ and $\Recv'$ induce the
  same click-pattern POVM if and only if $A_k=A_k'$ for every $k$. In
  particular the POVM depends on the instrument matrix only through the
  port Grams.
\item \emph{Realizability.} Every admissible family $\{A_k\}_{k=1}^{K}$ on
  $\mathbb{C}^{\Inp}$ (Definition~\ref{def:gram}) is the port-Gram family of
  some receiver instrument, in which port $k$ collects at most
  $\rk A_k\le|\Inp|$ modes.
\item \emph{Port-local unitaries are invisible.} If
  $\Gmat'=\big(\bigoplus_kV_k\oplus V_{\mathrm{rest}}\big)\Gmat$ with each
  $V_k$ unitary on $\mathbb{C}^{\Port{k}}$ and $V_{\mathrm{rest}}$ unitary on
  the remaining modes, then $A_k'=A_k$ for every $k$.
\item \emph{One mode per port.} There is a receiver instrument with
  $|\Port{k}|=1$ for every $k$ inducing the same POVM as $\Recv$ if and only if
  $\rk A_k\le1$ for every $k$. In that case, writing
  $A_k=\lambda_ku_ku_k^\dagger$ with $\lambda_k\ge0$ and $u_k$ a unit vector,
  the single row of port $k$ is $g_k=e^{i\varphi_k}\sqrt{\lambda_k}\,
  u_k^\dagger$, unique up to the phase $\varphi_k$.
\end{enumerate}
\end{digestclaim}

% claim receiver/cor:phases
\digestlabel{cor:phases}{5.2}\pdfbookmark[2]{Corollary 5.2 Unobservability of per-port phases}{digest-receiver/cor:phases}
\renewcommand*{\RestatedName}{\hyperref{contract.pdf}{}{claim-receiver/cor:phases}{Corollary~5.2}}%
\begin{digestclaim}[Unobservability of per-port phases]
\begin{enumerate}[label=(\roman*),leftmargin=3.0em]
\item Multiplying the output row of a port by a phase, or more generally
  applying any unitary \emph{inside} a single port's mode set, changes no
  click statistic (Theorem~\ref{thm:representation-independence}\,(iii)).
\item \emph{The hypothesis matters}: the unitary must act inside one port's
  modes and after all interference has taken place. A phase that is global
  on an element in the \emph{middle} of the receiver can become relative once
  that element's output paths are recombined with other paths, and it can
  then change the Grams.
\end{enumerate}
\end{digestclaim}

% claim receiver/cor:extinction
\digestlabel{cor:extinction}{5.3}\pdfbookmark[2]{Corollary 5.3 Extinction criterion for rank one at two input modes}{digest-receiver/cor:extinction}
\renewcommand*{\RestatedName}{\hyperref{contract.pdf}{}{claim-receiver/cor:extinction}{Corollary~5.3}}%
\begin{digestclaim}[Extinction criterion for rank one at two input modes]
Suppose $|\Inp|=2$. Then $\rk A_k\le1$ holds if and only if either $A_k=0$ or
there is a unit vector $v\in\mathbb{C}^{2}$ with $v^\dagger A_kv=0$, that is,
a polarization the port extinguishes completely. For $|\Inp|>2$ the two
conditions are \emph{not} equivalent: the existence of one null direction of a
positive semidefinite $A_k$ says only $\rk A_k\le|\Inp|-1$.
\end{digestclaim}

% claim receiver/def:coarse
\digestlabel{def:coarse}{5.5}\pdfbookmark[2]{Definition 5.5 Declared coarse-graining and the emitted POVM}{digest-receiver/def:coarse}
\renewcommand*{\RestatedName}{\hyperref{contract.pdf}{}{claim-receiver/def:coarse}{Definition~5.5}}%
\begin{digestclaim}[Declared coarse-graining and the emitted POVM]
A \emph{declared coarse-graining} is a surjection
$\cg:2^{[K]}\to\Vset$ onto the declared emitted alphabet, supplied as data by
the protocol declaration, $[K]$ indexing the detector ports of a
receiver instrument (Definition~\ref{def:instrument}). With $M_S$ the click-pattern elements of
Theorem~\ref{thm:click-povm} (\eqref{eq:MSdef}), the \emph{emitted POVM} is
\begin{equation*}\tag{\getrefnumber{C-eq:coarse}}\digesteqlabel{eq:coarse}
  F_v\coloneqq\sum_{S\in\cg^{-1}(v)}M_S ,\qquad v\in\Vset .
\end{equation*}
\end{digestclaim}

% claim receiver/thm:coarse-graining
\digestlabel{thm:coarse-graining}{5.6}\pdfbookmark[2]{Theorem 5.6 Coarse-graining of the click-pattern POVM}{digest-receiver/thm:coarse-graining}
\begin{mainclaim}
\renewcommand*{\RestatedName}{\hyperref{contract.pdf}{}{claim-receiver/thm:coarse-graining}{Theorem~5.6}}%
\begin{digestclaim}[Coarse-graining of the click-pattern POVM]
Let $\cg$ be a declared coarse-graining.
\begin{enumerate}[label=(\roman*),leftmargin=3.0em]
\item $\{F_v\}_{v\in\Vset}$ is a POVM on $\mathcal{F}(\mathbb{C}^{\Inp})$, and
  every $F_v$ is block diagonal in the photon number with
  $\blk{F}{\eqn}_v=\sum_{S\in\cg^{-1}(v)}\blk{M}{\eqn}_S$.
\item \emph{Bounds add.} If $\mathcal{O}\subseteq2^{[K]}$ and
  $\blk{M}{\gtN}_S\succeq p_S\,\id^{(>N)}$ for every $S\in\mathcal{O}$ with
  $p_S\ge0$, then
  $\sum_{S\in\mathcal{O}}\blk{M}{\gtN}_S\succeq
  \big(\sum_{S\in\mathcal{O}}p_S\big)\id^{(>N)}$. In particular, if
  $\cg^{-1}(v)\supseteq\mathcal{O}$ then
  $\blk{F}{\gtN}_v\succeq\big(\sum_{S\in\mathcal{O}}p_S\big)\id^{(>N)}$.
\item \emph{Coarsening is lossy in one direction only.} A bound on a coarse
  outcome does not descend to its fibres: $F_v\succeq p\,\id$ implies nothing
  about the individual $M_S$ with $\cg(S)=v$.
\end{enumerate}
\end{digestclaim}
\end{mainclaim}

\section*{6\quad The analytic high-photon bound}\pdfbookmark[1]{6 The analytic high-photon bound}{digest-group-receiver/def:bounddata}

% claim receiver/def:bounddata
\digestlabel{def:bounddata}{6.1}\pdfbookmark[2]{Definition 6.1 Admissible bound data and the bound function}{digest-receiver/def:bounddata}
\renewcommand*{\RestatedName}{\hyperref{contract.pdf}{}{claim-receiver/def:bounddata}{Definition~6.1}}%
\begin{digestclaim}[Admissible bound data and the bound function]
A tuple $\big(\gamma,\{\beta_k\},c,\{d_k\}\big)$ with $K\ge1$ entries is
\emph{admissible bound data} when
\begin{equation*}\tag{\getrefnumber{C-eq:bounddata}}\digesteqlabel{eq:bounddata}
  0\le\gamma\le\beta_k\le1,\qquad 0\le c\le d_k\le1
  \qquad(k=1,\dots,K),
\end{equation*}
and the associated \emph{bound function} is
\begin{equation*}\tag{\getrefnumber{C-eq:beta}}\digesteqlabel{eq:beta}
  \Bnd(n)\coloneqq 1+(K-1)\,\gamma\,c^{\,n}-\sum_{k=1}^{K}\beta_k\,d_k^{\,n},
  \qquad n\ge0 ,
\end{equation*}
with the convention $0^0=1$. Admissibility \eqref{eq:bounddata} is exact. The
one part of it that this contract decides on computed data is $0\le c$; the
other parts are not separately tested (Definition~\ref{def:tolmap}\,(iv)).
In floating point that part is decided by
the widened test $c\ge-\varepsilon_2$ of Definition~\ref{def:tolmap}\,(iv),
which decides it on \emph{every} input route and is not displaced by the
route-(ii) admission of Definition~\ref{def:tolmap}\,(v); the computed $c$
enters \eqref{eq:beta} unclamped.
\end{digestclaim}

% claim receiver/thm:multiclick-bound
\digestlabel{thm:multiclick-bound}{6.2}\pdfbookmark[2]{Theorem 6.2 Lower bound on every multiclick block}{digest-receiver/thm:multiclick-bound}
\renewcommand*{\RestatedName}{\hyperref{contract.pdf}{}{claim-receiver/thm:multiclick-bound}{Theorem~6.2}}%
\begin{digestclaim}[Lower bound on every multiclick block]
Let $B,B_1,\dots,B_K$ be Hermitian operators on a finite-dimensional space with
\begin{equation*}\tag{\getrefnumber{C-eq:bbkhyp}}\digesteqlabel{eq:bbkhyp}
  0\ \preceq\ B_k\ \preceq\ B\ \preceq\ \id\qquad(k=1,\dots,K),
\end{equation*}
let $\gamma,\beta_1,\dots,\beta_K\in(0,1]$ with $\gamma\le\beta_k$ for every
$k$, let $\Pvac{[K]}$ and $\Pvac{[K]\setminus\{k\}}$ ($k=1,\dots,K$) be
operators on the Fock space of that space of the form
\begin{equation*}\tag{\getrefnumber{C-eq:prefactorform}}\digesteqlabel{eq:prefactorform}
  \Pvac{[K]}=\gamma\,\Fock(\id-B),\qquad
  \Pvac{[K]\setminus\{k\}}=\beta_k\,\Fock\big(\id-B_k\big),
\end{equation*}
and put $M_{\mc}\coloneqq\id+(K-1)\Pvac{[K]}-\sum_k\Pvac{[K]\setminus\{k\}}$. Put
\begin{equation*}\tag{\getrefnumber{C-eq:cd}}\digesteqlabel{eq:cd}
  c\coloneqq1-\lmax(B),\qquad
  d_k\coloneqq1-\lmin\big(B_k\big).
\end{equation*}
Then $\big(\gamma,\{\beta_k\},c,\{d_k\}\big)$ is admissible bound data, and for
every $n\ge0$
\begin{equation*}\tag{\getrefnumber{C-eq:mcbound}}\digesteqlabel{eq:mcbound}
  \lmin\big(\blk{M}{\eqn}_{\mc}\big)\ \ge\ \Bnd(n) ,
  \qquad\text{equivalently}\qquad
  \blk{M}{\eqn}_{\mc}\ \succeq\ \Bnd(n)\,\id^{(=n)} .
\end{equation*}
Moreover $\Bnd(n)\le1$ for every $n$.
\end{digestclaim}

% claim receiver/lem:collapse
\digestlabel{lem:collapse}{6.5}\pdfbookmark[2]{Lemma 6.5 Collapse of the tail infimum to the first photon number}{digest-receiver/lem:collapse}
\renewcommand*{\RestatedName}{\hyperref{contract.pdf}{}{claim-receiver/lem:collapse}{Lemma~6.5}}%
\begin{digestclaim}[Collapse of the tail infimum to the first photon number]
Let $\big(\gamma,\{\beta_k\},c,\{d_k\}\big)$ be admissible bound data
satisfying
\begin{equation*}\tag{\getrefnumber{C-eq:collapsehyp}}\digesteqlabel{eq:collapsehyp}
  \sum_{k=1}^{K}\beta_k\ \ge\ 1+(K-1)\gamma
  \qquad\text{(equivalently, }\Bnd(0)\le0\text{)}.
\end{equation*}
Then for every $n\ge1$,
\begin{equation*}\tag{\getrefnumber{C-eq:collapsestep}}\digesteqlabel{eq:collapsestep}
  \Bnd(n)\ge0\ \Longrightarrow\ \Bnd(n+1)\ \ge\ \Bnd(n) .
\end{equation*}
Consequently, for every cutoff $N\ge0$,
\begin{equation*}\tag{\getrefnumber{C-eq:collapseinf}}\digesteqlabel{eq:collapseinf}
  \max\Big\{0,\ \inf_{n>N}\Bnd(n)\Big\}\;=\;\max\big\{0,\ \Bnd(N+1)\big\} ,
\end{equation*}
and if $\Bnd(N+1)\ge0$ then $\Bnd$ is non-decreasing on $\{N+1,N+2,\dots\}$ and
$\inf_{n>N}\Bnd(n)=\Bnd(N+1)$. Hypothesis \eqref{eq:collapsehyp} holds
automatically when $\gamma=\beta_k=1$, where it reads $K\ge K$.
\end{digestclaim}

% claim receiver/lem:windowtail
\digestlabel{lem:windowtail}{6.8}\pdfbookmark[2]{Lemma 6.8 Window-and-tail bound on the tail infimum}{digest-receiver/lem:windowtail}
\renewcommand*{\RestatedName}{\hyperref{contract.pdf}{}{claim-receiver/lem:windowtail}{Lemma~6.8}}%
\begin{digestclaim}[Window-and-tail bound on the tail infimum]
Let $\big(\gamma,\{\beta_k\},c,\{d_k\}\big)$ be admissible bound data and
$N\ge0$. Define, for an integer $m>N$,
\begin{equation*}\tag{\getrefnumber{C-eq:windowtail}}\digesteqlabel{eq:windowtail}
  \underline{\Bnd}_{N,m}\coloneqq
  \min\big\{\,W_{N,m},\ T_m\,\big\},\qquad
  W_{N,m}\coloneqq\min_{N<n\le m}\Bnd(n),\qquad
  T_m\coloneqq1-\sum_{k=1}^{K}\beta_kd_k^{\,m} .
\end{equation*}
The window $\{n:N<n\le m\}$ is non-empty for every $m>N$; at $m=N+1$ it is the
single index $N+1$. Then:
\begin{enumerate}[label=(\roman*),leftmargin=3.0em]
\item $\underline{\Bnd}_{N,m}\le\inf_{n>N}\Bnd(n)$ for every $m>N$: the
  quantity is a valid uniform bound;
\item $\underline{\Bnd}_{N,m}$ is non-decreasing in $m$, and
  $\lim_{m\to\infty}\underline{\Bnd}_{N,m}=\inf_{n>N}\Bnd(n)$ whenever
  $c<1$; moreover the \emph{residual gap}
  \begin{equation*}\tag{\getrefnumber{C-eq:resgap}}\digesteqlabel{eq:resgap}
    g_{N,m}\coloneqq\max\big\{0,\ W_{N,m}-T_m\big\}
  \end{equation*}
  is a finite non-negative real for every $m>N$ --- at $m=N+1$ it equals
  $(K-1)\gamma c^{\,N+1}$ --- and satisfies
  $\inf_{n>N}\Bnd(n)-\underline{\Bnd}_{N,m}\le g_{N,m}$, so the evaluation is
  convergent and $g_{N,m}$ bounds what widening the window could still recover;
\item if $c=0$ then $\Bnd(n)=1-\sum_k\beta_kd_k^{\,n}$ for $n\ge1$ is
  non-decreasing, $\underline{\Bnd}_{N,N+1}=\Bnd(N+1)$ is exact, and
  $g_{N,N+1}=0$;
\item the two lemmas are consistent: under the hypothesis of
  Lemma~\ref{lem:collapse} with $\Bnd(N+1)\ge0$ \emph{and $c<1$}, (i) and (ii)
  give $\underline{\Bnd}_{N,m}\le\Bnd(N+1)$ for every $m$ with equality in the
  limit $m\to\infty$, so the first branch of Definition~\ref{def:pmc} is the
  exact value the window-and-tail evaluation converges to, and using it loses
  nothing. The hypothesis $c<1$ is needed: at $c=1$ admissibility forces
  $d_k=1$ for every $k$, $\Bnd$ is constant, and under
  Lemma~\ref{lem:collapse}'s hypothesis with $\Bnd(N+1)\ge0$ that constant is
  $0$ while $T_m=1-\sum_k\beta_k=-(K-1)\gamma$ for every $m$, so the limit
  clause fails whenever $\gamma>0$ and $K\ge2$.
\end{enumerate}
\end{digestclaim}

\section*{7\quad Exact block eigenvalues}\pdfbookmark[1]{7 Exact block eigenvalues}{digest-group-receiver/def:exactbnd}

% claim receiver/def:exactbnd
\digestlabel{def:exactbnd}{7.1}\pdfbookmark[2]{Definition 7.1 The exact block bound function}{digest-receiver/def:exactbnd}
\renewcommand*{\RestatedName}{\hyperref{contract.pdf}{}{claim-receiver/def:exactbnd}{Definition~7.1}}%
\begin{digestclaim}[The exact block bound function]
Under the hypotheses of Theorem~\ref{thm:multiclick-bound}, the
\emph{exact block bound}
is
\begin{equation*}\tag{\getrefnumber{C-eq:exactbnd}}\digesteqlabel{eq:exactbnd}
  \Bnd^{\ast}(n)\ \coloneqq\ \lmin\big(\blk{M}{\eqn}_{\mc}\big)
  \ =\ \lmin\Big(\id^{(=n)}+(K-1)\gamma\,\Fock(\id-B)^{(n)}
   -\sum_{k=1}^{K}\beta_k\,\Fock\big(\id-B_k\big)^{(n)}\Big),
\end{equation*}
for $n\ge0$, the second equality being \eqref{eq:prefactorform} restricted to
the $n$-photon block. The operator inside is a Hermitian matrix of dimension
$\binom{D+n-1}{n}$, where $D$ is the dimension of the space $B$ acts on:
the channel input space of the setting family
(Definition~\ref{def:setting-family}), which is $|\InpC|$ at both instances
this module constructs. No new object is introduced: each
$\Fock(\cdot)^{(n)}=\big((\cdot)^{\otimes n}\big)\big|_{\Symn{n}}$ is the
construction \eqref{eq:blockMS} already uses for the blocks the artifact
emits at $n\le N$, in the occupation basis and order that
Definition~\ref{def:fock} and the interface conventions of this module fix,
evaluated at an index above the cutoff instead of below it.
\end{digestclaim}

% claim receiver/lem:exactbasic
\digestlabel{lem:exactbasic}{7.2}\pdfbookmark[2]{Lemma 7.2 Range of the exact block bound}{digest-receiver/lem:exactbasic}
\renewcommand*{\RestatedName}{\hyperref{contract.pdf}{}{claim-receiver/lem:exactbasic}{Lemma~7.2}}%
\begin{digestclaim}[Range of the exact block bound]
Let the hypotheses of Theorem~\ref{thm:multiclick-bound} hold. Then, for every
$n\ge0$,
\begin{equation*}\tag{\getrefnumber{C-eq:exactbasic}}\digesteqlabel{eq:exactbasic}
  \Bnd(n)\ \le\ \Bnd^{\ast}(n)\ \le\ 1 .
\end{equation*}
If in addition
\begin{equation*}\tag{\getrefnumber{C-eq:mcpos}}\digesteqlabel{eq:mcpos}
  M_{\mc}\ \succeq\ 0 ,
\end{equation*}
then $\Bnd^{\ast}(n)\ge0$ as well, so that $\Bnd^{\ast}(n)\in[0,1]$.
\end{digestclaim}

% claim receiver/lem:exactwindow
\digestlabel{lem:exactwindow}{7.4}\pdfbookmark[2]{Lemma 7.4 Exact window with an analytic remainder}{digest-receiver/lem:exactwindow}
\renewcommand*{\RestatedName}{\hyperref{contract.pdf}{}{claim-receiver/lem:exactwindow}{Lemma~7.4}}%
\begin{digestclaim}[Exact window with an analytic remainder]
Let the hypotheses of Theorem~\ref{thm:multiclick-bound} hold together with
\eqref{eq:mcpos}, let $N\ge0$, and let
$m^{\ast}>N$ be an integer. Let $\ell(n)$, $N<n\le m^{\ast}$, be reals with
\begin{equation*}\tag{\getrefnumber{C-eq:ellhyp}}\digesteqlabel{eq:ellhyp}
  \ell(n)\ \le\ \Bnd^{\ast}(n)\qquad\text{for }N<n\le m^{\ast},
\end{equation*}
and let $\tau$ be a real with
\begin{equation*}\tag{\getrefnumber{C-eq:tailadm}}\digesteqlabel{eq:tailadm}
  \tau\le0
  \qquad\text{or}\qquad
  \tau\le\inf_{n>m^{\ast}}\Bnd^{\ast}(n).
\end{equation*}
Then
\begin{equation*}\tag{\getrefnumber{C-eq:exactcomposite}}\digesteqlabel{eq:exactcomposite}
  P\ \coloneqq\ \max\Big\{0,\ \min\Big\{\ \min_{N<n\le m^{\ast}}\ell(n),
   \ \ \tau\ \Big\}\Big\}
\end{equation*}
satisfies $0\le P\le1$ and
\begin{equation*}\tag{\getrefnumber{C-eq:exactstatement}}\digesteqlabel{eq:exactstatement}
  \blk{M}{\gtN}_{\mc}\ \succeq\ P\,\id^{(>N)} .
\end{equation*}
Three choices are admissible:
\begin{enumerate}[label=(\roman*),leftmargin=3.0em]
\item $\tau=T_{m^{\ast}}$ of \eqref{eq:windowtail}, formed from the bound
  data $\big(\gamma,\{\beta_k\},c,\{d_k\}\big)$ that
  Theorem~\ref{thm:multiclick-bound} forms by \eqref{eq:cd}, satisfies
  \eqref{eq:tailadm} with no further hypothesis;
\item if the bound data satisfy \eqref{eq:collapsehyp}, then
  $\tau=\Bnd(m^{\ast}+1)$ satisfies \eqref{eq:tailadm};
\item $\ell(n)=\Bnd(n)$ satisfies \eqref{eq:ellhyp}, and the pointwise maximum
  of finitely many admissible choices of $\ell$ is admissible.
\end{enumerate}
\end{digestclaim}

\section*{8\quad Setting families and the emitted weight}\pdfbookmark[1]{8 Setting families and the emitted weight}{digest-group-receiver/def:setting-family}

% claim receiver/def:setting-family
\digestlabel{def:setting-family}{8.1}\pdfbookmark[2]{Definition 8.1 Setting family}{digest-receiver/def:setting-family}
\renewcommand*{\RestatedName}{\hyperref{contract.pdf}{}{claim-receiver/def:setting-family}{Definition~8.1}}%
\begin{digestclaim}[Setting family]
Let $D\ge1$ and write $\mathcal{F}_D\coloneqq\mathcal{F}(\mathbb{C}^{D})$ for
the Fock space of Definition~\ref{def:fock} over a $D$-dimensional
\emph{channel input space},
with photon-number projectors $\Pi_n$ onto $\Symn{n}(\mathbb{C}^{D})$. A
\emph{setting family}
of $K\ge1$ ports on $\mathcal{F}_D$ is a finite non-empty set $\Tset$ of
\emph{settings} together with, for every $t\in\Tset$:
\begin{enumerate}[label=(\alph*),leftmargin=2.6em]
\item a POVM $\{M_{S|t}\}_{S\subseteq[K]}$ on $\mathcal{F}_D$ every
  element of which commutes with the photon number, so that
  $M_{S|t}=\bigoplus_{n\ge0}\blk{M}{\eqn}_{S|t}$ with $\blk{M}{\eqn}_{S|t}$
  acting on $\Symn{n}(\mathbb{C}^{D})$. Its
  \emph{joint no-click elements}
  and its
  \emph{multiclick element}
  are
  \begin{equation*}\tag{\getrefnumber{C-eq:family-elements}}\digesteqlabel{eq:family-elements}
    \Pvac{J}^{(t)}\coloneqq\sum_{S\subseteq[K]\setminus J}M_{S|t}
    \quad(J\subseteq[K]),\qquad
    M_{\mc|t}\coloneqq\sum_{S\in\MCK}M_{S|t},
  \end{equation*}
  with $\MCK$ the multiclick set of Definition~\ref{def:clickalphabet};
\item Hermitian operators $B^{(t)},B^{(t)}_1,\dots,B^{(t)}_K$ on
  $\mathbb{C}^{D}$ and scalars $\gamma^{(t)},\beta^{(t)}_1,\dots,
  \beta^{(t)}_K\in(0,1]$ satisfying the hypotheses of
  Theorem~\ref{thm:multiclick-bound}, namely \eqref{eq:bbkhyp} and
  $\gamma^{(t)}\le\beta^{(t)}_k$ for every $k$, such that the joint
  no-click elements have the prefactor form \eqref{eq:prefactorform},
  \begin{equation*}\tag{\getrefnumber{C-eq:family-prefactor}}\digesteqlabel{eq:family-prefactor}
    \Pvac{[K]}^{(t)}=\gamma^{(t)}\,\Fock\big(\id-B^{(t)}\big),\qquad
    \Pvac{[K]\setminus\{k\}}^{(t)}=\beta^{(t)}_k\,\Fock\big(\id-B^{(t)}_k\big);
  \end{equation*}
\item a flag $\chi^{(t)}\in\{\mathrm{exact},\mathrm{general}\}$, the
  setting's
  \emph{exact flag},
  which may be $\mathrm{exact}$ only if $\gamma^{(t)}=\beta^{(t)}_k=1$ for
  every $k$ holds \emph{exactly, by the construction that produced the
  setting}, read off the declaration with no arithmetic performed, so that
  at an exact setting \eqref{eq:collapsehyp} reads $K\ge K$ and holds with
  equality by construction.
\end{enumerate}
The
\emph{bound data of the setting}
are $\big(\gamma^{(t)},\{\beta^{(t)}_k\},c^{(t)},\{d^{(t)}_k\}\big)$ with
$c^{(t)}\coloneqq1-\lmax(B^{(t)})$ and $d^{(t)}_k\coloneqq1-\lmin(B^{(t)}_k)$
as in \eqref{eq:cd}; they are admissible bound data
(Definition~\ref{def:bounddata}) by Theorem~\ref{thm:multiclick-bound}, and
$\Bnd^{(t)}$ denotes their bound function \eqref{eq:beta}.

Two consequences are used throughout. First, the multiclick element has the
inclusion--exclusion form
\begin{equation*}\tag{\getrefnumber{C-eq:family-mc}}\digesteqlabel{eq:family-mc}
  M_{\mc|t}=\id+(K-1)\,\Pvac{[K]}^{(t)}-\sum_{k=1}^{K}\Pvac{[K]\setminus\{k\}}^{(t)} ,
\end{equation*}
because completeness gives $M_{\mc|t}=\id-M_{\varnothing|t}-\sum_kM_{\{k\}|t}$
while \eqref{eq:family-elements} gives $M_{\varnothing|t}=\Pvac{[K]}^{(t)}$
and $M_{\{k\}|t}=\Pvac{[K]\setminus\{k\}}^{(t)}-\Pvac{[K]}^{(t)}$; so
Theorem~\ref{thm:multiclick-bound} applies to every setting, with
$\lmin(\blk{M}{\eqn}_{\mc|t})\ge\Bnd^{(t)}(n)$ for every $n\ge0$. Second,
$M_{\mc|t}$ is a sum of elements of a POVM, so hypothesis \eqref{eq:mcpos}
of Lemma~\ref{lem:exactbasic}, $M_{\mc|t}\succeq0$, holds at every setting.
\end{digestclaim}

% claim receiver/cor:unconditioned-family
\digestlabel{cor:unconditioned-family}{8.2}\pdfbookmark[2]{Corollary 8.2 The unconditioned receiver is a setting family}{digest-receiver/cor:unconditioned-family}
\renewcommand*{\RestatedName}{\hyperref{contract.pdf}{}{claim-receiver/cor:unconditioned-family}{Corollary~8.2}}%
\begin{digestclaim}[The unconditioned receiver is a setting family]
Let $\Recv$ be a receiver instrument (Definition~\ref{def:instrument})
with $\InpS=\varnothing$, so that $\Inp=\InpC$, and let $\{A_k\}$ be its
port Grams and $\{M_S\}$ its click-pattern POVM
(Theorem~\ref{thm:click-povm}). Then $\Tset=\{\off\}$ with
\begin{equation*}\tag{\getrefnumber{C-eq:unconditioned-data}}\digesteqlabel{eq:unconditioned-data}
  M_{S|\off}\coloneqq M_S,\qquad
  B^{(\off)}\coloneqq A_\Sigma,\qquad
  B^{(\off)}_k\coloneqq A_{[K]\setminus\{k\}},\qquad
  \gamma^{(\off)}=\beta^{(\off)}_k\coloneqq1,\qquad
  \chi^{(\off)}\coloneqq\mathrm{exact},
\end{equation*}
is a setting family (Definition~\ref{def:setting-family}) on
$\mathcal{F}(\mathbb{C}^{\Inp})$, with $D=|\Inp|$. Its bound data are
$c=1-\lmax(A_\Sigma)$ and $d_k=1-\lmin(A_{[K]\setminus\{k\}})$, and
\eqref{eq:collapsehyp} holds with equality.
\end{digestclaim}

% claim receiver/def:pmc
\digestlabel{def:pmc}{8.3}\pdfbookmark[2]{Definition 8.3 The uniform high-photon multiclick weight}{digest-receiver/def:pmc}
\renewcommand*{\RestatedName}{\hyperref{contract.pdf}{}{claim-receiver/def:pmc}{Definition~8.3}}%
\begin{digestclaim}[The uniform high-photon multiclick weight]
Fix a cutoff $N$ and a setting family
(Definition~\ref{def:setting-family}); for each setting $t\in\Tset$ let
$\Bnd^{(t)}$ be the bound function \eqref{eq:beta} of the
bound data of the setting. At every setting $t$ whose declaration
carries \emph{no} exact-window depth (Definition~\ref{def:pmcex}),
the
\emph{per-setting weight}
is
\begin{equation*}\tag{\getrefnumber{C-eq:pmcdef}}\digesteqlabel{eq:pmcdef}
  \pmc^{(t)}\coloneqq
  \begin{cases}
    \max\big\{0,\ \Bnd^{(t)}(N+1)\big\}
      &\text{if (X) holds at }t,\\[4pt]
    \max\big\{0,\ \underline{\Bnd}^{(t)}_{N,m_t}\big\}
      &\text{otherwise, for the declared window }m_t>N,
  \end{cases}
\end{equation*}
and the
\emph{uniform weight}
is $\pmc\coloneqq\min_{t\in\Tset}\pmc^{(t)}$, where the branch-1
\emph{predicate (X)}
is
\begin{description}[leftmargin=3.2em,style=nextline]
\item[(X) \emph{exact by construction, decidable from the declaration:}] the
  setting's exact flag is $\chi^{(t)}=\mathrm{exact}$
  (Definition~\ref{def:setting-family}\,(c)). Then
  $\gamma^{(t)}=\beta^{(t)}_k=1$ exactly, so \eqref{eq:collapsehyp} reads
  $K\ge K$ and holds with \emph{equality}, no arithmetic having been
  performed.
\end{description}
\emph{Where the exact window replaces this.} At a setting whose declaration
carries an exact-window depth $m^{\ast}_t$, \eqref{eq:pmcdef} does not apply
and $\pmc^{(t)}$ is \eqref{eq:pmcex} instead
(Definition~\ref{def:pmcex}); $\pmc=\min_{t\in\Tset}\pmc^{(t)}$ is unchanged,
and the two kinds of setting may be mixed in one declaration. A declaration
carrying no depth at any setting is governed by \eqref{eq:pmcdef} alone: no
$\Bnd^{\ast}$ is formed, no eigenvalue of any block above the cutoff is
computed, and Assumption~\ref{ass:eig} is not invoked.
\end{digestclaim}

% claim receiver/def:pmcex
\digestlabel{def:pmcex}{8.7}\pdfbookmark[2]{Definition 8.7 Exact-window weight of a setting with a declared depth}{digest-receiver/def:pmcex}
\renewcommand*{\RestatedName}{\hyperref{contract.pdf}{}{claim-receiver/def:pmcex}{Definition~8.7}}%
\begin{digestclaim}[Exact-window weight of a setting with a declared depth]
Let $t\in\Tset$ be a setting whose declaration carries an
\emph{exact-window depth}
$m^{\ast}_t$: an integer with $m^{\ast}_t>N$ and, where predicate~(X)
fails at $t$, also $m^{\ast}_t\le m_t$. Write $\widehat\Bnd^{\ast(t)}(n)$ for
the value an implementation computes for $\Bnd^{\ast(t)}(n)$ of
\eqref{eq:exactbnd} at that setting, let $\varepsilon_{\mathrm{eig}}$ be the
declared
\emph{eigenvalue slack},
and put
\begin{equation*}\tag{\getrefnumber{C-eq:elldef}}\digesteqlabel{eq:elldef}
  \ell^{(t)}(n)\ \coloneqq\
  \max\big\{\ \Bnd^{(t)}(n),\ \ \widehat\Bnd^{\ast(t)}(n)
   -\varepsilon_{\mathrm{eig}}\ \big\},
  \qquad N<n\le m^{\ast}_t .
\end{equation*}
Then
\begin{equation*}\tag{\getrefnumber{C-eq:pmcex}}\digesteqlabel{eq:pmcex}
  \pmc^{(t)}\coloneqq
  \begin{cases}
    \max\Big\{0,\ \min\big\{\ \min\limits_{N<n\le m^{\ast}_t}\ell^{(t)}(n),
      \ \ \Bnd^{(t)}\big(m^{\ast}_t+1\big)\ \big\}\Big\}
      &\text{if (X) holds at }t,\\[10pt]
    \max\Big\{0,\ \min\big\{\ \min\limits_{N<n\le m^{\ast}_t}\ell^{(t)}(n),
      \ \ \min\limits_{m^{\ast}_t<n\le m_t}\Bnd^{(t)}(n),
      \ \ T^{(t)}_{m_t}\ \big\}\Big\}
      &\text{otherwise,}
  \end{cases}
\end{equation*}
the middle minimum in the second case running over the empty index set ---
and so equal to $+\infty$ and inactive --- when $m^{\ast}_t=m_t$.

\emph{Per setting.} The depth is optional and per setting, and each
per-setting weight $\pmc^{(t)}$ is its own setting's: a depth changes the
weight of the setting that declares it and of no other. A declaration that
carries a depth at some settings and not at others therefore obtains the
exact-window value as the weight of each setting with one, while the uniform
weight $\pmc=\min_{t\in\Tset}\pmc^{(t)}$ is bounded by the settings without
one.
\end{digestclaim}

% claim receiver/prop:exactnoworse
\digestlabel{prop:exactnoworse}{8.9}\pdfbookmark[2]{Proposition 8.9 Dominance of the exact-window weight over the analytic weight}{digest-receiver/prop:exactnoworse}
\renewcommand*{\RestatedName}{\hyperref{contract.pdf}{}{claim-receiver/prop:exactnoworse}{Proposition~8.9}}%
\begin{digestclaim}[Dominance of the exact-window weight over the analytic weight]
At every setting carrying a declared exact-window depth, the value
\eqref{eq:pmcex} is at least the value \eqref{eq:pmcdef} gives at that
setting on the same bound data.
\end{digestclaim}

% claim receiver/def:enclosed-weight
\digestlabel{def:enclosed-weight}{8.12}\pdfbookmark[2]{Definition 8.12 The weight evaluated on enclosed bound data}{digest-receiver/def:enclosed-weight}
\renewcommand*{\RestatedName}{\hyperref{contract.pdf}{}{claim-receiver/def:enclosed-weight}{Definition~8.12}}%
\begin{digestclaim}[The weight evaluated on enclosed bound data]
Fix a cutoff $N$, a setting family
(Definition~\ref{def:setting-family}) and a setting $t\in\Tset$, and write
$\big(\gamma,\{\beta_k\},c,\{d_k\}\big)$ for the
bound data of the setting, the index $t$ omitted.
\begin{enumerate}[label=(\roman*),leftmargin=2.4em]
\item Reals $\big(\check\gamma,\{\check\beta_k\},\check c,\{\check d_k\}\big)$
  are
  \emph{enclosed bound data} of $t$ when
  \begin{equation*}\tag{\getrefnumber{C-eq:enclosed-directions}}\digesteqlabel{eq:enclosed-directions}
    0\le\check\gamma\le\gamma,\qquad \beta_k\le\check\beta_k\le1,\qquad
    0\le\check c\le c,\qquad d_k\le\check d_k\le1
    \qquad(k=1,\dots,K).
  \end{equation*}
  $\check\Bnd$ denotes the bound function \eqref{eq:beta} of this
  tuple, and $\check T_m$ and $\underline{\check\Bnd}_{N,m}$ the quantities
  $T_m$ and $\underline{\Bnd}_{N,m}$ of \eqref{eq:windowtail} formed from it.
\item Where $t$ carries an exact-window depth $m^{\ast}_t$
  (Definition~\ref{def:pmcex}), reals $w(n)$, $N<n\le m^{\ast}_t$, are
  \emph{rigorous window bounds} of $t$ when
  $w(n)\le\Bnd^{\ast(t)}(n)$ for each such $n$, $\Bnd^{\ast(t)}$ being the
  exact block bound \eqref{eq:exactbnd} at $t$.
\item The
  \emph{enclosed weight} $\check p^{(t)}_{\mc}$ of $t$ formed
  from them is \eqref{eq:pmcdef} where $t$ carries no exact-window depth and
  \eqref{eq:pmcex} where it carries one, its branch chosen by
  predicate (X) at $t$ as there, each read with $\check\Bnd$,
  $\underline{\check\Bnd}_{N,m_t}$ and $\check T_{m_t}$ in place of
  $\Bnd^{(t)}$, $\underline{\Bnd}^{(t)}_{N,m_t}$ and $T^{(t)}_{m_t}$, and with
  \begin{equation*}\tag{\getrefnumber{C-eq:enclosed-ell}}\digesteqlabel{eq:enclosed-ell}
    \check\ell(n)\ \coloneqq\ \max\big\{\check\Bnd(n),\ w(n)\big\},
    \qquad N<n\le m^{\ast}_t ,
  \end{equation*}
  in place of $\ell^{(t)}(n)$ of \eqref{eq:elldef}. No computed eigenvalue of
  a block and no eigenvalue slack enters it.
\end{enumerate}
\end{digestclaim}

% claim receiver/thm:enclosed-weight
\digestlabel{thm:enclosed-weight}{8.13}\pdfbookmark[2]{Theorem 8.13 The weight on enclosed bound data bounds the multiclick element}{digest-receiver/thm:enclosed-weight}
\renewcommand*{\RestatedName}{\hyperref{contract.pdf}{}{claim-receiver/thm:enclosed-weight}{Theorem~8.13}}%
\begin{digestclaim}[The weight on enclosed bound data bounds the multiclick element]
Let a setting family (Definition~\ref{def:setting-family}) be given, let
$N\ge0$, and let $t\in\Tset$, with operators $B=B^{(t)}$ and $B_k=B^{(t)}_k$,
bound data of the setting $\big(\gamma,\{\beta_k\},c,\{d_k\}\big)$,
bound function $\Bnd=\Bnd^{(t)}$ and $T_m$ of \eqref{eq:windowtail}
formed from those data. Write
$\operatorname{clip}(x)\coloneqq\min\{1,\max\{0,x\}\}$ for real $x$.
\begin{enumerate}[label=(\roman*),leftmargin=2.4em]
\item \emph{Construction.} Let $u$, $\ell_1,\dots,\ell_K$, $g$ and
  $b_1,\dots,b_K$ be reals with $u\ge\lmax(B)$, $\ell_k\le\lmin(B_k)$,
  $g\le\gamma$ and $b_k\ge\beta_k$. Then
  \begin{equation*}\tag{\getrefnumber{C-eq:enclosed-construction}}\digesteqlabel{eq:enclosed-construction}
    \big(\operatorname{clip}(g),\ \{\operatorname{clip}(b_k)\},\
    \operatorname{clip}(1-u),\ \{\operatorname{clip}(1-\ell_k)\}\big)
  \end{equation*}
  are enclosed bound data of $t$. Where $t$ carries an
  exact-window depth $m^{\ast}_t$ and reals $\lambda_n$ satisfy
  $\lambda_n\le\lmin\big(\blk{M}{\eqn}_{\mc|t}\big)$ for
  $N<n\le m^{\ast}_t$, the reals $\max\{0,\lambda_n\}$ are
  rigorous window bounds of $t$.
\item \emph{Comparison.} Enclosed bound data
  $\big(\check\gamma,\{\check\beta_k\},\check c,\{\check d_k\}\big)$ of $t$
  are admissible bound data (Definition~\ref{def:bounddata}), and their
  bound function $\check\Bnd$ and $\check T_m$ satisfy
  \begin{equation*}\tag{\getrefnumber{C-eq:enclosed-comparison}}\digesteqlabel{eq:enclosed-comparison}
    \check\Bnd(n)\ \le\ \Bnd(n)\quad(n\ge0),\qquad
    \check T_m\ \le\ T_m\quad(m\ge0).
  \end{equation*}
\item \emph{The bound.} For enclosed bound data of $t$ and, where $t$
  carries an exact-window depth, rigorous window bounds of $t$, the
  enclosed weight $\check p^{(t)}_{\mc}$ formed from them
  (Definition~\ref{def:enclosed-weight}) satisfies
  $0\le\check p^{(t)}_{\mc}\le1$ and
  \begin{equation*}\tag{\getrefnumber{C-eq:enclosed-bound}}\digesteqlabel{eq:enclosed-bound}
    \blk{M}{\gtN}_{\mc|t}\ \succeq\ \check p^{(t)}_{\mc}\,\id^{(>N)} ,
  \end{equation*}
  and every real $p$ with $0\le p\le\check p^{(t)}_{\mc}$ satisfies
  $0\le p\le1$ and $\blk{M}{\gtN}_{\mc|t}\succeq p\,\id^{(>N)}$.
\end{enumerate}
\end{digestclaim}

% claim receiver/thm:emitted-bound
\digestlabel{thm:emitted-bound}{8.14}\pdfbookmark[2]{Theorem 8.14 Multiclick bounds carried by the emitted weights}{digest-receiver/thm:emitted-bound}
\begin{mainclaim}
\renewcommand*{\RestatedName}{\hyperref{contract.pdf}{}{claim-receiver/thm:emitted-bound}{Theorem~8.14}}%
\begin{digestclaim}[Multiclick bounds carried by the emitted weights]
Let a setting family (Definition~\ref{def:setting-family}) be given. For
every setting $t\in\Tset$ let $\pmc^{(t)}$ be as in
Definition~\ref{def:pmc} if $t$ carries no exact-window depth and as in
Definition~\ref{def:pmcex} if it carries one, or, at any setting, let it be
the enclosed weight $\check p^{(t)}_{\mc}$ of
Definition~\ref{def:enclosed-weight} formed from enclosed bound data of
$t$ and, where $t$ carries a depth, from rigorous window bounds of $t$,
or a real $p$ with $0\le p\le\check p^{(t)}_{\mc}$;
and let $\pmc=\min_{t\in\Tset}\pmc^{(t)}$ as in Definition~\ref{def:pmc}. Then
\begin{equation*}\tag{\getrefnumber{C-eq:pmcsetting}}\digesteqlabel{eq:pmcsetting}
  \blk{M}{\gtN}_{\mc|t}\ \succeq\ \pmc^{(t)}\,\id^{(>N)}
  \qquad\text{for every }t\in\Tset ,
\end{equation*}
where $\blk{M}{\gtN}_{\mc|t}$ is the restriction of the multiclick element
$M_{\mc|t}$ of \eqref{eq:family-elements} at setting $t$ to the photon-number
blocks $n>N$; and consequently, with the one weight $\pmc$ at every setting,
\begin{equation*}\tag{\getrefnumber{C-eq:pmcstatement}}\digesteqlabel{eq:pmcstatement}
  \blk{M}{\gtN}_{\mc|t}\ \succeq\ \pmc\,\id^{(>N)}
  \qquad\text{for every }t\in\Tset .
\end{equation*}
At a setting governed by Definition~\ref{def:pmc} or by
Definition~\ref{def:enclosed-weight} both inequalities hold at
that setting with no hypothesis beyond this contract's ledger; at a setting
governed by Definition~\ref{def:pmcex} they hold there \emph{under
Assumption~\ref{ass:eig}}, which is what makes the second entry of
\eqref{eq:elldef} a lower bound. At a
one-setting family, $\Tset=\{\off\}$, one has $\pmc=\pmc^{(\off)}$ and the
two inequalities are one. Moreover $0\le\pmc^{(t)}\le1$ at every setting,
hence $0\le\pmc\le1$, and each $\pmc^{(t)}$ of Definition~\ref{def:pmc}
or~\ref{def:pmcex} is computable by finitely many
arithmetic operations on
$\{\lmax(B^{(t)}),\lmin(B^{(t)}_k),\gamma^{(t)},\beta^{(t)}_k\}$ and the
declared $N$, $m_t$ and $m^{\ast}_t$ --- where $B^{(t)},B^{(t)}_k$ are the
operators of the setting (Definition~\ref{def:setting-family}\,(b)) ---
together with, where the setting carries a declared depth, the declared
$\varepsilon_{\mathrm{eig}}$ and the $m^{\ast}_t-N$ extreme eigenvalues
$\widehat\Bnd^{\ast(t)}(n)$ of the Hermitian matrices \eqref{eq:exactbnd},
each of dimension $\binom{D+n-1}{n}$; each enclosed weight by finitely many
on the enclosed bound data, the rigorous window bounds and the declared $N$,
$m_t$ and $m^{\ast}_t$; and $\pmc$ by finitely many comparisons among them.
\end{digestclaim}
\end{mainclaim}

% claim receiver/cor:emitted-unconditioned
\digestlabel{cor:emitted-unconditioned}{8.15}\pdfbookmark[2]{Corollary 8.15 Emitted multiclick weight of a probe-free receiver}{digest-receiver/cor:emitted-unconditioned}
\renewcommand*{\RestatedName}{\hyperref{contract.pdf}{}{claim-receiver/cor:emitted-unconditioned}{Corollary~8.15}}%
\begin{digestclaim}[Emitted multiclick weight of a probe-free receiver]
Let $\Recv$ be a receiver instrument with $\InpS=\varnothing$
(Definition~\ref{def:instrument}), let $N$ be the declared cutoff, and let
$\pmc$ be computed by Definition~\ref{def:pmc}, or by
Definition~\ref{def:pmcex} where an exact-window depth is declared, from the
one-setting family of Corollary~\ref{cor:unconditioned-family}, or let it be
a real with $0\le\pmc\le\check p_{\mc}$, where $\check p_{\mc}$ is an
enclosed weight of that family's one setting
(Definition~\ref{def:enclosed-weight}), $\check p_{\mc}$ itself included. Then
\begin{equation*}\tag{\getrefnumber{C-eq:emitted-unconditioned}}\digesteqlabel{eq:emitted-unconditioned}
  \blk{M}{\gtN}_{\mc}\ \succeq\ \pmc\,\id^{(>N)},\qquad 0\le\pmc\le1,
\end{equation*}
with no hypothesis beyond the ledger where $\pmc$ is given by
Definition~\ref{def:pmc} or~\ref{def:enclosed-weight} and under
Assumption~\ref{ass:eig} where it is given by Definition~\ref{def:pmcex}; and
the strengthening slot
$\{(\MCK,\pmc)\}$ satisfies the guarantee of Definition~\ref{def:slot}.
\end{digestclaim}

% claim receiver/def:slot
\digestlabel{def:slot}{8.16}\pdfbookmark[2]{Definition 8.16 The strengthening slot}{digest-receiver/def:slot}
\renewcommand*{\RestatedName}{\hyperref{contract.pdf}{}{claim-receiver/def:slot}{Definition~8.16}}%
\begin{digestclaim}[The strengthening slot]
The
\emph{strengthening slot}
is the finite list of pairs in which the artifact carries the high-photon
evidence,
\begin{equation*}\tag{\getrefnumber{C-eq:slot}}\digesteqlabel{eq:slot}
  \Big\{\big(\mathcal{O}_i,\ (p^{(t)}_i)_{t\in\Tset}\big)\Big\}_{i=1}^{r},
  \qquad \mathcal{O}_i\subseteq2^{[K]}\ \text{pairwise disjoint},\quad
  p^{(t)}_i\ge0,
\end{equation*}
one weight per entry and per setting $t\in\Tset$ of the setting family
(Definition~\ref{def:setting-family}), with the guarantee
$\sum_{S\in\mathcal{O}_i}\blk{M}{\gtN}_{S|t}\succeq p^{(t)}_i\,\id^{(>N)}$
for every setting $t\in\Tset$ and every $i$. The
\emph{uniform weight of an entry}
$i$ is $p_i\coloneqq\min_{t\in\Tset}p^{(t)}_i$; since
$0\le p_i\le p^{(t)}_i$, it carries the guarantee
$\sum_{S\in\mathcal{O}_i}\blk{M}{\gtN}_{S|t}\succeq p_i\,\id^{(>N)}$ for
every setting $t\in\Tset$ and every $i$, one weight at every setting. At a
one-setting family, $\Tset=\{\off\}$, an entry has the single weight
$p^{(\off)}_i=p_i$ and is written $(\mathcal{O}_i,p_i)$. The artifact
carries $r=1$ with
$(\mathcal{O}_1,(p^{(t)}_1)_{t\in\Tset})=(\MCK,(\pmc^{(t)})_{t\in\Tset})$,
the per-setting weights of Theorem~\ref{thm:emitted-bound}, each given at its
setting by Definition~\ref{def:pmc} or~\ref{def:pmcex} or as a non-negative
real at most the enclosed weight of Definition~\ref{def:enclosed-weight},
whose uniform weight is $\pmc$; at a
one-setting family that entry is $(\MCK,\pmc)$.
\end{digestclaim}

% claim receiver/prop:slotvalid
\digestlabel{prop:slotvalid}{8.17}\pdfbookmark[2]{Proposition 8.17 Range of the strengthening-slot weights}{digest-receiver/prop:slotvalid}
\begin{mainclaim}
\renewcommand*{\RestatedName}{\hyperref{contract.pdf}{}{claim-receiver/prop:slotvalid}{Proposition~8.17}}%
\begin{digestclaim}[Range of the strengthening-slot weights]
For any list \eqref{eq:slot} satisfying the stated guarantee with pairwise
disjoint $\mathcal{O}_i$, one has at every setting $t\in\Tset$ that
$p^{(t)}_i\ge0$ for every $i$ and $\sum_{i}p^{(t)}_i\le1$; hence the uniform
weights satisfy $p_i\ge0$ for every $i$ and $\sum_{i}p_i\le1$. In particular
the list this module emits --- by Corollary~\ref{cor:emitted-unconditioned}
and, for a setting family a submodule constructs, by that submodule's
instance of Theorem~\ref{thm:emitted-bound} --- satisfies
$\pmc^{(t)}\in[0,1]$ at every setting, so the failure conditions
``$p_y<0$'' and ``$\sum_yp_y>1$'' of a squashing request cannot be triggered
by this module's output; and what the further failure condition of a
squashing request, ``a bound does not hold at every declared setting it is
stated for and every photon number above the cutoff'', requires to hold --- each $\pmc^{(t)}$ at
its own setting $t$ and $\pmc$ at every setting, for every photon number
above the cutoff --- is Theorem~\ref{thm:emitted-bound}.
\end{digestclaim}
\end{mainclaim}

\section*{9\quad The reference, enclosed and exact}\pdfbookmark[1]{9 The reference, enclosed and exact}{digest-group-receiver/def:complex-enclosure}

% claim receiver/def:complex-enclosure
\digestlabel{def:complex-enclosure}{9.1}\pdfbookmark[2]{Definition 9.1 Complex enclosures, and enclosures of square roots and angles}{digest-receiver/def:complex-enclosure}
\renewcommand*{\RestatedName}{\hyperref{contract.pdf}{}{claim-receiver/def:complex-enclosure}{Definition~9.1}}%
\begin{digestclaim}[Complex enclosures, and enclosures of square roots and angles]
Write $\mathbb{Q}(i)=\{x+iy:x,y\in\mathbb{Q}\}$ and, for $z\in\mathbb{Q}(i)$,
$\lvert z\rvert_{1}\coloneqq\lvert\mathrm{Re}\,z\rvert+\lvert\mathrm{Im}\,z\rvert$,
a rational with $\lvert z\rvert\le\lvert z\rvert_{1}$. A
\emph{complex enclosure} of a complex number $x$ is a pair $(x_c,r)$ of
$x_c\in\mathbb{Q}(i)$, its \emph{centre}, and a rational $r\ge0$, its
\emph{radius}, with $\lvert x-x_c\rvert\le r$. It is \emph{exact} when $r=0$,
and then $x=x_c$. A complex enclosure of a complex matrix $X$ is a pair
$(X_c,r)$ of a matrix $X_c$ of the same shape with entries in
$\mathbb{Q}(i)$ and a matrix $r$ of non-negative rationals such that
$\big((X_c)_{il},r_{il}\big)$ is a complex enclosure of $X_{il}$ for all
$i,l$.

A \emph{square-root enclosure} of a rational $p\ge0$ is a complex enclosure
of $\sqrt p$ with a rational centre, exact where $p$ is the square of a
rational. An
\emph{angle enclosure} of a real number $\theta$ that is
given either as a rational, such as a declared value
(\xref{foundations/def:binary64}\,(iii)), or as $c\pi$ with $c$ a rational,
is a pair of complex enclosures with rational centres of $\cos\theta$ and of
$\sin\theta$, both exact where $\theta=c\pi$ with $2c$ an integer.
\end{digestclaim}

% claim receiver/thm:enclosed-blocks
\digestlabel{thm:enclosed-blocks}{9.5}\pdfbookmark[2]{Theorem 9.5 Rigorous enclosures of the photon-number blocks}{digest-receiver/thm:enclosed-blocks}
\renewcommand*{\RestatedName}{\hyperref{contract.pdf}{}{claim-receiver/thm:enclosed-blocks}{Theorem~9.5}}%
\begin{digestclaim}[Rigorous enclosures of the photon-number blocks]
Let $\Recv$ be a receiver instrument with
port Grams $\{A_k\}_{k=1}^{K}$ on
$\mathbb{C}^{\Inp}$ (Definitions~\ref{def:instrument} and~\ref{def:gram}),
$\{M_S\}_{S\subseteq[K]}$ its click-pattern POVM
(Theorem~\ref{thm:click-povm}), $\cg$ a declared coarse-graining onto
$\Vset$ and $\{F_v\}_{v\in\Vset}$ the emitted POVM
(Definition~\ref{def:coarse}). Fix $N\ge0$, and read every block
$\Symn{n}(\mathbb{C}^{\Inp})$, $n\le N$, in its occupation basis listed in an
order fixed for that block. Enclosures are
complex enclosures
(Definition~\ref{def:complex-enclosure}), and ``the rules'' are those of
Lemma~\ref{C-lem:enclosure-arithmetic}\,(i)--(iv).
\begin{enumerate}[label=(\roman*),leftmargin=3.0em]
\item \emph{Port Grams of a compiled receiver.} Suppose that the
  instrument matrix of $\Recv$ is the
  compiled instrument matrix $\Gmat=D_\eta\Gmat^{\sigma}_{\Grph}$ of a
  well-formed graph for a linear extension $\sigma$
  (Definition~\ref{def:compile}), and that enclosures are given of the matrix
  $\hat\Gmat_e$ of every element $e$ and of $\sqrt{\eta_k}$ for every port
  $k$. Enclose every other entry of $D_\eta$, and every entry of each
  $\Emd{\{k\}}$, exactly by itself. Then the rules, applied to the products
  \eqref{eq:compile} in the order of $\sigma$, to the restriction of
  $\Gmat$ to the columns $\Inp$, and to \eqref{eq:gramdef}, give enclosures
  of the $A_k$.
\item \emph{Blocks.} Let $(C_k,r_k)$ be enclosures of the $A_k$, and for
  every $n\le N$ and occupation vectors $\vec m,\vec m'$ of weight $n$ let
  $(a_{\vec m\vec m'},e_{\vec m\vec m'})$ be a square-root enclosure
  of $1/(\vec m!\,\vec m'!)$. For $J\subseteq[K]$ enclose $\id-A_J$ by the
  rules from the $(C_k,r_k)$, the identity being exact; enclose the entry
  $(\vec m,\vec m')$ of $\Pvac{J}^{(n)}$ by the rules applied to the product
  of $\operatorname{per}\big((\id-A_J)[\vec m|\vec m']\big)$ with
  $\sqrt{1/(\vec m!\,\vec m'!)}$; enclose $\blk{M}{\eqn}_S$ by the rules
  applied to \eqref{eq:blockMS}, the signs being exact; and enclose
  $\blk{F}{\eqn}_v$ by the rules applied to
  $\sum_{S\in\cg^{-1}(v)}\blk{M}{\eqn}_S$. Each enclosure so obtained is a
  complex enclosure of the block it names.
\item \emph{The reference.} Replace each enclosure of a block
  $\blk{F}{\eqn}_v$ of (ii) by the one with a Hermitian centre of
  Lemma~\ref{C-lem:enclosure-arithmetic}\,(v). Let $\mathfrak{m}_N$ be a
  sector map (\xref{foundations/def:sector-map}) whose blocks are the
  photon numbers $n=0,\dots,N$ of the declared inputs, in increasing order,
  of dimensions $\binom{|\Inp|+n-1}{n}$, with any block kind and block
  identifiers, each block read in its fixed order. Then the family
  $\{F_v\}$ restricted to the blocks $n\le N$, read with these enclosures, is
  a reference with constant $1$ against $\mathfrak{m}_N$ over $\Vset$
  (\xref{foundations/def:exact-reference}); the same holds for $\{M_S\}$ over
  the click patterns.
\item \emph{Exactness.} An entry of an enclosure of (i)--(iii) is exact where
  every enclosure it is computed from is exact; the reference of (iii) is an
  exact reference when every one of its radii is zero.
\end{enumerate}
\end{digestclaim}

% claim receiver/def:port-dilation
\digestlabel{def:port-dilation}{9.6}\pdfbookmark[2]{Definition 9.6 Port dilation of a receiver instrument}{digest-receiver/def:port-dilation}
\renewcommand*{\RestatedName}{\hyperref{contract.pdf}{}{claim-receiver/def:port-dilation}{Definition~9.6}}%
\begin{digestclaim}[Port dilation of a receiver instrument]
Let $\Recv$ be a receiver instrument with
port Grams $\{A_k\}_{k=1}^{K}$ on
$\mathbb{C}^{\Inp}$ (Definitions~\ref{def:instrument} and~\ref{def:gram}). A
\emph{port dilation} of $\Recv$ is a triple
$(\mathbb{O};\mathbb{O}_1,\dots,\mathbb{O}_K;\Omega)$ of a finite mode set
$\mathbb{O}$ with a fixed order, pairwise disjoint subsets
$\mathbb{O}_1,\dots,\mathbb{O}_K$ of $\mathbb{O}$, its \emph{ports}, and a
linear map $\Omega:\mathbb{C}^{\Inp}\to\mathbb{C}^{\mathbb{O}}$ with
\begin{equation*}\tag{\getrefnumber{C-eq:port-dilation}}\digesteqlabel{eq:port-dilation}
  \Omega^{\dagger}\Omega=\id_{\Inp}
  \qquad\text{and}\qquad
  \Omega^{\dagger}\Emd{\mathbb{O}_k}\Omega=A_k\quad(k=1,\dots,K),
\end{equation*}
where $\Emd{\mathbb{O}_k}$ is the coordinate projector of
$\mathbb{C}^{\mathbb{O}}$ onto the modes of $\mathbb{O}_k$. The modes of
$\mathbb{O}$ in no port are its \emph{loss modes}, and the \emph{row} of a mode
$i\in\mathbb{O}$ is the row of $\Omega$ that $i$ indexes. The \emph{click
pattern} of an occupation vector $\vec m\in\mathbb{N}_0^{\mathbb{O}}$
(Definition~\ref{def:fock}) is
\begin{equation*}\tag{\getrefnumber{C-eq:click-pattern-of}}\digesteqlabel{eq:click-pattern-of}
  S_{\vec m}\coloneqq\big\{k\in[K]:\ m_i\ge1\text{ for some }i\in\mathbb{O}_k\big\}.
\end{equation*}
\end{digestclaim}

% claim receiver/thm:factor-blocks
\digestlabel{thm:factor-blocks}{9.7}\pdfbookmark[2]{Theorem 9.7 Photon-number blocks in enclosed factor form}{digest-receiver/thm:factor-blocks}
\renewcommand*{\RestatedName}{\hyperref{contract.pdf}{}{claim-receiver/thm:factor-blocks}{Theorem~9.7}}%
\begin{digestclaim}[Photon-number blocks in enclosed factor form]
Let $\Recv$ be a receiver instrument with
port Grams $\{A_k\}_{k=1}^{K}$ on
$\mathbb{C}^{\Inp}$ and click-pattern POVM $\{M_S\}_{S\subseteq[K]}$
(Theorem~\ref{thm:click-povm}), and let
$(\mathbb{O};\mathbb{O}_1,\dots,\mathbb{O}_K;\Omega)$ be a
port dilation of $\Recv$ (Definition~\ref{def:port-dilation}), with
$S_{\vec m}$ the click pattern of an occupation vector $\vec m$ of
$\mathbb{O}$. Read every block $\Symn{n}(\mathbb{C}^{\Inp})$ in its
occupation basis listed in an order fixed for that block, and for an
occupation vector $\vec m\in\mathbb{N}_0^{\mathbb{O}}$ of weight $n$ put
\begin{equation*}\tag{\getrefnumber{C-eq:factor-column}}\digesteqlabel{eq:factor-column}
  c_{\vec m}\coloneqq\Fock(\Omega)^{\dagger}\ket{\vec m}
  \ \in\ \Symn{n}(\mathbb{C}^{\Inp}).
\end{equation*}
Enclosures are complex enclosures
(Definition~\ref{def:complex-enclosure}), and ``the rules'' are those of
Lemma~\ref{C-lem:enclosure-arithmetic}\,(i)--(iv).
\begin{enumerate}[label=(\roman*),leftmargin=3.0em]
\item \emph{The factor form.} For every $n\ge0$ and all reals $w_S\ge0$,
  $S\subseteq[K]$,
  \begin{equation*}\tag{\getrefnumber{C-eq:factor-form}}\digesteqlabel{eq:factor-form}
    \sum_{S\subseteq[K]}w_S\,\blk{M}{\eqn}_S
    \;=\;\sum_{\vec m:\,|\vec m|=n}w_{S_{\vec m}}\,c_{\vec m}c_{\vec m}^{\dagger},
  \end{equation*}
  the sum on the right over the occupation vectors of $\mathbb{O}$ of weight
  $n$. In particular $\blk{M}{\eqn}_S$ is the sum of the
  $c_{\vec m}c_{\vec m}^{\dagger}$ with $S_{\vec m}=S$, and the block
  $\blk{F}{\eqn}_v$ of the emitted POVM of a
  declared coarse-graining $\cg$ (Definition~\ref{def:coarse}) is
  \eqref{eq:factor-form} with $w_S=1$ for $S\in\cg^{-1}(v)$ and $w_S=0$
  otherwise. The entry of $c_{\vec m}$ at an occupation vector $\vec m'$ of
  $\Inp$ is
  $\overline{\operatorname{per}\Omega[\vec m|\vec m']}\big/\sqrt{\vec m!\,\vec m'!}$,
  with $\Omega[\vec m|\vec m']$ as in
  Lemma~\ref{C-lem:fock-permanent-maps}\,(ii).
\item \emph{Enclosed columns.} Let enclosures of the entries of $\Omega$ be
  given, and, for all occupation vectors $\vec m$ of $\mathbb{O}$ and
  $\vec m'$ of $\Inp$ of weight $n$, a square-root enclosure of
  $1/(\vec m!\,\vec m'!)$. The rules, applied to the conjugate of
  $\operatorname{per}\Omega[\vec m|\vec m']$ times that square root, enclose
  every entry of every $c_{\vec m}$. An entry so enclosed is exact where every
  enclosure it is computed from is exact, and it is enclosed exactly as $0$
  where every term of the permanent, a product of entries of $\Omega$, has a
  factor enclosed exactly as $0$; in particular every entry of $c_{\vec m}$
  is, where $\vec m$ has a photon in a mode whose row is enclosed exactly as
  zero.
\item \emph{The enclosure of the factor form.} Let $n\ge0$, let the weights $w_S\ge0$ be
  rationals, let enclosures of the columns $c_{\vec m}$ be given, such as
  those of (ii), and let $\mathcal{C}$ be the set of the occupation vectors
  $\vec m$ of $\mathbb{O}$ of weight $n$ with $w_{S_{\vec m}}>0$ whose column
  is not enclosed exactly as zero, listed in a fixed order. Let $(L_c,\varrho)$
  be the matrix box of complex
  $\dim\Symn{n}(\mathbb{C}^{\Inp})\times|\mathcal{C}|$ matrices whose column
  $\vec m$ has the centres and the radii of the enclosure of $c_{\vec m}$,
  and $W\coloneqq\operatorname{diag}(w_{S_{\vec m}})_{\vec m\in\mathcal{C}}$.
  Then $(L_c,\varrho,W)$ is a factor enclosure
  (\xref{foundations/def:factor-enclosure}) of the block
  \eqref{eq:factor-form}, with $|\mathcal{C}|$ columns; where $\mathcal{C}$
  is empty, that block is zero.
\item \emph{Port dilations from factors.} Let
  $\mathbb{O}_0,\mathbb{O}_1,\dots,\mathbb{O}_K$ be pairwise disjoint finite
  mode sets and $Y_k:\mathbb{C}^{\Inp}\to\mathbb{C}^{\mathbb{O}_k}$ linear
  maps with $Y_k^{\dagger}Y_k=A_k$ for $k=1,\dots,K$ and
  $Y_0^{\dagger}Y_0=\defic$, the loss deficiency
  (Definition~\ref{def:gram}). Then the map $\Omega$ into
  $\mathbb{O}\coloneqq\mathbb{O}_0\sqcup\mathbb{O}_1\sqcup\dots\sqcup\mathbb{O}_K$,
  in any fixed order, whose rows on $\mathbb{O}_k$ are those of $Y_k$ for
  $k=0,\dots,K$, is with the ports $\mathbb{O}_1,\dots,\mathbb{O}_K$ a port
  dilation of $\Recv$. If a matrix $X$ is $E^{\dagger}DE$ with $D$ diagonal with
  positive rational diagonal entries $d_i$ and $E$ with entries in
  $\mathbb{Q}(i)$, then $Y\coloneqq D^{1/2}E$ satisfies $Y^{\dagger}Y=X$, and
  the rules applied to the products of square-root enclosures of the $d_i$
  with the entries of $E$ enclose $Y$, exactly in every row $i$ in which
  $d_i$ is the square of a rational.
\end{enumerate}
\end{digestclaim}

% claim receiver/def:exact-route
\digestlabel{def:exact-route}{9.8}\pdfbookmark[2]{Definition 9.8 Gaussian-rational port Grams and the exact route}{digest-receiver/def:exact-route}
\renewcommand*{\RestatedName}{\hyperref{contract.pdf}{}{claim-receiver/def:exact-route}{Definition~9.8}}%
\begin{digestclaim}[Gaussian-rational port Grams and the exact route]
Let $\Recv$ be a receiver instrument with
port Grams $\{A_k\}_{k=1}^{K}$ on
$\mathbb{C}^{\Inp}$ (Definitions~\ref{def:instrument} and~\ref{def:gram}),
and write $\mathbb{Q}(i)=\{x+iy:x,y\in\mathbb{Q}\}$. A declaration of a
receiver is on the
\emph{exact route} when (i) the port Grams $\{A_k\}$ that its
declared data determine, computed in exact arithmetic, have entries in
$\mathbb{Q}(i)$, are Hermitian, and are an admissible Gram family
exactly (Definition~\ref{def:gram}), that is,
\begin{equation*}\tag{\getrefnumber{C-eq:exact-admissible}}\digesteqlabel{eq:exact-admissible}
  A_k\succeq0\quad(k=1,\dots,K)
  \qquad\text{and}\qquad
  \id-A_{[K]}\succeq0 ,
\end{equation*}
conditions on Gaussian-rational matrices; and (ii) its
probe declaration, where it has one, is of the Gaussian tier
(Definition~\ref{def:probe}), whose mean photon numbers $\bar n^{(t)}_m$ are
declared values and so rationals (\xref{foundations/def:binary64}\,(iii)).
Under (i) the port Grams are those of a receiver instrument
(Theorem~\ref{thm:representation-independence}\,(ii)), and every receiver
instrument with these port Grams induces the same click-pattern POVM
(Theorem~\ref{thm:representation-independence}\,(i)); where the
instrument matrix $\Gmat$, declared or compiled, satisfies
$\Gmat^\dagger\Gmat\preceq\id$ exactly, the declared instrument is one of
them. A declared family admitted only at the receiver tolerances
(Definition~\ref{def:tolmap}) and not admissible exactly is not on the exact
route. A port Gram family or an instrument matrix declared directly has
entries that are dyadic rationals (\xref{foundations/def:binary64}\,(iii));
its port Grams then have entries in $\mathbb{Q}(i)$, by \eqref{eq:gramdef}
for a declared instrument matrix, and such a declaration satisfies (i)
exactly when its port Grams are Hermitian, as they always are for a declared
instrument matrix, and admissible exactly.
\end{digestclaim}

% claim receiver/thm:exact-blocks
\digestlabel{thm:exact-blocks}{9.9}\pdfbookmark[2]{Theorem 9.9 Exact entries of the blocks on the exact route}{digest-receiver/thm:exact-blocks}
\renewcommand*{\RestatedName}{\hyperref{contract.pdf}{}{claim-receiver/thm:exact-blocks}{Theorem~9.9}}%
\begin{digestclaim}[Exact entries of the blocks on the exact route]
Let $\Recv$ be a receiver instrument whose
port Grams $\{A_k\}$ have entries in
$\mathbb{Q}(i)$, as on the exact route
(Definition~\ref{def:exact-route}), with click-pattern POVM $\{M_S\}$, a
declared coarse-graining $\cg$ and emitted POVM $\{F_v\}$. Let $n\ge0$
and let $\vec m,\vec m'$ be occupation vectors of weight $n$ such that
$\vec m!\,\vec m'!$ is the square of an integer; this holds on the diagonal,
$\vec m=\vec m'$, and for every $\vec m,\vec m'$ when $n\le1$.
\begin{enumerate}[label=(\roman*),leftmargin=3.0em]
\item The entry $(\vec m,\vec m')$ of every block $\Pvac{J}^{(n)}$,
  $\blk{M}{\eqn}_S$ and $\blk{F}{\eqn}_v$ lies in $\mathbb{Q}(i)$.
\item If the enclosures of the $A_k$ in Theorem~\ref{thm:enclosed-blocks}\,(ii)
  are exact, that is, each $A_k$ is enclosed by itself with radius zero,
  then the enclosure of the entry $(\vec m,\vec m')$ of each of these blocks
  in Theorem~\ref{thm:enclosed-blocks}\,(ii), and of the blocks
  $\blk{F}{\eqn}_v$ in Theorem~\ref{thm:enclosed-blocks}\,(iii), is exact,
  its centre being the entry.
\item Under the hypothesis of (ii), every block $n\le1$ of the reference of
  Theorem~\ref{thm:enclosed-blocks}\,(iii) has radius zero everywhere, and for
  a cutoff $N\le1$ that reference is an exact reference
  (\xref{foundations/def:exact-reference}).
\end{enumerate}
\end{digestclaim}

\section*{10\quad Residuals, tolerances and floating-point evaluation}\pdfbookmark[1]{10 Residuals, tolerances and floating-point evaluation}{digest-group-receiver/def:residuals}

% claim receiver/def:residuals
\digestlabel{def:residuals}{10.1}\pdfbookmark[2]{Definition 10.1 The six residuals}{digest-receiver/def:residuals}
\renewcommand*{\RestatedName}{\hyperref{contract.pdf}{}{claim-receiver/def:residuals}{Definition~10.1}}%
\begin{digestclaim}[The six residuals]
The \emph{residuals} of an evaluation of this contract are those of the six quantities
below that exist on the input route taken and for the evaluation requested
(which exist is stated after \eqref{eq:rhodef}):
\begin{align*}
  r_1&\coloneqq\big\lVert\big(\Gmat^\dagger\Gmat-\id\big)_{+}\big\rVert
    &&\text{(contractivity defect: the positive part must be }0),
    \tag{\getrefnumber{C-eq:res1}}\digesteqlabel{eq:res1}\\
  r_2&\coloneqq\max_k\ \max\big\{0,\,-\lmin(A_k)\big\}
    &&\text{(Gram positivity)},\tag{\getrefnumber{C-eq:res2}}\digesteqlabel{eq:res2}\\
  r_3&\coloneqq\big\lVert\,\id_{\Inp}-A_\Sigma
     -\Gmat_{\Inp}^\dagger\Emd{\mathrm{un}}\Gmat_{\Inp}
     -\big(\id_{\Inp}-\Gmat_{\Inp}^\dagger\Gmat_{\Inp}\big)\big\rVert
    &&\text{(deficiency identity \eqref{eq:deficiencyid})},\tag{\getrefnumber{C-eq:res3}}\digesteqlabel{eq:res3}\\
  r_4&\coloneqq\Big\lVert\,\id-\sum_{S\subseteq[K]}
     \blk{M}{\eqn}_{S|t}\Big\rVert\ \ \text{for each emitted }n
     \text{ and each setting }t
    &&\text{(POVM completeness, blockwise)},\tag{\getrefnumber{C-eq:res4}}\digesteqlabel{eq:res4}\\
  r_5&\coloneqq\max_{v,\,n,\,t}\ \varrho\big(\blk{M}{\eqn}_{v|t}\big)
    &&\text{(block Hermiticity and positivity)},\tag{\getrefnumber{C-eq:res5}}\digesteqlabel{eq:res5}\\
  r_6&\coloneqq\max_{t,\ N<n\le m^{\ast}_t}
     \big(\Bnd^{(t)}(n)-\widehat\Bnd^{\ast(t)}(n)\big)_{+}
    &&\text{(window consistency, Lem.~\ref{lem:exactbasic})},\tag{\getrefnumber{C-eq:res6}}\digesteqlabel{eq:res6}
\end{align*}
in which the \emph{Hermiticity-and-positivity defect} of a square matrix $X$ is
\begin{equation*}\tag{\getrefnumber{C-eq:rhodef}}\digesteqlabel{eq:rhodef}
  \varrho(X)\ \coloneqq\ \max\Big\{\,\big\lVert X-X^\dagger\big\rVert,\ \
    \big(-\lmin\big(\tfrac12(X+X^\dagger)\big)\big)_{+}\,\Big\} ,
\end{equation*}
where $\Emd{\mathrm{un}}$ projects onto the declared modes in no port,
$(\cdot)_+$ is the positive part, and $\lVert\cdot\rVert$ is the spectral norm
throughout \eqref{eq:res1}--\eqref{eq:rhodef}. The maximum in \eqref{eq:res5}
runs over exactly the blocks $\blk{M}{\eqn}_{v|t}$ of the coarse-grained
elements, before any dark-count channel, that the artifact emits: every
emitted outcome letter
$v\in\Vset$, every photon number $0\le n\le N$, and every setting $t\in\Tset$
of the setting family (Definition~\ref{def:setting-family}; one setting when
$\InpS=\varnothing$, Corollary~\ref{cor:unconditioned-family}); the emitted
bound of the one-setting family is Corollary~\ref{cor:emitted-unconditioned},
and that of a setting family a submodule constructs is that submodule's
instance of Theorem~\ref{thm:emitted-bound}. Here $\blk{M}{\eqn}_{v|t}$ is the
emitted block: the $n$-photon block of the coarse-grained conditional element
$F_{v|t}$, written $\blk{F}{\eqn}_v$ in Theorem~\ref{thm:coarse-graining}\,(i) at a single
setting.

\emph{Which residuals exist on which route.} $r_2$, $r_4$ and $r_5$ are
functions of the port Grams and of the emitted blocks, which every route
produces, so all three exist on all three routes. $r_1$ \eqref{eq:res1} and
$r_3$ \eqref{eq:res3} are written on $\Gmat$ and $\Gmat_{\Inp}$, and route (i)
supplies no instrument matrix: there the Grams are the declared datum and no
$\Gmat$ exists, so $r_1$ and $r_3$ are \emph{not defined}. Synthesizing a
$\Gmat$ from Theorem~\ref{thm:representation-independence}\,(ii) to evaluate them on would test the
synthesis rather than the declaration. This agrees with
Theorem~\ref{thm:exact-and-evaluated}\,(iv), where route (i) is the case in which ``there is
no compilation''. On routes (ii) and (iii) all five of $r_1,\dots,r_5$ exist.

\emph{$r_6$ depends on the request, not on the route.} It compares the two
evaluations of a single proved inequality, $\Bnd^{(t)}(n)\le
\Bnd^{\ast(t)}(n)$ (Lemma~\ref{lem:exactbasic}), so its exact value is $0$
like the other five, and it exists on all three routes but \emph{only at the
settings whose declaration carries an exact-window depth}
(Definition~\ref{def:pmcex}), over exactly the indices $N<n\le m^{\ast}_t$ that
depth names. A declaration that requests no exact window computes no
$\widehat\Bnd^{\ast}$ and has nothing to compare, so $r_6$ is not defined
there, as $r_1$ and $r_3$ are not on route (i).

No residual bounds the distance between a computed quantity and the exact
one (Theorem~\ref{thm:exact-and-evaluated}\,(ii)).
\end{digestclaim}

% claim receiver/prop:res5
\digestlabel{prop:res5}{10.2}\pdfbookmark[2]{Proposition 10.2 Properties of the Hermiticity-and-positivity defect}{digest-receiver/prop:res5}
\renewcommand*{\RestatedName}{\hyperref{contract.pdf}{}{claim-receiver/prop:res5}{Proposition~10.2}}%
\begin{digestclaim}[Properties of the Hermiticity-and-positivity defect]
Let $X$ be a complex $D\times D$ matrix and $\varrho(X)$ its
Hermiticity-and-positivity defect \eqref{eq:rhodef}.
\begin{enumerate}[label=(\roman*),leftmargin=3.0em]
\item $\varrho(X)=0$ if and only if $X=X^\dagger$ and $X\succeq0$.
\item If $X=X^\dagger$ then $\varrho(X)=\big(-\lmin(X)\big)_+$, which is the
  form \eqref{eq:res2} takes on a Gram.
\item For every unit vector $\ket{\psi}$,
  $\Re\bra{\psi}X\ket{\psi}\ge-\varrho(X)$ and
  $\big|\Im\bra{\psi}X\ket{\psi}\big|\le\tfrac12\varrho(X)$. Hence
  $\varrho(X)\le\varepsilon$ bounds every diagonal expectation of $X$ away from
  the negative reals by $\varepsilon$ and off the real axis by
  $\varepsilon/2$.
\item Every block $\blk{M}{\eqn}_{v|t}$ of the coarse-grained elements, before
  any dark-count channel, that the artifact emits satisfies $\varrho=0$
  \emph{exactly}: for
  each $v$, $n$ and $t$, $\blk{M}{\eqn}_{v|t}$ is a diagonal block of the
  positive operator $F_{v|t}$ (Definition~\ref{def:setting-family}\,(a) and
  Theorem~\ref{thm:coarse-graining}), and a diagonal block --- the compression $\Pi X\Pi$ of a positive
  operator $X$ by an orthogonal projector $\Pi$, read on the range of $\Pi$ ---
  is Hermitian and positive semidefinite. So $r_5$ is a residual in exactly the
  sense of the other five: the exact value is $0$, and what is emitted is the
  floating-point evaluation.
\end{enumerate}
\end{digestclaim}

% claim receiver/def:tolmap
\digestlabel{def:tolmap}{10.4}\pdfbookmark[2]{Definition 10.4 Receiver tolerances and the exact conditions they decide}{digest-receiver/def:tolmap}
\renewcommand*{\RestatedName}{\hyperref{contract.pdf}{}{claim-receiver/def:tolmap}{Definition~10.4}}%
\begin{digestclaim}[Receiver tolerances and the exact conditions they decide]
Several conditions this contract states are \emph{exact} statements about
declared data that this contract nevertheless decides in floating point. Each is decided by a declared number, fixed here,
and by no other threshold. The \emph{receiver tolerances} are the declared
$\varepsilon_1,\dots,\varepsilon_6$, $\varepsilon_{\mathrm{gap}}$ and
$\varepsilon_{\mathrm{eig}}$, each in the declared domain $[0,1)$; every quantitative statement below uses
that upper bound and says where. Items (i) to (v) spend $\varepsilon_1$,
$\varepsilon_2$ and $\varepsilon_4$; $\varepsilon_3$, $\varepsilon_5$,
$\varepsilon_6$ and $\varepsilon_{\mathrm{gap}}$ decide no exact condition on
declared data and bound computed quantities only (kind (c) below); and
$\varepsilon_{\mathrm{eig}}$ decides nothing at all. It is subtracted from a
computed eigenvalue in \eqref{eq:elldef}, so it is the one declared number this
contract spends on a \emph{hypothesis} (Assumption~\ref{ass:eig}) rather than
on a test.
\begin{enumerate}[label=(\roman*),leftmargin=3.0em]
\item \emph{Route-(i) Gram admissibility}, $A_k\succeq0$ and
  $\sum_kA_k\preceq\id$ (Definition~\ref{def:gram}), is decided at
  $\varepsilon_2$: the family is admitted when $\lmin(A_k)\ge-\varepsilon_2$
  for every $k$ and $\lmax(A_\Sigma)\le1+\varepsilon_2$. $\varepsilon_2$ is the
  tolerance the contract already attaches to Gram positivity through
  \eqref{eq:res2}, so the same declared number bounds the residual and decides
  the admission, in the same units.
\item \emph{The route-(i) rank consistency check} $\rk A_k\le|\Port{k}|$ is
  decided by the $\varepsilon_2$-rank
  \begin{equation*}\tag{\getrefnumber{C-eq:eps2rank}}\digesteqlabel{eq:eps2rank}
    \rk_{\varepsilon_2}(X)\ \coloneqq\
    \#\big\{\,i:\ \lambda_i\big(\tfrac12(X+X^\dagger)\big)>\varepsilon_2\,\big\}
    \qquad\text{(eigenvalues with multiplicity)},
  \end{equation*}
  the check being $\rk_{\varepsilon_2}(A_k)\le|\Port{k}|$ for every $k$. The
  threshold is \emph{absolute} and it is $\varepsilon_2$; no separate tolerance
  is added.
\item \emph{Dark-channel admissibility} for a \emph{declared} $R$ is decided
  in the detector-imperfections submodule, Definition~\ref{def:darktol}, which
  extends this list; the tolerances that definition spends are the
  $\varepsilon_4$ and $\varepsilon_2$ of this list.
\item \emph{Bound-data admissibility}, $0\le c$ in \eqref{eq:bounddata}, is
  decided at $\varepsilon_2$: the data are admitted when
  $c=1-\lmax(B)\ge-\varepsilon_2$. \emph{Nothing is clamped}: the computed
  $c$ enters \eqref{eq:beta} exactly as it stands. The remaining inequalities of
  \eqref{eq:bounddata}, $\gamma\le\beta_k\le1$ and $c\le d_k\le1$, are
  \emph{not} separately tested and no threshold decides them: they are
  redundant by theorem, discharged by the admissibility paragraph of
  Theorem~\ref{thm:multiclick-bound} from \eqref{eq:bbkhyp} and from
  $\gamma\le\beta_k$, both hypotheses of the setting family
  (Definition~\ref{def:setting-family}\,(b)), and their floating-point
  evaluation is governed by Theorem~\ref{thm:exact-and-evaluated}\,(ii) like every other
  computed quantity.

  \emph{Item (iv) decides this admission alone, on every input route, and
  item~(v) does not decide it.} Whenever $\varepsilon_1>\varepsilon_2$ the band
  \begin{equation*}\tag{\getrefnumber{C-eq:tolmapband}}\digesteqlabel{eq:tolmapband}
    c\in\big[-\varepsilon_1,\,-\varepsilon_2\big)
  \end{equation*}
  is non-empty, admitted by item~(v)'s test and refused by item (iv)'s, and
  item (iv) wins on it: the band is refused.
\item \emph{Route-(ii) contractivity, and (E1) at catalogue entry C10.}
  $\Gmat^\dagger\Gmat\preceq\id$ (Definition~\ref{def:instrument}) is decided
  at $\varepsilon_1$: the matrix is admitted when
  $\lmax(\Gmat^\dagger\Gmat)\le1+\varepsilon_1$. The same test decides the (E1)
  predicate of the catalogue's generic block entry C10
  (Definition~\ref{def:c10-tolerance} of the component-catalogue submodule).

  The tolerance that bounds the residual $r_1$ also decides this admission.
  $\Gmat^\dagger\Gmat$ is Hermitian positive semidefinite, so
  $\big(\Gmat^\dagger\Gmat-\id\big)_+$ has eigenvalues
  $\big(\lambda_i(\Gmat^\dagger\Gmat)-1\big)_+$ and therefore
  \begin{equation*}\tag{\getrefnumber{C-eq:r1islmax}}\digesteqlabel{eq:r1islmax}
    r_1=\big\lVert\big(\Gmat^\dagger\Gmat-\id\big)_{+}\big\rVert
      =\big(\lmax(\Gmat^\dagger\Gmat)-1\big)_{+} ,
  \end{equation*}
  so ``$r_1\le\varepsilon_1$'' and ``$\lmax(\Gmat^\dagger\Gmat)\le
  1+\varepsilon_1$'' are the \emph{same test on the same number}, and no
  separate tolerance is added.
\end{enumerate}
\emph{The five items, together with the two extensions the submodules declare
in the same form, Definition~\ref{def:darktol}, which is item (iii)'s content,
and Definition~\ref{def:c10-tolerance}, are the complete set of exact conditions
about \emph{declared} data that this module and its submodules decide under a
threshold.} Every other validation they state is one of three kinds, and none
of them adds a threshold this definition has not named:
\begin{enumerate}[label=(\alph*),leftmargin=2.2em]
\item a comparison of declared data against a declared domain: every parameter
  domain of Table~\ref{C-tab:catalogue} except C10's, the detector-layer ranges,
  $m_t>N$, $m^{\ast}_t>N$, $m^{\ast}_t\le m_t$, $N\ge0$, the domain of the
  receiver tolerances itself, the test of
  the reserved tolerance field of Assumption~\ref{ass:exact}, and the diagonal
  non-negativity of $\Gamma^{(t)}$;
\item a combinatorial predicate on the declaration decided by finite
  inspection: (E2) and (E3) of Definition~\ref{def:graph}, route exclusivity,
  the port map's disjointness, totality and surjectivity of $\cg$, the
  single-valued $\tau$ of Definition~\ref{def:modes}, the exclusivity of the
  two dark-count routes of Definition~\ref{def:darkadm}, predicate~(X) of
  Definition~\ref{def:pmc}, and the conditional presence of
  $\varepsilon_{\mathrm{eig}}$ and $\varepsilon_6$ at a declaration carrying an
  exact-window depth;
\item the residual tests themselves: $r_3$, $r_4$, $r_5$, $r_6$ against
  $\varepsilon_3$, $\varepsilon_4$, $\varepsilon_5$, $\varepsilon_6$, and the
  residual gap $g_{N,m_t}$ against $\varepsilon_{\mathrm{gap}}$. These are
  threshold decisions, but on \emph{computed} quantities rather than on
  declared data, and each is a declared-threshold decision, for the residuals of
  Definition~\ref{def:residuals} and for the residual gap \eqref{eq:resgap},
  so nothing undeclared decides them either. The tests $r_1\le\varepsilon_1$
  and $r_2\le\varepsilon_2$ are of this kind too, except where they coincide
  with an item: on route (ii), where $\Gmat$ is declared, the test
  $r_1\le\varepsilon_1$ \emph{is} item (v) by \eqref{eq:r1islmax}, and on
  route (i), where the Grams are declared, $r_2\le\varepsilon_2$ is the
  positivity half of item (i) written as a residual; on route (iii) $r_1$
  tests the compiled $\Gmat$, and on routes (ii) and (iii) $r_2$ tests
  computed Grams. Neither is
  $\varepsilon_{\mathrm{eig}}$, which is not a test of any kind: it is
  subtracted from $\widehat\Bnd^{\ast}$ in \eqref{eq:elldef} and changes an
  emitted value, not an admission.
\end{enumerate}
\end{digestclaim}

% claim receiver/def:evaluation-model
\digestlabel{def:evaluation-model}{10.9}\pdfbookmark[2]{Definition 10.9 The evaluation model: floating point, enclosures and rigour}{digest-receiver/def:evaluation-model}
\renewcommand*{\RestatedName}{\hyperref{contract.pdf}{}{claim-receiver/def:evaluation-model}{Definition~10.9}}%
\begin{digestclaim}[The evaluation model: floating point, enclosures and rigour]
The claims of this contract are exact statements about exactly declared data
(\xref{foundations/def:binary64}\,(iii)). This contract \emph{evaluates} them
in binary64 floating-point arithmetic (\xref{foundations/def:binary64}),
reports those of the six residuals of Definition~\ref{def:residuals}
that exist on the input route taken and for the evaluation requested, and
asserts no rigorous enclosure of any quantity it computes in floating point
(Assumption~\ref{ass:eig} is the one exception): this is its
\emph{floating-point evaluation}. For a declaration whose
model is an admissible receiver exactly it moreover computes, in exact
rational arithmetic, complex enclosures
(Definition~\ref{def:complex-enclosure}) of the photon-number blocks of the
emitted family: those of Theorem~\ref{thm:enclosed-blocks}, and in the
submodules those of the setting families a submodule constructs, each
enclosed by that submodule, and of their composition with the dark counts
(Corollary~\ref{cor:exact-dark-composed}). They are exact where
the declaration is on the exact route
(Definition~\ref{def:exact-route}) and the entry is known exactly
(Theorem~\ref{thm:exact-blocks}). From the same enclosures of the port Grams
it computes, in exact rational arithmetic, at every setting at which the rules
of Lemma~\ref{C-lem:enclosure-arithmetic} give them, complex enclosures of the
setting's operators $B^{(t)}$, $B^{(t)}_k$ and prefactors $\gamma^{(t)}$,
$\beta^{(t)}_k$ (Definition~\ref{def:setting-family}) and, at a setting
carrying an exact-window depth, of its blocks $\blk{M}{\eqn}_{\mc|t}$ in the
window; from these, an upper bound $u$ on $\lmax(B^{(t)})$ and lower bounds
$\ell_k$ on $\lmin(B^{(t)}_k)$ and $\lambda_n$ on the least eigenvalue of
each window block, by the least-eigenvalue bounds over an enclosure
(\xref{foundations/lem:least-eigenvalue-bound},
\xref{foundations/cor:least-eigenvalue-over-enclosure}), and a lower bound $g$
on $\gamma^{(t)}$ and upper bounds $b_k$ on $\beta^{(t)}_k$, each the real
part of a centre less or plus its radius; from these, the reals
$g,b_k,1-u,1-\ell_k$ each intersected with $[0,1]$ and
$\max\{0,\lambda_n\}$; and from those, by the formulas of
Definition~\ref{def:enclosed-weight}\,(iii), a weight: this is its
\emph{rigorous evaluation} of the weights. A weight emitted at a setting is
the weight so formed where the rigorous evaluation is formed there, and
otherwise the floating-point evaluation of Definition~\ref{def:pmc}
or~\ref{def:pmcex}.
These enclosures and bounds are not evaluations, and a decision taken on
them is exact; their interval arithmetic is proved in this contract
(Lemma~\ref{C-lem:enclosure-arithmetic}) and their eigenvalue bounds in the
foundations contract, and no interval-arithmetic assumption is imported. This
\emph{evaluation model}
is a \emph{restriction} this contract places on itself and asserts nothing a
run could violate.
\end{digestclaim}

% claim receiver/thm:exact-and-evaluated
\digestlabel{thm:exact-and-evaluated}{10.10}\pdfbookmark[2]{Theorem 10.10 Exact statements, floating-point and rigorous evaluation}{digest-receiver/thm:exact-and-evaluated}
\renewcommand*{\RestatedName}{\hyperref{contract.pdf}{}{claim-receiver/thm:exact-and-evaluated}{Theorem~10.10}}%
\begin{digestclaim}[Exact statements, floating-point and rigorous evaluation]
For the evaluation model of Definition~\ref{def:evaluation-model}:
\begin{enumerate}[label=(\roman*),leftmargin=3.0em]
\item \emph{The asserted mathematics is exact, except where a claim says so in
  its own statement.} \emph{Every asserted identity and inequality} of this
  module and its submodules is between exactly specified quantities and
  contains no error term, no tolerance and no numerical hypothesis
  (Theorem~\ref{thm:multiclick-bound}, Lemma~\ref{lem:collapse},
  Lemma~\ref{lem:windowtail}, Lemma~\ref{lem:exactbasic} and
  Lemma~\ref{lem:exactwindow} included, the last two being stated on the exact
  $\Bnd^{\ast}$ and carrying no tolerance), \emph{except where a claim's own
  statement names a receiver tolerance or
  Assumption~\ref{ass:eig}}.
\item \emph{The floating-point evaluation is not.} The emitted $\{A_k\}$,
  block matrices and those bound values that the floating-point evaluation
  produces are floating-point evaluations of those quantities; the enclosures
  of item (vi) and the rigorous evaluation of item (vii) are not. This
  contract asserts no bound on the difference, and the residuals of
  Definition~\ref{def:residuals} are \emph{diagnostics}: passing them is
  necessary for the artifact to be self-consistent and is not sufficient for
  any enclosure statement.
\item \emph{The evaluation label.} An artifact carries the rigorous evaluation
  label exactly where every per-setting weight it emits is formed by the
  rigorous evaluation of Definition~\ref{def:evaluation-model}, as the
  enclosed weight of Definition~\ref{def:enclosed-weight} or a
  non-negative real below it; every other artifact, on every input route, the
  route-(iii) compiler included, carries the floating-point evaluation label.
  The rigorous label concerns the weights alone: the blocks are evaluated in
  floating point under (ii) whatever the label. No path in this module
  produces a label stronger than these two.
\item \emph{Where route (i) stands.} If the Grams are \emph{declared}, as
  binary64 values, which are dyadic rationals
  (\xref{foundations/def:binary64}\,(iii)), the compilation error of (ii) does
  not exist: there is no compilation. The eigenvalue quantities of
  Theorem~\ref{thm:multiclick-bound} are still evaluated in floating point by
  the floating-point evaluation, so (ii) and (iii) apply to a weight it forms;
  what changes is that the rigorous evaluation of (vii) then encloses no
  matrix product, the enclosures it starts from being the declared Grams
  themselves. This is also why $r_1$ and $r_3$ do not exist on this route: they
  are written on a $\Gmat$ that is not supplied
  (Definition~\ref{def:residuals}). Nothing further is claimed.
\item \emph{The one place a floating-point number is asserted to bound an
  exact one, and it is opt-in.} Where a setting $t$ carries an exact-window
  depth and its emitted weight $\pmc^{(t)}$ is that of
  Definition~\ref{def:pmcex}, the weight depends on the
  computed $\widehat\Bnd^{\ast}$ through \eqref{eq:elldef}, and through the
  minimum so may the uniform $\pmc$; the inequality
  $\widehat\Bnd^{\ast(t)}(n)-\varepsilon_{\mathrm{eig}}\le\Bnd^{\ast(t)}(n)$
  that makes each such emitted value a lower bound is Assumption~\ref{ass:eig},
  not a theorem of this contract: (ii) governs $\widehat\Bnd^{\ast}$ exactly as
  it governs every other computed quantity, and $r_6$ is a diagnostic in
  exactly the sense (ii) fixes and is blind to the direction the assumption
  constrains. Where a setting's emitted weight is an enclosed weight, no
  computed $\widehat\Bnd^{\ast}$ and no $\varepsilon_{\mathrm{eig}}$ enters
  it and Assumption~\ref{ass:eig} is not invoked there, whatever depth is
  declared. Where no depth is declared, no
  $\Bnd^{\ast}$ is computed, Assumption~\ref{ass:eig} is not invoked, and the
  emitted weights are the values \eqref{eq:pmcdef} defines, expression for
  expression, evaluated in floating point or on enclosed bound data. Item
  (iii) covers all of these: an artifact with a weight resting on
  Assumption~\ref{ass:eig} carries the floating-point label, so no artifact
  resting on Assumption~\ref{ass:eig} can claim more.
\item \emph{The model's blocks are enclosed, and the evaluation approximates
  them.} For a declaration whose model is an admissible receiver exactly, the
  photon-number blocks of the click-pattern and
  emitted POVMs of that model are quantities the
  identities and inequalities of (i) are about. Their
  complex enclosures
  (Definition~\ref{def:complex-enclosure}) of
  Theorem~\ref{thm:enclosed-blocks}, where they are formed, are computed in
  exact rational arithmetic and contain them: every entry of a
  block lies within its radius of its centre, an entry whose radius is zero
  is its centre, and on the exact route the entries that
  Theorem~\ref{thm:exact-blocks} names have radius zero when the port Grams
  are enclosed exactly. A positive semidefiniteness or a positive
  definiteness established through
  \xref{foundations/lem:enclosure-psd}\,(i) on such an enclosure with a
  Hermitian centre holds of the block itself. The floating-point block
  matrices are evaluations of the same blocks and fall under (ii): this
  contract asserts no bound on their difference from the blocks, so no
  decision about the blocks is taken on them.
\item \emph{The rigorously evaluated weights are bounds.} Let $t$ be a
  setting of a setting family (Definition~\ref{def:setting-family}) at
  which the rigorous evaluation is formed from complex enclosures of
  $B^{(t)}$, of every $B^{(t)}_k$, of $\gamma^{(t)}$, of every
  $\beta^{(t)}_k$ and, at a setting carrying a depth, of the window blocks
  $\blk{M}{\eqn}_{\mc|t}$; for a receiver without a probe, whose setting
  family is that of Corollary~\ref{cor:unconditioned-family}, the enclosures
  it forms from the enclosures of the port Grams are such. Then the four
  clipped reals it forms
  are enclosed bound data of the setting and, at a setting carrying a
  depth, the window values it forms are rigorous window bounds; the
  weight it forms from them is their enclosed weight, a rational computed
  exactly; and it and every real between $0$
  and it bound the multiclick element of the setting from below on every
  photon-number block above the cutoff and lie in $[0,1]$. Its validity rests
  on no number computed in floating point: such a number enters it, if at
  all, only as the exact rational it is, of which nothing is assumed, as the
  approximate eigenbasis of \xref{foundations/lem:least-eigenvalue-bound}
  does. So (ii) does not govern it, and it rests on no assumption about
  arithmetic.
\end{enumerate}
\end{digestclaim}

\section*{11\quad Assumptions}\pdfbookmark[1]{11 Assumptions}{digest-group-receiver/ass:trust}

% claim receiver/ass:trust
\digestlabel{ass:trust}{11.1}\pdfbookmark[2]{Assumption 11.1 Trusted and exactly characterised receiver}{digest-receiver/ass:trust}
\SetShortId{A-TRUST}
\renewcommand*{\RestatedName}{\hyperref{contract.pdf}{}{claim-receiver/ass:trust}{Assumption~11.1}}%
\begin{digestassumption}[Trusted and exactly characterised receiver]
(i) The declared instrument describes the physical receiver \emph{exactly}:
its mode space, its declared element graph or the supplied
instrument matrix or port Gram family, the
port map, the efficiencies, the dark-count probabilities and the probe states;
(ii) the receiver is under the legitimate party's control: it is not
selected, tuned or otherwise influenced by the adversary.
\end{digestassumption}

% claim receiver/ass:passive
\digestlabel{ass:passive}{11.3}\pdfbookmark[2]{Assumption 11.3 Passive linear optics}{digest-receiver/ass:passive}
\SetShortId{A-PASSIVE}
\renewcommand*{\RestatedName}{\hyperref{contract.pdf}{}{claim-receiver/ass:passive}{Assumption~11.3}}%
\begin{digestassumption}[Passive linear optics]
(i) Every optical element is passive linear optics: it maps annihilation
operators to combinations of annihilation operators,
$b_i=\sum_j\Gmat_{ij}a_j$, with $\Gmat^\dagger\Gmat\preceq\id$;
(ii) no element mixes $a$ with $a^\dagger$, and none provides gain.
\end{digestassumption}

% claim receiver/ass:vac
\digestlabel{ass:vac}{11.4}\pdfbookmark[2]{Assumption 11.4 Vacuum on undeclared input modes}{digest-receiver/ass:vac}
\SetShortId{A-VAC}
\renewcommand*{\RestatedName}{\hyperref{contract.pdf}{}{claim-receiver/ass:vac}{Assumption~11.4}}%
\begin{digestassumption}[Vacuum on undeclared input modes]
(i) Every mode of $\Modes\setminus\Inp$, and every mode of the dilation
$\Env$, carries the vacuum state at the input;
(ii) every mode that can carry a non-vacuum state and can couple to a port is
declared in $\Inp$: in $\InpC$ if the adversary can populate it, in $\InpS$
if the legitimate party does.
\end{digestassumption}

% claim receiver/ass:modes
\digestlabel{ass:modes}{11.5}\pdfbookmark[2]{Assumption 11.5 Completeness of the declared modes}{digest-receiver/ass:modes}
\SetShortId{A-MODES}
\renewcommand*{\RestatedName}{\hyperref{contract.pdf}{}{claim-receiver/ass:modes}{Assumption~11.5}}%
\begin{digestassumption}[Completeness of the declared modes]
(i) Every optical path from outside the receiver to a detector is represented
by a declared mode of $\Modes$.
\end{digestassumption}

% claim receiver/ass:thresh
\digestlabel{ass:thresh}{11.6}\pdfbookmark[2]{Assumption 11.6 Threshold detectors without memory}{digest-receiver/ass:thresh}
\SetShortId{A-THRESH}
\renewcommand*{\RestatedName}{\hyperref{contract.pdf}{}{claim-receiver/ass:thresh}{Assumption~11.6}}%
\begin{digestassumption}[Threshold detectors without memory]
(i) Each detector reports one bit per detection gate: click or no click;
(ii) there is no photon-number resolution, no dead time, no afterpulsing and
no timing jitter;
(iii) there is no correlation of any kind between gates or between detectors
beyond the declared dark-count channel;
(iv) dark counts are independent across rounds, and across ports unless a
correlated channel is declared explicitly (Definition~\ref{def:darkadm}).
\end{digestassumption}

% claim receiver/ass:port
\digestlabel{ass:port}{11.7}\pdfbookmark[2]{Assumption 11.7 Mode-insensitive port response}{digest-receiver/ass:port}
\SetShortId{A-PORT}
\renewcommand*{\RestatedName}{\hyperref{contract.pdf}{}{claim-receiver/ass:port}{Assumption~11.7}}%
\begin{digestassumption}[Mode-insensitive port response]
(i) Port $k$ clicks if and only if the total photon number in the modes
$\Port{k}$ is nonzero; equivalently, its no-click element is the joint vacuum
projector $\Pi^0_k=\Fock(\id-\Emd{\{k\}})$ on those modes.
\end{digestassumption}

% claim receiver/ass:exact
\digestlabel{ass:exact}{11.8}\pdfbookmark[2]{Assumption 11.8 Exact declared parameters}{digest-receiver/ass:exact}
\SetShortId{A-EXACT}
\renewcommand*{\RestatedName}{\hyperref{contract.pdf}{}{claim-receiver/ass:exact}{Assumption~11.8}}%
\begin{digestassumption}[Exact declared parameters]
(i) Every declared component parameter, efficiency, dark-count probability and
probe mean photon number is exact.
\end{digestassumption}

% claim receiver/ass:nosel
\digestlabel{ass:nosel}{11.9}\pdfbookmark[2]{Assumption 11.9 Device fixed independently of the adversary}{digest-receiver/ass:nosel}
\SetShortId{A-NOSEL}
\renewcommand*{\RestatedName}{\hyperref{contract.pdf}{}{claim-receiver/ass:nosel}{Assumption~11.9}}%
\begin{digestassumption}[Device fixed independently of the adversary]
(i) The component list, its parameters, the port map and the
probe declaration are fixed before the protocol run;
(ii) none of them is chosen by the adversary, and none can be influenced by
her, in particular not through a degree of freedom the model does not carry,
such as wavelength.
\end{digestassumption}

% claim receiver/ass:round
\digestlabel{ass:round}{11.10}\pdfbookmark[2]{Assumption 11.10 Round homogeneity of the instrument, probe and dark counts}{digest-receiver/ass:round}
\SetShortId{A-ROUND}
\renewcommand*{\RestatedName}{\hyperref{contract.pdf}{}{claim-receiver/ass:round}{Assumption~11.10}}%
\begin{digestassumption}[Round homogeneity of the instrument, probe and dark counts]
The family of the instrument, the probe declaration and the dark-count channel
is round-homogeneous (\xref{foundations/def:round-homogeneity}) in both
clauses of that definition:
(i) the same instrument, the same probe declaration and the same dark-count
channel describe every round of the protocol;
(ii) the receiver acts in every round as they describe, whatever the inputs
and outputs of the other rounds, with no memory across rounds.
\end{digestassumption}

% claim receiver/ass:eig
\digestlabel{ass:eig}{11.11}\pdfbookmark[2]{Assumption 11.11 Declared slack on the computed block eigenvalues}{digest-receiver/ass:eig}
\SetShortId{A-EIG}
\renewcommand*{\RestatedName}{\hyperref{contract.pdf}{}{claim-receiver/ass:eig}{Assumption~11.11}}%
\begin{digestassumption}[Declared slack on the computed block eigenvalues]
(i) At every setting $t$ carrying a declared exact-window depth $m^{\ast}_t$
whose emitted weight is that of Definition~\ref{def:pmcex} formed by the
floating-point evaluation of Definition~\ref{def:evaluation-model}, and every
index $N<n\le m^{\ast}_t$, the value
$\widehat\Bnd^{\ast(t)}(n)$ an implementation computes for the exact block
minimum eigenvalue $\Bnd^{\ast(t)}(n)$ of \eqref{eq:exactbnd} satisfies
\begin{equation*}\tag{\getrefnumber{C-eq:eigassumption}}\digesteqlabel{eq:eigassumption}
  \widehat\Bnd^{\ast(t)}(n)-\varepsilon_{\mathrm{eig}}
  \ \le\ \Bnd^{\ast(t)}(n) .
\end{equation*}
\end{digestassumption}

\part*{Engine requirements}\pdfbookmark[0]{Engine requirements}{digest-part-2}

\section*{12\quad Requirements: compilation, representation and outcomes}\pdfbookmark[1]{12 Requirements: compilation, representation and outcomes}{digest-group-receiver/req:graph-refusal}

% claim receiver/req:graph-refusal
\digestlabel{req:graph-refusal}{12.1}\pdfbookmark[2]{Requirement 12.1 Refusal of a graph that is not well formed}{digest-receiver/req:graph-refusal}
\renewcommand*{\RestatedName}{\hyperref{contract.pdf}{}{claim-receiver/req:graph-refusal}{Requirement~12.1}}%
\begin{digestrequirement}[Refusal of a graph that is not well formed]
The engine refuses a declared element graph (input \iface{element_graph})
that fails any of conditions (E1)--(E3) of Definition~\ref{def:graph}, with
the failing condition named, together with the offending element or edge
for (E1) and the offending element pair for (E3); no matrix is
emitted, and no ordering is invented to repair it. The refusal follows the
\iface{fail_closed} convention.
\end{digestrequirement}

% claim receiver/req:probe-diagonal-admission
\digestlabel{req:probe-diagonal-admission}{12.2}\pdfbookmark[2]{Requirement 12.2 Admission of Gaussian-tier probe states only}{digest-receiver/req:probe-diagonal-admission}
\renewcommand*{\RestatedName}{\hyperref{contract.pdf}{}{claim-receiver/req:probe-diagonal-admission}{Requirement~12.2}}%
\begin{digestrequirement}[Admission of Gaussian-tier probe states only]
The engine admits a probe declaration (input \iface{probe_declaration})
only in the Gaussian tier of Definition~\ref{def:probe}, whose states
are occupation-diagonal by construction, and refuses a declared probe
state outside it under the \iface{fail_closed} convention.
\end{digestrequirement}

% claim receiver/req:declared-port-map
\digestlabel{req:declared-port-map}{12.3}\pdfbookmark[2]{Requirement 12.3 Declared port map}{digest-receiver/req:declared-port-map}
\renewcommand*{\RestatedName}{\hyperref{contract.pdf}{}{claim-receiver/req:declared-port-map}{Requirement~12.3}}%
\begin{digestrequirement}[Declared port map]
The engine takes the detector ports as a declared input
(\iface{port_map}) and never infers them from the instrument matrix:
by Remark~\ref{C-rem:unmonitored}, the partition of the output modes into
``in some port'' and ``in no port'' has no observable content beyond the port
Grams it produces.
\end{digestrequirement}

% claim receiver/req:one-mode-per-port
\digestlabel{req:one-mode-per-port}{12.4}\pdfbookmark[2]{Requirement 12.4 One mode per port only where the rank criterion holds}{digest-receiver/req:one-mode-per-port}
\renewcommand*{\RestatedName}{\hyperref{contract.pdf}{}{claim-receiver/req:one-mode-per-port}{Requirement~12.4}}%
\begin{digestrequirement}[One mode per port only where the rank criterion holds]
The engine uses a one-mode-per-port representation of a receiver only where
the criterion of Theorem~\ref{thm:representation-independence}\,(iv), $\rk A_k\le1$ for every $k$,
holds, and never assumes it by default. The criterion is decided on the
emitted port Grams, that is after every detector efficiency is applied
(Remark~\ref{C-rem:effrank}), with each rank decided as
Requirement~\ref{req:tolerance-tests} fixes.
\end{digestrequirement}

% claim receiver/req:declared-alphabet
\digestlabel{req:declared-alphabet}{12.5}\pdfbookmark[2]{Requirement 12.5 Declared emitted alphabet and coarse-graining}{digest-receiver/req:declared-alphabet}
\renewcommand*{\RestatedName}{\hyperref{contract.pdf}{}{claim-receiver/req:declared-alphabet}{Requirement~12.5}}%
\begin{digestrequirement}[Declared emitted alphabet and coarse-graining]
The engine takes the emitted alphabet $\Vset$ as an input
(\iface{emitted_outcome_alphabet}) and the declared coarse-graining
$\cg$ of Definition~\ref{def:coarse} with it
(\iface{outcome_coarse_graining}); it never fixes an outcome alphabet
internally. A receiver that fixed six outcomes over four detectors
internally would make the outcome alphabet an implementation detail of the
receiver rather than a declared property of the protocol, and every
downstream index would then depend on an undeclared convention.
\end{digestrequirement}

\section*{13\quad Requirements: the emitted weight}\pdfbookmark[1]{13 Requirements: the emitted weight}{digest-group-receiver/req:window-declaration}

% claim receiver/req:window-declaration
\digestlabel{req:window-declaration}{13.1}\pdfbookmark[2]{Requirement 13.1 Declared window at every setting whose exact flag is general}{digest-receiver/req:window-declaration}
\renewcommand*{\RestatedName}{\hyperref{contract.pdf}{}{claim-receiver/req:window-declaration}{Requirement~13.1}}%
\begin{digestrequirement}[Declared window at every setting whose exact flag is general]
A declaration supplying no window $m_t$ (\iface{tail_window}) at a setting
where predicate (X) of Definition~\ref{def:pmc} fails is refused
rather than defaulted; by Remark~\ref{C-rem:flag-instances} that condition is
decidable at input validation from declared data alone. The engine
therefore requires $m_t$ exactly at the settings whose exact flag is
$\mathrm{general}$, and at no others. A declared window satisfies $m_t>N$ (Definition~\ref{def:pmc}); one
that does not is refused at every setting, whatever its flag. A window
$m_t>N$ supplied at a setting whose flag is $\mathrm{exact}$ is unused, not
an error.
\end{digestrequirement}

% claim receiver/req:exact-window-declaration
\digestlabel{req:exact-window-declaration}{13.2}\pdfbookmark[2]{Requirement 13.2 Declaration of an exact-window depth and its tolerances}{digest-receiver/req:exact-window-declaration}
\renewcommand*{\RestatedName}{\hyperref{contract.pdf}{}{claim-receiver/req:exact-window-declaration}{Requirement~13.2}}%
\begin{digestrequirement}[Declaration of an exact-window depth and its tolerances]
Declaring an exact-window depth (Definition~\ref{def:pmcex}) at any
setting requires $\varepsilon_{\mathrm{eig}}$ and $\varepsilon_6$ (the bound
on $r_6$ of Definition~\ref{def:residuals}); the engine refuses a
declaration missing either, rather than defaulting it, exactly as a missing
window $m_t$ is refused (Requirement~\ref{req:window-declaration}), with
\iface{exact_window_declared_without_required_tolerance}. The engine refuses
a depth violating $m^{\ast}_t>N$, or violating $m^{\ast}_t\le m_t$ at a
setting where predicate (X) of Definition~\ref{def:pmc} fails, with
\iface{exact_window_depth_outside_declared_range}.
\end{digestrequirement}

% claim receiver/req:exact-window-record
\digestlabel{req:exact-window-record}{13.3}\pdfbookmark[2]{Requirement 13.3 Computed block eigenvalues carried as computed}{digest-receiver/req:exact-window-record}
\renewcommand*{\RestatedName}{\hyperref{contract.pdf}{}{claim-receiver/req:exact-window-record}{Requirement~13.3}}%
\begin{digestrequirement}[Computed block eigenvalues carried as computed]
Nothing is clamped, and the computed value is emitted: at every setting
carrying an exact-window depth, the engine carries each
$\widehat\Bnd^{\ast(t)}(n)$, $N<n\le m^{\ast}_t$, in the artifact exactly as
computed, in \iface{bound_quantifier_domain}, and
$\varepsilon_{\mathrm{eig}}$ enters only as the declared slack subtracted in
\eqref{eq:elldef} (Definition~\ref{def:pmcex}), so that a reader can
recompute \eqref{eq:elldef} and see what the slack cost.
\end{digestrequirement}

% claim receiver/req:weight-provenance
\digestlabel{req:weight-provenance}{13.4}\pdfbookmark[2]{Requirement 13.4 Provenance of each emitted multiclick weight}{digest-receiver/req:weight-provenance}
\renewcommand*{\RestatedName}{\hyperref{contract.pdf}{}{claim-receiver/req:weight-provenance}{Requirement~13.4}}%
\begin{digestrequirement}[Provenance of each emitted multiclick weight]
The engine evaluates an exact window, that of Definition~\ref{def:pmcex} or
the rigorous window bounds of Definition~\ref{def:enclosed-weight}, only
where a declaration requests it by an exact-window depth, and never by
default. Where it is not requested, the engine may also report block
eigenvalues as a diagnostic; the values entering the strengthening slot
(Definition~\ref{def:slot}), \iface{high_photon_bounds}, are the
per-setting weights $\pmc^{(t)}$ of
Theorem~\ref{thm:emitted-bound} --- at a receiver without a probe the single
value $\pmc$ --- each the enclosed weight of
Definition~\ref{def:enclosed-weight}, rounded down, where
Requirement~\ref{req:rigorous-weight} forms it, and otherwise the
floating-point evaluation of Definition~\ref{def:pmc} or~\ref{def:pmcex}; and
the artifact says, for each, which evaluation produced it, which of
Definitions~\ref{def:pmc} and~\ref{def:pmcex} gives its formula and at what
depth, in \iface{bound_quantifier_domain}.
\end{digestrequirement}

% claim receiver/req:rigorous-weight
\digestlabel{req:rigorous-weight}{13.5}\pdfbookmark[2]{Requirement 13.5 Rigorous evaluation of the multiclick weight}{digest-receiver/req:rigorous-weight}
\renewcommand*{\RestatedName}{\hyperref{contract.pdf}{}{claim-receiver/req:rigorous-weight}{Requirement~13.5}}%
\begin{digestrequirement}[Rigorous evaluation of the multiclick weight]
\begin{enumerate}[label=(\roman*),leftmargin=2.4em]
\item \emph{Where it applies.} On every declaration whose port Grams the
  engine encloses by Requirement~\ref{req:exact-route-decision}\,(iii), the
  engine evaluates the weight of every setting by (ii) to (iv), in exact
  rational arithmetic: the rigorous evaluation of
  Definition~\ref{def:evaluation-model}. At a setting where it cannot, it
  emits the weight of Definition~\ref{def:pmc} or~\ref{def:pmcex} evaluated in
  floating point, with the reason \iface{port_grams_not_enclosed} where the
  declaration has no model admissible exactly and
  \iface{bound_inputs_not_enclosed} where a reciprocal that (ii) forms
  receives no enclosure.
\item \emph{The enclosures.} At each setting $t$ the engine forms, from the
  enclosures of the port Grams by the rules of
  Lemma~\ref{C-lem:enclosure-arithmetic},
  complex enclosures of $B^{(t)}$ and of every
  $B^{(t)}_k$, of $\gamma^{(t)}$ and of every $\beta^{(t)}_k$, and, where $t$
  carries an exact-window depth $m^{\ast}_t$, of
  $\blk{M}{\eqn}_{\mc|t}$ for $N<n\le m^{\ast}_t$, each enclosure of one of
  these Hermitian matrices made Hermitian-centred by part~(v) of that lemma.
  For a receiver without a probe these are, by
  Corollary~\ref{cor:unconditioned-family}, $B=A_\Sigma$ and
  $B_k=A_{[K]\setminus\{k\}}$, sums of enclosed Grams, $\gamma=\beta_k=1$
  exactly, and the blocks of \eqref{eq:family-mc}, formed from the enclosures
  of Theorem~\ref{thm:enclosed-blocks}\,(ii) taken with $m^{\ast}_t$ as its
  cutoff. For a setting family a submodule constructs they are the
  enclosures of its prefactors, its operators and its window blocks that the
  submodule states, formed by the same rules.
\item \emph{The four inputs and the window bounds.} For an
  entrywise enclosure $(C,r)$ of a Hermitian matrix $X$ and a rational
  $\rho\ge0$ with $\rho^{2}\ge\sum_{i,l}r_{il}^{2}$, an upper bound on
  $\lmax(C)$ from \xref{foundations/lem:least-eigenvalue-bound}\,(iv) plus
  $\rho$ bounds $\lmax(X)$ from above, and a lower bound on $\lmin(C)$ from
  part~(ii) of that lemma less $\rho$ bounds $\lmin(X)$ from below
  (\xref{foundations/cor:least-eigenvalue-over-enclosure}\,(i),~(ii)); for a
  complex enclosure $(x_c,r)$ of a real $x$,
  $\mathrm{Re}\,x_c-r\le x\le\mathrm{Re}\,x_c+r$. By these rules the engine
  forms an upper bound $u$ on $\lmax(B^{(t)})$, lower bounds $\ell_k$ on
  $\lmin(B^{(t)}_k)$, a lower bound $g$ on $\gamma^{(t)}$, upper bounds $b_k$
  on $\beta^{(t)}_k$ and lower bounds $\lambda_n$ on
  $\lmin(\blk{M}{\eqn}_{\mc|t})$. Its enclosed bound data are
  \eqref{eq:enclosed-construction} of Theorem~\ref{thm:enclosed-weight}\,(i),
  each of the four intersected with $[0,1]$, and its
  rigorous window bounds are $\max\{0,\lambda_n\}$. The output of a
  floating-point eigensolver may enter only as the basis $V$ of
  \xref{foundations/lem:least-eigenvalue-bound}, read as the exact rational
  matrix it is; a basis for which the engine forms no rational $\theta<1$
  satisfying the hypothesis of that lemma is replaced by the identity, which
  admits $\theta=0$. No computed eigenvalue is used as a bound, and
  $\varepsilon_{\mathrm{eig}}$ is not subtracted.
\item \emph{The weight.} The engine computes the enclosed weight
  $\check p^{(t)}_{\mc}$ (Definition~\ref{def:enclosed-weight}) as an exact
  rational and emits, as the setting's weight in the
  strengthening slot \iface{high_photon_bounds}
  (Definition~\ref{def:slot}), the largest binary64 number at or below it,
  rounded down as \xref{foundations/req:rounding-direction} rounds an emitted
  lower bound; by Theorem~\ref{thm:enclosed-weight}\,(iii) the emitted number
  bounds the multiclick element of its setting on every block above the
  cutoff. The uniform weight is the minimum of the emitted per-setting
  weights.
\item \emph{What stays.} The floating-point evaluation of the blocks, the
  residuals and their tests, and, at a setting carrying a depth, the computed
  values of Requirement~\ref{req:exact-window-record}, are formed as they are
  without this requirement. At a setting evaluated by (ii) to (iv), the
  floating-point weight of Definition~\ref{def:pmc} or~\ref{def:pmcex} enters
  no emitted number, and it may be recorded beside the emitted one as a
  diagnostic.
\item \emph{The record.} Per setting, in \iface{bound_quantifier_domain}: the
  evaluation that produced the emitted weight, \iface{rigorous} or
  \iface{floating_point}, with the reason of (i) where it is
  \iface{floating_point}; and, where it is \iface{rigorous}, the enclosed
  bound data, the rigorous window bounds and the enclosed weight as exact
  rationals, so that a reader can recompute the weight. A setting whose weight
  is \iface{rigorous} is not one at which Assumption~\ref{ass:eig} is invoked
  (Requirement~\ref{C-req:eig-ledger-row}), and the artifact's evaluation label
  follows Requirement~\ref{req:evaluation-label}.
\end{enumerate}
\end{digestrequirement}

\section*{14\quad Requirements: residuals, tolerances and the evaluation label}\pdfbookmark[1]{14 Requirements: residuals, tolerances and the evaluation label}{digest-group-receiver/req:residual-emission}

% claim receiver/req:residual-emission
\digestlabel{req:residual-emission}{14.1}\pdfbookmark[2]{Requirement 14.1 Emission of the residuals that exist on the route and request}{digest-receiver/req:residual-emission}
\renewcommand*{\RestatedName}{\hyperref{contract.pdf}{}{claim-receiver/req:residual-emission}{Requirement~14.1}}%
\begin{digestrequirement}[Emission of the residuals that exist on the route and request]
The engine emits with every artifact, in \iface{validation_results}, those
residuals of Definition~\ref{def:residuals} that exist on the input route taken
and for the evaluation requested, and no other: the keys of $r_1$ and $r_3$ are
absent, not zero, on route (i), and the key of $r_6$ is absent, not zero, at a
declaration that requests no exact window. The blocks over which
\eqref{eq:res5} maximises are the floating-point blocks
$\blk{M}{\eqn}_{v|t}$ of the coarse-grained elements that
Definition~\ref{def:residuals} names, before any dark-count channel, which
the \iface{photon_blocks} entry carries beside its dark-composed blocks
(Requirement~\ref{req:dark-composed-reference}); $r_5$ does not run over the
dark-composed blocks, on which the exact value $0$ of
Proposition~\ref{prop:res5}\,(iv) is not asserted.
$\varepsilon_1$ and $\varepsilon_3$ are declared on every route, route (i)
included, where they are unused (\iface{numerical_tolerances});
$\varepsilon_6$ is required exactly where an exact-window depth is
declared, on the same condition as $r_6$, and where it is declared elsewhere
it is unused and not an error.
\end{digestrequirement}

% claim receiver/req:residual-refusal
\digestlabel{req:residual-refusal}{14.2}\pdfbookmark[2]{Requirement 14.2 Refusal on a residual above its tolerance, and no repair}{digest-receiver/req:residual-refusal}
\renewcommand*{\RestatedName}{\hyperref{contract.pdf}{}{claim-receiver/req:residual-refusal}{Requirement~14.2}}%
\begin{digestrequirement}[Refusal on a residual above its tolerance, and no repair]
The engine compares each residual that exists (Definition~\ref{def:residuals})
with its receiver tolerance $\varepsilon_i$, and a
residual above its tolerance is a \emph{refusal}, not a warning; in particular
$r_5>\varepsilon_5$ refuses with the reason
\iface{photon_block_positivity_residual_above_tolerance}, and
$r_6>\varepsilon_6$ with the reason
\iface{window_eigenvalue_consistency_residual_above_tolerance}. In the same
way, at every setting where predicate (X) of
Definition~\ref{def:pmc} fails, the engine compares the residual gap
$g_{N,m_t}$ of \eqref{eq:resgap} on the declared window $m_t$ with
$\varepsilon_{\mathrm{gap}}$, and a gap above it refuses. The direction
is refuse-only: no block is Hermitized, symmetrized, projected onto the
positive cone, clamped or repaired in any way; the artifact carries the
computed block itself, and $r_5$ reports that block's own defect. The same
holds of $r_6$: it refuses and repairs nothing (Remark~\ref{C-rem:exacteig}
states the direction of failure it is blind to).
\end{digestrequirement}

% claim receiver/req:tolerance-tests
\digestlabel{req:tolerance-tests}{14.3}\pdfbookmark[2]{Requirement 14.3 Decisions at the receiver tolerances, and refusal of failed bound data}{digest-receiver/req:tolerance-tests}
\renewcommand*{\RestatedName}{\hyperref{contract.pdf}{}{claim-receiver/req:tolerance-tests}{Requirement~14.3}}%
\begin{digestrequirement}[Decisions at the receiver tolerances, and refusal of failed bound data]
The engine decides each exact condition that Definition~\ref{def:tolmap} lists
by the test given there and by no other threshold, and every rank statement
this module decides on floating-point data by \eqref{eq:eps2rank}; no other
rule is licensed, there or at Theorem~\ref{thm:representation-independence}\,(iv). Bound data failing
item (iv) of Definition~\ref{def:tolmap} are a failed validation, so
\iface{fail_closed} applies at this site: no artifact is emitted. The reason
string is the implementer's to name, under the rule the refusal rules of the
module's interface state for the validations they do not tabulate, none of
the tabulated strings
covering this check.
\end{digestrequirement}

% claim receiver/req:tolerance-domain
\digestlabel{req:tolerance-domain}{14.4}\pdfbookmark[2]{Requirement 14.4 Validation of the receiver tolerances against their domain}{digest-receiver/req:tolerance-domain}
\renewcommand*{\RestatedName}{\hyperref{contract.pdf}{}{claim-receiver/req:tolerance-domain}{Requirement~14.4}}%
\begin{digestrequirement}[Validation of the receiver tolerances against their domain]
The engine validates each of the receiver tolerances
$\varepsilon_1,\dots,\varepsilon_6$, $\varepsilon_{\mathrm{gap}}$ and
$\varepsilon_{\mathrm{eig}}$ of Definition~\ref{def:tolmap} at input against
the declared domain $[0,1)$, and refuses a value outside that interval with the
reason \iface{tolerance_outside_declared_domain}.
\end{digestrequirement}

% claim receiver/req:evaluation-label
\digestlabel{req:evaluation-label}{14.5}\pdfbookmark[2]{Requirement 14.5 The evaluation label on every artifact}{digest-receiver/req:evaluation-label}
\renewcommand*{\RestatedName}{\hyperref{contract.pdf}{}{claim-receiver/req:evaluation-label}{Requirement~14.5}}%
\begin{digestrequirement}[The evaluation label on every artifact]
The engine writes \iface{bound_evaluation} on every artifact of this module,
on every input route: \iface{rigorous} where the weight of every setting was
formed by the rigorous evaluation of Requirement~\ref{req:rigorous-weight},
and \iface{floating_point} otherwise, an artifact with a weight resting on
Assumption~\ref{ass:eig} among them; which evaluation formed each setting's
weight is recorded as Requirement~\ref{req:rigorous-weight}\,(vi) says. No
path in this module produces another label
(Theorem~\ref{thm:exact-and-evaluated}\,(iii)): the label \iface{certified},
which would name a tier and not an evaluation, and the retired fallback reason
\iface{receiver_grams_not_enclosed} are never produced. A run whose
declaration carries a weight into which a per-setting weight enters
that the record of Requirement~\ref{req:rigorous-weight}\,(vi) names as
evaluated in floating point, the
uniform weight of a slot entry counting
the weight of every setting as entering it, is not reported at the certified
tier
(\xref{meat-core/frank-wolfe/req:certified-evaluation}\,(vii)); in every other
respect the tier a run reports from these artifacts follows
\xref{foundations/req:tiers}.
\end{digestrequirement}

% claim receiver/req:exact-route-decision
\digestlabel{req:exact-route-decision}{14.6}\pdfbookmark[2]{Requirement 14.6 Deciding the exact route and enclosing the port Grams}{digest-receiver/req:exact-route-decision}
\renewcommand*{\RestatedName}{\hyperref{contract.pdf}{}{claim-receiver/req:exact-route-decision}{Requirement~14.6}}%
\begin{digestrequirement}[Deciding the exact route and enclosing the port Grams]
\begin{enumerate}[label=(\roman*),leftmargin=2.4em]
\item \emph{Declared Grams or matrix.} On a route that
  declares the port Grams or the instrument matrix, the engine computes
  the port Grams in exact rational arithmetic from the
  declared values, by \eqref{eq:gramdef} where the
  instrument matrix is declared, and decides in exact rational arithmetic
  whether every $A_k$ is Hermitian and whether the family satisfies
  \eqref{eq:exact-admissible}, that is, whether it is an
  admissible Gram family exactly. The declaration has a model
  admissible exactly, and is then on the exact route
  (Definition~\ref{def:exact-route}), exactly when both hold.
\item \emph{Component graph.} On the route that declares a component graph,
  the engine decides in exact rational arithmetic whether the matrix of every
  generic block, whose entries are declared values, is a contraction,
  $\Gmat_\Sigma^{\dagger}\Gmat_\Sigma\preceq\id$. Where every one is, the
  declaration has a model admissible exactly: the compiled instrument matrix
  is a contraction (Theorem~\ref{thm:compilation-well-defined}\,(ii)) and its
  port Grams are an admissible Gram family (Theorem~\ref{thm:port-grams}\,(i)),
  at the exact values of the declared parameters. For a declaration that
  has one and whose parameters its exact arithmetic represents, the engine
  moreover compiles the instrument matrix of Definition~\ref{def:compile} and
  the port Grams exactly and decides whether every port Gram has entries in
  $\mathbb{Q}(i)$, which places the declaration on the exact route. Every
  transmittance, loss, reflectivity, extinction and efficiency is the
  rational number the artifact carries for it: the declared value, or, for a
  loss declared in decibels, the converted efficiency that
  \iface{presentation_rule} makes the declared datum, a binary64 number and
  so a dyadic rational (\xref{foundations/def:binary64}\,(iii)). An angle
  parameter given in the presentation \iface{angle_presentation}, as a
  rational coefficient $c$ of $\pi$, is $c\pi$ exactly; an angle given in
  radians is a declared value, the dyadic rational it denotes. The
  arithmetic represents at least the field generated over $\mathbb{Q}(i)$ by
  $\sqrt2$ and by $\sqrt p$ and $\sqrt{1-p}$ for every declared rational
  parameter $p$ in $[0,1]$, which contains every entry of every catalogue
  entry whose angles are given as coefficients of $\pi$, whose axis angles
  are multiples of $\pi/8$, whose retardances are multiples of $\pi/2$ and
  whose rotations and phases are multiples of $\pi/4$. A declaration whose
  parameters the arithmetic does not represent, a non-zero angle given in
  radians among them, is treated as off the exact route, with its cause
  recorded by (iv), even where its port Grams happen to have entries in
  $\mathbb{Q}(i)$ (two half-wave plates in sequence at one such angle, for
  instance, compose to $-\id$).
\item \emph{Enclosure of the port Grams.} For every declaration with a
  model admissible exactly, the engine computes, in exact rational
  arithmetic, complex enclosures
  (Definition~\ref{def:complex-enclosure}) of its port Grams at the exact
  values of the declared parameters. On the exact route as decided by (i)
  or (ii) each port Gram is enclosed exactly, by itself. Otherwise, on the
  route of (ii), the
  enclosures are those of Theorem~\ref{thm:enclosed-blocks}\,(i), with the
  matrix of every element enclosed by the rules of
  Lemma~\ref{C-lem:enclosure-arithmetic} applied to the formula of its
  catalogue entry from square-root enclosures
  and angle enclosures of its parameters, and
  $\sqrt{\eta_k}$ by a square-root enclosure; where the exact arithmetic of
  (ii) has computed a port Gram, its entries, which are expressions in square
  roots of rationals, may be enclosed from those expressions instead. A
  phase $e^{i\varphi}$ is $\cos\varphi+i\sin\varphi$, and the angle $\delta/2$
  of a retarder is again a rational or a rational multiple of $\pi$. No
  floating-point value enters an enclosure.
\item \emph{Record.} The \iface{photon_blocks} entry of the artifact records
  whether (i) or (ii) places the declaration on the exact route and, where
  neither does, why:
  a parameter or an angle that the exact arithmetic of (ii) does not
  represent (a non-zero angle given in radians among them), port Grams that
  are not Gaussian-rational, or no model admissible exactly, which is
  declared port Grams that are not Hermitian, a declared family that is not
  an admissible Gram family exactly (a declaration admitted only at the
  receiver tolerances), or a generic block that is not a contraction
  exactly. Being off the exact route is not a refusal of the receiver and
  takes nothing away from the floating-point blocks (Definition~\ref{def:evaluation-model});
  where the declaration has a model admissible exactly, its port Grams are
  enclosed by (iii).
\end{enumerate}
\end{digestrequirement}

\section*{15\quad Requirements: the assumptions}\pdfbookmark[1]{15 Requirements: the assumptions}{digest-group-receiver/req:decidable-assumptions}

% claim receiver/req:decidable-assumptions
\digestlabel{req:decidable-assumptions}{15.1}\pdfbookmark[2]{Requirement 15.1 Validation of the assumptions on the declared data}{digest-receiver/req:decidable-assumptions}
\renewcommand*{\RestatedName}{\hyperref{contract.pdf}{}{claim-receiver/req:decidable-assumptions}{Requirement~15.1}}%
\begin{digestrequirement}[Validation of the assumptions on the declared data]
The engine enforces by validation, wherever it is decidable, every clause of
an assumption of this contract that constrains the declared data. The
assumptions that the engine never verifies, as the contract says of each
(Remark~\ref{C-rem:unchecked-assumptions}), form \iface{never_verified_set},
and no code path decides what they say beyond those clauses.
\end{digestrequirement}

% claim receiver/req:reserved-tolerance-field
\digestlabel{req:reserved-tolerance-field}{15.2}\pdfbookmark[2]{Requirement 15.2 Refusal of a non-zero parameter tolerance}{digest-receiver/req:reserved-tolerance-field}
\renewcommand*{\RestatedName}{\hyperref{contract.pdf}{}{claim-receiver/req:reserved-tolerance-field}{Requirement~15.2}}%
\begin{digestrequirement}[Refusal of a non-zero parameter tolerance]
The engine validates the reserved \iface{tolerance} field of every declared
parameter as absent or exactly zero, and refuses it otherwise.
\end{digestrequirement}

\part*{Component catalogue}\pdfbookmark[0]{Component catalogue}{digest-part-3}

\section*{16\quad Catalogue entries and their admissibility}\pdfbookmark[1]{16 Catalogue entries and their admissibility}{digest-group-receiver/component-catalogue/lem:attenuated}

% claim receiver/component-catalogue/def:catalogue-entries
\digestlabel{def:catalogue-entries}{16.2}\pdfbookmark[2]{Definition 16.2 Entries of the component catalogue}{digest-receiver/component-catalogue/def:catalogue-entries}
\renewcommand*{\RestatedName}{\hyperref{contract.pdf}{}{claim-receiver/component-catalogue/def:catalogue-entries}{Definition~16.2}}%
\begin{digestclaim}[Entries of the component catalogue]
The \emph{catalogue entries} are the entry types C1--C10 below.
Throughout, a \emph{path} is a spatial index $p$ and carries the two
polarization modes $(p,H)$ and $(p,V)$; a matrix written on two paths acts as
$\cdot\otimes\id_2$ on polarization unless stated otherwise, and a matrix
written on polarization acts on one path.

\emph{C1 --- Beamsplitter.} Support: two paths $(p_1,p_2)$, four modes.
Parameters: transmittance $x\in[0,1]$, reflection phase
$\varphi\in[0,2\pi)$, per-output losses $y_1,y_2\in[0,1]$. Core, on the ordered
path pair with rows $=$ outputs:
\begin{equation*}\tag{\getrefnumber{C-eq:bs}}\digesteqlabel{eq:bs}
  \Bsp(x,\varphi)=
  \begin{pmatrix}
    \sqrt{x} & \sqrt{1-x}\,e^{i\varphi}\\[2pt]
    \sqrt{1-x}\,e^{-i\varphi} & -\sqrt{x}
  \end{pmatrix},
  \qquad
  \Gmat_{\mathrm{C1}}=
  \Big[\operatorname{diag}\big(\sqrt{1-y_1},\sqrt{1-y_2}\big)\,
       \Bsp(x,\varphi)\Big]\otimes\id_2 .
\end{equation*}
A \emph{polarization-dependent beamsplitter} declares $(x_H,\varphi_H)$ and
$(x_V,\varphi_V)$ separately, with the per-output losses $y_1,y_2$ of C1, and
is
$\operatorname{diag}\big(\sqrt{1-y_1},\sqrt{1-y_2},\sqrt{1-y_1},\sqrt{1-y_2}\big)
\big[\Bsp(x_H,\varphi_H)\oplus\Bsp(x_V,\varphi_V)\big]$ in the reordered basis
$(p_1H,p_2H,p_1V,p_2V)$; this is the matrix of entry C2 at
$(t_H,\varphi_H,t_V,\varphi_V)=(x_H,\varphi_H,x_V,\varphi_V)$ written in that
basis.

\emph{C2 --- Polarizing beamsplitter.} Support: two paths, four modes.
Parameters: $t_H,t_V\in[0,1]$ (ideal $t_H=1$, $t_V=0$), phases
$\varphi_H,\varphi_V\in[0,2\pi)$, losses $y_1,y_2\in[0,1]$. In the ordering
$(p_1H,p_1V,p_2H,p_2V)$,
\begin{equation*}\tag{\getrefnumber{C-eq:pbs}}\digesteqlabel{eq:pbs}
  \Gmat_{\mathrm{C2}}=
  \operatorname{diag}\big(\sqrt{1-y_1},\sqrt{1-y_1},
                          \sqrt{1-y_2},\sqrt{1-y_2}\big)\;
  \Xi^\dagger\big[\Bsp(t_H,\varphi_H)\oplus\Bsp(t_V,\varphi_V)\big]\,\Xi ,
\end{equation*}
where $\Xi$ is the permutation from $(p_1H,p_1V,p_2H,p_2V)$ to
$(p_1H,p_2H,p_1V,p_2V)$.

\emph{C3 --- \emph{Retarder}, with half- and quarter-wave presets.} Support: the two
polarization modes of one path. Parameters: retardance $\delta\in[0,2\pi)$,
fast-axis angle $\theta\in[0,\pi)$, scalar loss $y\in[0,1]$. With
$\Rot(\theta)=\left(\begin{smallmatrix}\cos\theta&-\sin\theta\\
\sin\theta&\cos\theta\end{smallmatrix}\right)$,
\begin{equation*}\tag{\getrefnumber{C-eq:retarder}}\digesteqlabel{eq:retarder}
  \Wplate(\delta,\theta)
  \coloneqq e^{-i\delta/2}\,\Rot(\theta)\operatorname{diag}(1,e^{i\delta})
    \Rot(\theta)^{\transpose}
  = e^{-i\delta/2}
  \begin{pmatrix}
    \cos^2\theta+e^{i\delta}\sin^2\theta &
    (1-e^{i\delta})\sin\theta\cos\theta\\[2pt]
    (1-e^{i\delta})\sin\theta\cos\theta &
    \sin^2\theta+e^{i\delta}\cos^2\theta
  \end{pmatrix},
\end{equation*}
and $\Gmat_{\mathrm{C3}}=\sqrt{1-y}\,\Wplate(\delta,\theta)$.
The presets are
\begin{equation*}\tag{\getrefnumber{C-eq:hwpqwp}}\digesteqlabel{eq:hwpqwp}
  \mathrm{HWP}(\theta)\coloneqq\Wplate(\pi,\theta)
  =-i\begin{pmatrix}\cos2\theta&\sin2\theta\\
                    \sin2\theta&-\cos2\theta\end{pmatrix},
  \qquad
  \mathrm{QWP}(\theta)\coloneqq\Wplate(\tfrac{\pi}{2},\theta),
\end{equation*}
so that $\mathrm{HWP}(\pi/8)=-\tfrac{i}{\sqrt2}
\left(\begin{smallmatrix}1&1\\1&-1\end{smallmatrix}\right)$.

\emph{C4 --- Circular retarder, or polarization \emph{rotator}.} Support: the two
polarization modes of one path. Parameter: rotation $\rho\in[0,2\pi)$.
$\Gmat_{\mathrm{C4}}=\Rot(\rho)$.

\emph{C5 --- Mirror or fold.} Support: one path. Parameters: reflectivities
$R_H,R_V\in[0,1]$, phases $\varphi_H,\varphi_V\in[0,2\pi)$, handedness flip
$\mathrm{fl}\in\{\mathrm{false},\mathrm{true}\}$ (default $\mathrm{false}$).
\begin{equation*}\tag{\getrefnumber{C-eq:mirror}}\digesteqlabel{eq:mirror}
  \Gmat_{\mathrm{C5}}=\operatorname{diag}\big(\sqrt{R_H},\sqrt{R_V}\big)\,
  \operatorname{diag}\big(e^{i\varphi_H},e^{i\varphi_V}\big)\,
  F,\qquad
  F=\begin{cases}\operatorname{diag}(1,-1)&\mathrm{fl}=\mathrm{true}\\
                 \id_2&\text{otherwise.}\end{cases}
\end{equation*}

\emph{C6 --- Phase shifter.} Support: any declared mode subset $\Sigma$.
Parameters: $\varphi_m\in[0,2\pi)$ per $m\in\Sigma$.
$\Gmat_{\mathrm{C6}}=\operatorname{diag}(e^{i\varphi_m})_{m\in\Sigma}$.

\emph{C7 --- Loss.} Support: any declared mode subset $\Sigma$. Parameters:
$\eta_m\in[0,1]$ per $m\in\Sigma$, optionally declared as
$L_{\mathrm{dB}}\ge0$ with $\eta=10^{-L_{\mathrm{dB}}/10}$.
$\Gmat_{\mathrm{C7}}=\operatorname{diag}(\sqrt{\eta_m})_{m\in\Sigma}$.

\emph{C8 --- Polarizer.} Support: the two polarization modes of one path.
Parameters: transmission axis $\theta\in[0,\pi)$, pass amplitude $\sqrt{t}$ and
leak amplitude $\sqrt{t_\perp}$ with $t,t_\perp\in[0,1]$. With
$\ket{e_\theta}=\cos\theta\ket{H}+\sin\theta\ket{V}$,
\begin{equation*}\tag{\getrefnumber{C-eq:polarizer}}\digesteqlabel{eq:polarizer}
  \Gmat_{\mathrm{C8}}
  =\sqrt{t}\,\dyad{e_\theta}+\sqrt{t_\perp}\,\dyad{e_\theta^{\perp}}
  =\Rot(\theta)\operatorname{diag}\big(\sqrt{t},\sqrt{t_\perp}\big)
   \Rot(\theta)^{\transpose}.
\end{equation*}

\emph{C9 --- Mode routing.} Support: any declared mode subset.
Parameter: a permutation $\pi$ of the modes of $\Sigma$;
$(\Gmat_{\mathrm{C9}})_{ij}=\delta_{i,\pi(j)}$.

\emph{C10 --- \emph{Generic block}.} Support: any declared mode subset $\Sigma$.
Parameter: a matrix $\Gmat_\Sigma\in\mathbb{C}^{\Sigma\times\Sigma}$, admitted
after the explicit check $\Gmat_\Sigma^\dagger\Gmat_\Sigma\preceq\id_\Sigma$.

Table~\ref{C-tab:catalogue} collects the supports and the parameter domains.

\begin{table}[htbp]
\centering
{\small
\begin{tabular}{@{}l >{\raggedright\arraybackslash}p{0.18\textwidth}
                  >{\raggedright\arraybackslash}p{0.30\textwidth}
                  >{\raggedright\arraybackslash}p{0.32\textwidth}@{}}
\hline
ID & Support & Parameters (domain) &
Lossless exactly when \\
\hline
C1 & 2 paths (4 modes) & $x\in[0,1]$, $\varphi\in[0,2\pi)$,
$y_1,y_2\in[0,1]$ & $y_1=y_2=0$ \\
C2 & 2 paths (4 modes) & $t_H,t_V\in[0,1]$, $\varphi_H,\varphi_V\in[0,2\pi)$,
$y_1,y_2\in[0,1]$ & $y_1=y_2=0$ \\
C3 & 1 path, polarization & $\delta\in[0,2\pi)$, $\theta\in[0,\pi)$,
$y\in[0,1]$ & $y=0$ \\
C4 & 1 path, polarization & $\rho\in[0,2\pi)$ & always \\
C5 & 1 path, polarization & $R_H,R_V\in[0,1]$,
$\varphi_H,\varphi_V\in[0,2\pi)$, $\mathrm{fl}$ & $R_H=R_V=1$ \\
C6 & any $\Sigma$ & $\varphi_m\in[0,2\pi)$ & always \\
C7 & any $\Sigma$ & $\eta_m\in[0,1]$ & every $\eta_m=1$ \\
C8 & 1 path, polarization & $\theta\in[0,\pi)$, $t,t_\perp\in[0,1]$ &
$t=t_\perp=1$ (then $\Gmat_\Sigma=\id$) \\
C9 & any $\Sigma$ & a permutation of $\Sigma$ & always \\
C10 & any $\Sigma$ & $\Gmat_\Sigma$ with
$\Gmat_\Sigma^\dagger\Gmat_\Sigma\preceq\id$ (checked) &
$\Gmat_\Sigma^\dagger\Gmat_\Sigma=\id$ (checked) \\
\hline
\end{tabular}}
\caption{The catalogue entries. Support shapes are stated in modes; the
parameter domains are the declared domains of
Definition~\ref{def:graph}\,(E1); the last column names the subfamily on
which Theorem~\ref{thm:catalogue-admissibility} proves the entry lossless.}
\end{table}
\end{digestclaim}

% claim receiver/component-catalogue/thm:catalogue-admissibility
\digestlabel{thm:catalogue-admissibility}{16.3}\pdfbookmark[2]{Theorem 16.3 Admissibility of the catalogue entries}{digest-receiver/component-catalogue/thm:catalogue-admissibility}
\begin{mainclaim}
\renewcommand*{\RestatedName}{\hyperref{contract.pdf}{}{claim-receiver/component-catalogue/thm:catalogue-admissibility}{Theorem~16.3}}%
\begin{digestclaim}[Admissibility of the catalogue entries]
Let $\Gmat_\Sigma$ be the matrix of an entry type C1--C10 of
Definition~\ref{def:catalogue-entries} at a parameter tuple in its declared
domain. Then:
\begin{enumerate}[label=(\roman*),leftmargin=2.2em]
\item $\Gmat_\Sigma^\dagger\Gmat_\Sigma\preceq\id_\Sigma$, so $\Gmat_\Sigma$
  is a local element in the sense of Definition~\ref{def:embed};
\item $\Gmat_\Sigma$ is lossless,
  $\Gmat_\Sigma^\dagger\Gmat_\Sigma=\id_\Sigma$, exactly on the subfamily
  named in Table~\ref{C-tab:catalogue};
\item for C1--C9, a matrix whose parameter tuple passes the domain comparison
  of (E1) (Definition~\ref{def:graph}) is contractive by (i); C10 is the only
  entry whose declared domain is the contractivity condition itself, so that
  its check of Theorem~\ref{thm:compilation-well-defined}\,(iv)(E1) is a
  spectral test on the supplied matrix.
\end{enumerate}
\end{digestclaim}
\end{mainclaim}

% claim receiver/component-catalogue/lem:catalogue-dilations
\digestlabel{lem:catalogue-dilations}{16.4}\pdfbookmark[2]{Lemma 16.4 Loss factors of the catalogue entries and the dilation of a graph}{digest-receiver/component-catalogue/lem:catalogue-dilations}
\renewcommand*{\RestatedName}{\hyperref{contract.pdf}{}{claim-receiver/component-catalogue/lem:catalogue-dilations}{Lemma~16.4}}%
\begin{digestclaim}[Loss factors of the catalogue entries and the dilation of a graph]
For a diagonal matrix $D$ write $D^{\flat}$ for the matrix of the rows of $D$
whose diagonal entry is not zero, so that
$(D^{\flat})^{\dagger}D^{\flat}=D^{\dagger}D$.
\begin{enumerate}[label=(\roman*),leftmargin=2.2em]
\item \emph{Loss factors.} Let $\Gmat_\Sigma$ be the matrix of an entry type
  of Definition~\ref{def:catalogue-entries} at a parameter tuple in its
  declared domain, in the notation of that definition, and let $\Theta$ be
  the matrix given for its type below, each of its rows indexed by a mode of
  its own:
  \begin{itemize}
  \item C1: $\Theta=\big(\operatorname{diag}(\sqrt{y_1},\sqrt{y_2})^{\flat}\,\Bsp(x,\varphi)\big)\otimes\id_2$;
  \item C2: $\Theta=\operatorname{diag}(\sqrt{y_1},\sqrt{y_1},\sqrt{y_2},\sqrt{y_2})^{\flat}\,
    \Xi^{\dagger}\big[\Bsp(t_H,\varphi_H)\oplus\Bsp(t_V,\varphi_V)\big]\Xi$;
  \item C3: $\Theta=\operatorname{diag}(\sqrt{y},\sqrt{y})^{\flat}$;
  \item C4, C6 and C9: $\Theta$ with no row;
  \item C5: $\Theta=\operatorname{diag}\big(\sqrt{1-R_H},\sqrt{1-R_V}\big)^{\flat}$;
  \item C7: $\Theta=\operatorname{diag}\big(\sqrt{1-\eta_m}\big)_{m\in\Sigma}^{\flat}$;
  \item C8: $\Theta=\operatorname{diag}\big(\sqrt{1-t},\sqrt{1-t_\perp}\big)^{\flat}\,\Rot(\theta)^{\transpose}$;
  \item C10, with the entries of $\Gmat_\Sigma$ in $\mathbb{Q}(i)$:
    $\Theta=D^{1/2}E$, where
    $\id_\Sigma-\Gmat_\Sigma^{\dagger}\Gmat_\Sigma=E^{\dagger}DE$ with $D$
    diagonal with positive rational diagonal entries and $E$ with entries in
    $\mathbb{Q}(i)$; such a decomposition exists, with at most $|\Sigma|$
    rows.
  \end{itemize}
  Then $\Gmat_\Sigma^{\dagger}\Gmat_\Sigma+\Theta^{\dagger}\Theta=\id_\Sigma$,
  and $\Theta$ is a \emph{loss factor} of the entry. For C1--C9, $\Theta$ has
  no row exactly on the lossless subfamily named in
  Table~\ref{C-tab:catalogue}.
\item \emph{The dilation of a compiled graph.} Let $\Recv$ be a
  receiver instrument whose instrument matrix is the
  compiled instrument matrix
  $\Gmat=D_\eta\Gmat^{\sigma}_{\Grph}$ of a well-formed graph $\Grph$ of
  catalogue elements for a linear extension
  $\sigma=(e_1,\dots,e_L)$ (Definition~\ref{def:compile}), with ports
  $\Port{1},\dots,\Port{K}$ and efficiencies $\eta_k$. For each $j$ let
  $\Theta_j$ be a loss factor of (i) of the element $e_j$, its rows indexed
  by a mode set $\Env_j$, and let
  $\hat\Theta_j:\mathbb{C}^{\Modes}\to\mathbb{C}^{\Env_j}$ act as $\Theta_j$
  on the modes of the support $\Sigma_{e_j}$ and as zero on the others. Let
  $\Theta_\eta\coloneqq\big(\sum_{k}\sqrt{1-\eta_k}\,\Emd{\{k\}}\big)^{\flat}$,
  its rows indexed by a mode set $\Env_\eta$, one for each mode of a port
  with $\eta_k<1$, the sets $\Modes$, $\Env_\eta$ and the $\Env_j$ being
  pairwise disjoint. With $\iota:\mathbb{C}^{\Inp}\to\mathbb{C}^{\Modes}$ the
  inclusion of the declared input modes, put $X_0\coloneqq\iota$ and
  $X_j\coloneqq\hat\Gmat_{e_j}X_{j-1}$ for $j=1,\dots,L$, and let
  \begin{equation*}\tag{\getrefnumber{C-eq:graph-dilation}}\digesteqlabel{eq:graph-dilation}
    \Omega:\mathbb{C}^{\Inp}\to\mathbb{C}^{\Modes\sqcup\Env_\eta\sqcup\Env_L\sqcup\dots\sqcup\Env_1}
  \end{equation*}
  have the rows $D_\eta X_L$ on $\Modes$, $\Theta_\eta X_L$ on $\Env_\eta$
  and $\hat\Theta_jX_{j-1}$ on $\Env_j$. Then $\Omega$, with the ports
  $\Port{1},\dots,\Port{K}$ of $\Recv$, is a port dilation of $\Recv$
  (Definition~\ref{def:port-dilation}), and its rows on $\Modes$ are those of
  $\Gmat_{\Inp}$.
\end{enumerate}
\end{digestclaim}

% claim receiver/component-catalogue/def:c10-tolerance
\digestlabel{def:c10-tolerance}{16.5}\pdfbookmark[2]{Definition 16.5 Admission test of the generic block}{digest-receiver/component-catalogue/def:c10-tolerance}
\renewcommand*{\RestatedName}{\hyperref{contract.pdf}{}{claim-receiver/component-catalogue/def:c10-tolerance}{Definition~16.5}}%
\begin{digestclaim}[Admission test of the generic block]
The (E1) predicate of catalogue entry C10 of
Definition~\ref{def:catalogue-entries} is decided by the route-(ii)
contractivity test of Definition~\ref{def:tolmap}\,(v) applied to the supplied
block $\Gmat_\Sigma$: the block is admitted when
$\lmax(\Gmat_\Sigma^\dagger\Gmat_\Sigma)\le1+\varepsilon_1$. C10 is the one
entry whose (E1) is a spectral condition rather than a parameter-domain
comparison (Theorem~\ref{thm:catalogue-admissibility}\,(iii)). Its lossless
sub-case $\Gmat_\Sigma^\dagger\Gmat_\Sigma=\id$ is an \emph{equality} and is
decided two-sidedly, as
$\lVert\Gmat_\Sigma^\dagger\Gmat_\Sigma-\id\rVert\le\varepsilon_1$.
\end{digestclaim}

\part*{Detector imperfections}\pdfbookmark[0]{Detector imperfections}{digest-part-4}

\section*{18\quad Detector efficiency}\pdfbookmark[1]{18 Detector efficiency}{digest-group-receiver/detector-imperfections/thm:efficiency}

% claim receiver/detector-imperfections/thm:efficiency
\digestlabel{thm:efficiency}{18.1}\pdfbookmark[2]{Theorem 18.1 Detector efficiency in the port Grams}{digest-receiver/detector-imperfections/thm:efficiency}
\begin{mainclaim}
\renewcommand*{\RestatedName}{\hyperref{contract.pdf}{}{claim-receiver/detector-imperfections/thm:efficiency}{Theorem~18.1}}%
\begin{digestclaim}[Detector efficiency in the port Grams]
Let $\Recv$ be a receiver instrument with port
Grams $\{A_k\}$ and an instrument matrix $\Gmat$ that does not yet
contain a detector-efficiency layer (for a declared element graph, the
product $\Gmat^{\sigma}_{\Grph}$ of Definition~\ref{def:compile}), let
$\eta_1,\dots,\eta_K\in[0,1]$, and let
$D_\eta=\sum_k\sqrt{\eta_k}\,\Emd{\{k\}}+\Emd{\Modes\setminus\bigcup_k\Port{k}}$
be the terminal detector-efficiency layer of
Definition~\ref{def:compile}.
\begin{enumerate}[label=(\roman*),leftmargin=3.0em]
\item \emph{Uniform within a port: a scalar on the Gram.} $D_\eta\Gmat$ is a
  contraction, and its port Grams are
  \begin{equation*}\tag{\getrefnumber{C-eq:etagram}}\digesteqlabel{eq:etagram}
    A_k\ \longmapsto\ \eta_k\,A_k .
  \end{equation*}
  For $\eta_k>0$ this preserves $\rk A_k$, so it leaves the one-mode-per-port
  criterion of Theorem~\ref{thm:representation-independence}\,(iv) untouched. Every formula of
  Theorems~\ref{thm:click-povm}, \ref{thm:coarse-graining}, \ref{thm:multiclick-bound}
  and~\ref{thm:emitted-bound} therefore covers inefficient detectors under the
  single substitution \eqref{eq:etagram}; in particular
  $A_\Sigma\mapsto\sum_k\eta_kA_k$, which is in general not a multiple of
  $A_\Sigma$ unless the $\eta_k$ agree.
\item \emph{Non-uniform within a port: not a scalar, and the rank can drop.} If
  attenuation varies \emph{within} a port, that is, for a diagonal contraction
  $D$ that is not constant on $\Port{k}$, the effective Gram is
  \begin{equation*}\tag{\getrefnumber{C-eq:noneffgram}}\digesteqlabel{eq:noneffgram}
    A_k^{\mathrm{eff}}=\Gmat_{\Inp}^\dagger D^\dagger\Emd{\{k\}}D\,
    \Gmat_{\Inp},
  \end{equation*}
  which is in general \emph{not} a multiple of $A_k$ and satisfies only
  $\rk A_k^{\mathrm{eff}}\le\rk A_k$, with strict inequality possible. All
  formulas of the theorems named in (i) still apply, to $A_k^{\mathrm{eff}}$.
\item \emph{One structural site.} Detector efficiency is a parameter of the
  terminal layer L1 and enters the compiled instrument matrix
  $D_\eta\Gmat^{\sigma}_{\Grph}$ only through the single left factor
  $D_\eta$ of Definition~\ref{def:compile}, which is applied after every chain
  element. It is therefore not an audit but an invariant of the compiler: there
  is exactly one site at which it can enter.
\item \emph{The composition record.} The emitted efficiency record
  lists every declared loss parameter $q_i$, under its name, with the element
  that declares it and its position in the compiled product: the
  transmittance $x$ of a beamsplitter (C1) and $t_H,t_V$ of a polarizing
  beamsplitter (C2), which is also the entry of a
  polarization-dependent beamsplitter, its $x_H,x_V$ recorded as
  $t_H,t_V$; the output losses $y_1,y_2$ of C1 and C2 and the loss $y$ of a
  retarder (C3), whose transmittances are $1-y_1$, $1-y_2$ and $1-y$; the
  reflectivities $R_H,R_V$ of a mirror (C5); the transmittances $\eta_m$ of a
  loss element (C7); the pass and leak transmittances $t,t_\perp$ of a
  polarizer (C8); and the detector efficiencies $\eta_k$. Two elements may
  carry equal numerical values; they are still two losses and appear as two
  rows.
\item \emph{The deficiency identity.} The identity \eqref{eq:deficiencyid} of
  Lemma~\ref{lem:deficiency} holds for every receiver instrument, before and
  after the layer $D_\eta$, and is the correct check that loss was accounted
  for exactly once: its two summands are the amplitude that left through the
  dilation and the amplitude that reached an unmonitored declared mode, and
  neither carries a separate efficiency term; after the layer, the efficiency
  loss is part of the dilation summand.
\end{enumerate}
\end{digestclaim}
\end{mainclaim}

\section*{19\quad Dark counts}\pdfbookmark[1]{19 Dark counts}{digest-group-receiver/detector-imperfections/def:darkadm}

% claim receiver/detector-imperfections/def:darkadm
\digestlabel{def:darkadm}{19.1}\pdfbookmark[2]{Definition 19.1 Dark-count channel and its independent product form}{digest-receiver/detector-imperfections/def:darkadm}
\renewcommand*{\RestatedName}{\hyperref{contract.pdf}{}{claim-receiver/detector-imperfections/def:darkadm}{Definition~19.1}}%
\begin{digestclaim}[Dark-count channel and its independent product form]
A \emph{dark-count channel}
is a matrix $R=(R_{S|S'})_{S,S'\subseteq[K]}$ that is
\begin{enumerate}[label=(D\arabic*),leftmargin=3.0em]
\item \emph{column-stochastic}: $R_{S|S'}\ge0$ and $\sum_SR_{S|S'}=1$ for every
  $S'$;
\item \emph{monotone}: $R_{S|S'}>0$ implies $S'\subseteq S$; a dark count
  adds a click and never removes one.
\end{enumerate}
The \emph{independent product form}
generated by per-port dark-count probabilities $p_{\dc|k}\in[0,1]$ is
\begin{equation*}\tag{\getrefnumber{C-eq:darkR}}\digesteqlabel{eq:darkR}
  R_{S|S'}\coloneqq
  \mathds{1}\big[S'\subseteq S\big]
  \prod_{k\in S\setminus S'}p_{\dc|k}
  \prod_{k\notin S}\big(1-p_{\dc|k}\big).
\end{equation*}
For $\mathcal{O}\subseteq2^{[K]}$ we write
$R_{\mathcal{O}|S'}\coloneqq\sum_{S\in\mathcal{O}}R_{S|S'}$.
A dark-count channel is supplied on one of two routes: on the
\emph{generated dark-count route}
it is \eqref{eq:darkR}, and on the
\emph{declared dark-count route}
it is a declared matrix $R$.
(D1) and (D2) hold by construction for \eqref{eq:darkR}
(Theorem~\ref{thm:dark-counts}\,(i)); for a \emph{declared} $R$ they are conditions to
be checked, and in floating point they are decided at the
receiver tolerances $\varepsilon_4$ and $\varepsilon_2$ respectively, by
item (iii) of Definition~\ref{def:tolmap}, whose content is
Definition~\ref{def:darktol}.
\end{digestclaim}

% claim receiver/detector-imperfections/thm:dark-counts
\digestlabel{thm:dark-counts}{19.2}\pdfbookmark[2]{Theorem 19.2 Dark counts on a photon-number block-diagonal POVM}{digest-receiver/detector-imperfections/thm:dark-counts}
\begin{mainclaim}
\renewcommand*{\RestatedName}{\hyperref{contract.pdf}{}{claim-receiver/detector-imperfections/thm:dark-counts}{Theorem~19.2}}%
\begin{digestclaim}[Dark counts on a photon-number block-diagonal POVM]
Let $R$ be a dark-count channel, let $\{M_S\}_{S\subseteq[K]}$ be the
click-pattern POVM of Theorem~\ref{thm:click-povm} or, at a setting $t$ of a
setting family of $K$ ports, its POVM $\{M_{S|t}\}_{S\subseteq[K]}$
(Definition~\ref{def:setting-family}\,(a)), written $\{M_S\}$, and put
$\widetilde{M}_S\coloneqq\sum_{S'\subseteq[K]}R_{S|S'}M_{S'}$.
\begin{enumerate}[label=(\roman*),leftmargin=3.0em]
\item \emph{The product form is admissible.} \eqref{eq:darkR} satisfies (D1)
  and (D2).
\item \emph{Structure is preserved.} $\{\widetilde M_S\}_{S\subseteq[K]}$ is a
  POVM, every $\widetilde M_S$ is block diagonal in the photon number with
  $\blk{\widetilde M}{\eqn}_S=\sum_{S'}R_{S|S'}\blk{M}{\eqn}_{S'}$, and $R$ acts
  on the outcome label alone: it commutes with every operation on the optical
  modes.
\item \emph{Multiclick absorption.} (D2) implies
  $R_{\MCK|S'}=1$ for every $S'\in\MCK$: dark counts never move weight out of
  the multiclick set.
\item \emph{The bound improves.} With
  \begin{equation*}\tag{\getrefnumber{C-eq:rmin}}\digesteqlabel{eq:rmin}
    r_{\min}\coloneqq\min_{S'\in2^{[K]}\setminus\MCK}R_{\MCK|S'} ,
  \end{equation*}
  the minimum running over \emph{every} non-multiclick column, the no-click
  column $S'=\varnothing$ included,
  one has, writing $\widetilde M_{\MCK}=\sum_{S\in\MCK}\widetilde M_S$,
  \begin{equation*}\tag{\getrefnumber{C-eq:darkimprove}}\digesteqlabel{eq:darkimprove}
    \widetilde M_{\MCK}\ \succeq\ (1-r_{\min})\,M_{\MCK}+r_{\min}\,\id ,
  \end{equation*}
  and consequently, if $\blk{M}{\gtN}_{\MCK}\succeq\pmc\,\id^{(>N)}$, then
  \begin{equation*}\tag{\getrefnumber{C-eq:darkpmc}}\digesteqlabel{eq:darkpmc}
    \blk{\widetilde M}{\gtN}_{\MCK}\ \succeq\
    \big((1-r_{\min})\,\pmc+r_{\min}\big)\,\id^{(>N)} ,
  \end{equation*}
  which is $\ge\pmc$ since $\pmc\le1$.
\end{enumerate}
\end{digestclaim}
\end{mainclaim}

% claim receiver/detector-imperfections/def:darktol
\digestlabel{def:darktol}{19.3}\pdfbookmark[2]{Definition 19.3 Receiver tolerances deciding dark-channel admissibility}{digest-receiver/detector-imperfections/def:darktol}
\renewcommand*{\RestatedName}{\hyperref{contract.pdf}{}{claim-receiver/detector-imperfections/def:darktol}{Definition~19.3}}%
\begin{digestclaim}[Receiver tolerances deciding dark-channel admissibility]
This definition extends the list of Definition~\ref{def:tolmap} by its
item~(iii), under the same rule: an exact condition about declared data
is decided in floating point by a receiver
tolerance, fixed here, and by no other threshold.

\emph{Dark-channel admissibility} (Definition~\ref{def:darkadm}) for a
\emph{declared} $R$: (D1) column-stochasticity is decided at
$\varepsilon_4$, admitted when $R_{S|S'}\ge-\varepsilon_4$ for every
entry and $\big|\sum_SR_{S|S'}-1\big|\le\varepsilon_4$ for every column;
and (D2) monotonicity is decided at $\varepsilon_2$: an entry counts as
positive when $R_{S|S'}>\varepsilon_2$, so the channel is refused exactly
when some entry with $S'\not\subseteq S$ exceeds $\varepsilon_2$.
$\varepsilon_4$ bounds the completeness residual \eqref{eq:res4}, which is
the same kind of quantity as a column sum's defect from one; $\varepsilon_2$
is the threshold at which this contract already calls a spectral quantity
zero.
\end{digestclaim}

% claim receiver/detector-imperfections/prop:dark-slack-bound
\digestlabel{prop:dark-slack-bound}{19.4}\pdfbookmark[2]{Proposition 19.4 Multiclick bound under a dark-count channel admitted with slack}{digest-receiver/detector-imperfections/prop:dark-slack-bound}
\renewcommand*{\RestatedName}{\hyperref{contract.pdf}{}{claim-receiver/detector-imperfections/prop:dark-slack-bound}{Proposition~19.4}}%
\begin{digestclaim}[Multiclick bound under a dark-count channel admitted with slack]
Let $\{M_S\}_{S\subseteq[K]}$ be the click-pattern POVM of
Theorem~\ref{thm:click-povm} or, at a setting $t$ of a setting family
of $K$ ports, its POVM $\{M_{S|t}\}_{S\subseteq[K]}$
(Definition~\ref{def:setting-family}\,(a)), written $\{M_S\}$, and for a matrix $R=(R_{S|S'})_{S,S'\subseteq[K]}$ let
$\widetilde M_S$, $\widetilde M_{\MCK}$ and $r_{\min}$ be formed as in
Theorem~\ref{thm:dark-counts} and \eqref{eq:rmin}. Put
\begin{equation*}\tag{\getrefnumber{C-eq:darkdelta}}\digesteqlabel{eq:darkdelta}
  \delta_{\dc}\ \coloneqq\
  \begin{cases}
    0
      & \text{generated route: }R\text{ is the product form }\eqref{eq:darkR},
      \\[3pt]
    \varepsilon_4+(K+1)\,\varepsilon_2
      & \text{declared route: }R\text{ admitted under
        Def.~\ref{def:tolmap}\,(iii),}
  \end{cases}
\end{equation*}
the two cases being the two routes of Definition~\ref{def:darkadm}, and
\begin{equation*}\tag{\getrefnumber{C-eq:rbar}}\digesteqlabel{eq:rbar}
  \bar r\ \coloneqq\ \min\{\,r_{\min},\,1\,\} .
\end{equation*}
Let $R$ be declared and admitted under item (iii) of
Definition~\ref{def:tolmap} (Definition~\ref{def:darktol}), on either side of
either test and whatever the sign of $r_{\min}$. Then:
\begin{enumerate}[label=(\roman*),leftmargin=3.0em]
\item \emph{(the multiclick columns)} for every $S'\in\MCK$,
  \begin{equation*}\tag{\getrefnumber{C-eq:d2slack}}\digesteqlabel{eq:d2slack}
    R_{\MCK|S'}\ \ge\ (1-\varepsilon_4)-(K+1)\varepsilon_2\ =\ 1-\delta_{\dc} ;
  \end{equation*}
\item \emph{(the upper side)} for every column $S'$,
  \begin{equation*}\tag{\getrefnumber{C-eq:rminupper}}\digesteqlabel{eq:rminupper}
    R_{\MCK|S'}\ \le\ (1+\varepsilon_4)+(K+1)\varepsilon_4\ =\ 1+(K+2)\varepsilon_4 ,
  \end{equation*}
  and consequently
  \begin{equation*}\tag{\getrefnumber{C-eq:rbarcost}}\digesteqlabel{eq:rbarcost}
    0\ \le\ r_{\min}-\bar r\ \le\ (K+2)\,\varepsilon_4 ;
  \end{equation*}
\item \emph{(the composed inequality)}
  \begin{equation*}\tag{\getrefnumber{C-eq:darkimproveslack}}\digesteqlabel{eq:darkimproveslack}
    \widetilde M_{\MCK}\ \succeq\
    (1-\delta_{\dc})\,M_{\MCK}+\bar r\,\big(\id-M_{\MCK}\big)
    \ \succeq\ \Big[(1-\bar r)\,M_{\MCK}+\bar r\,\id\Big]-\delta_{\dc}\,\id ,
  \end{equation*}
  and consequently, whenever $\blk{M}{\gtN}_{\MCK}\succeq\pmc\,\id^{(>N)}$,
  \begin{equation*}\tag{\getrefnumber{C-eq:darkpmcslack}}\digesteqlabel{eq:darkpmcslack}
    \blk{\widetilde M}{\gtN}_{\MCK}\ \succeq\
    \big((1-\bar r)\,\pmc+\bar r-\delta_{\dc}\big)\,\id^{(>N)} .
  \end{equation*}
\end{enumerate}
\end{digestclaim}

\section*{20\quad The dark-composed family and the noise floor}\pdfbookmark[1]{20 The dark-composed family and the noise floor}{digest-group-receiver/detector-imperfections/def:dark-composed}

% claim receiver/detector-imperfections/def:dark-composed
\digestlabel{def:dark-composed}{20.1}\pdfbookmark[2]{Definition 20.1 The dark-composed family}{digest-receiver/detector-imperfections/def:dark-composed}
\renewcommand*{\RestatedName}{\hyperref{contract.pdf}{}{claim-receiver/detector-imperfections/def:dark-composed}{Definition~20.1}}%
\begin{digestclaim}[The dark-composed family]
Let $\Tset$ be a finite non-empty set of settings and, for every $t\in\Tset$,
let $\{M_{S|t}\}_{S\subseteq[K]}$ be a POVM on a Fock space
$\mathcal{F}(\mathbb{C}^{I})$ (Definition~\ref{def:fock}) every element of
which commutes with the photon number. Let $R$ be a
dark-count channel (Definition~\ref{def:darkadm}) and $\cg$ a
declared coarse-graining onto $\Vset$ (Definition~\ref{def:coarse}).
The
\emph{dark-composed
family} is
\begin{equation*}\tag{\getrefnumber{C-eq:dark-composed}}\digesteqlabel{eq:dark-composed}
  \widetilde M_{S|t}\coloneqq\sum_{S'\subseteq[K]}R_{S|S'}\,M_{S'|t},
  \qquad
  \widetilde F_{v|t}\coloneqq\sum_{S\in\cg^{-1}(v)}\widetilde M_{S|t}
  =\sum_{S'\subseteq[K]}R_{\cg^{-1}(v)|S'}\,M_{S'|t},
\end{equation*}
for $S\subseteq[K]$, $v\in\Vset$ and $t\in\Tset$. Each
$\{\widetilde M_{S|t}\}_S$ and each $\{\widetilde F_{v|t}\}_v$ is a POVM
every element of which commutes with the photon number, with
$\blk{\widetilde F}{\eqn}_{v|t}=\sum_{S'}R_{\cg^{-1}(v)|S'}\blk{M}{\eqn}_{S'|t}$:
the coefficients are non-negative scalars, they sum to one over $S$ in every
column by (D1), and the fibres of $\cg$ partition $2^{[K]}$.
\end{digestclaim}

% claim receiver/detector-imperfections/cor:exact-dark-composed
\digestlabel{cor:exact-dark-composed}{20.2}\pdfbookmark[2]{Corollary 20.2 The dark-composed reference, enclosed or exact}{digest-receiver/detector-imperfections/cor:exact-dark-composed}
\renewcommand*{\RestatedName}{\hyperref{contract.pdf}{}{claim-receiver/detector-imperfections/cor:exact-dark-composed}{Corollary~20.2}}%
\begin{digestclaim}[The dark-composed reference, enclosed or exact]
Let $\{\widetilde M_{S|t}\}$ and $\{\widetilde F_{v|t}\}$ be the
dark-composed family of POVMs $\{M_{S|t}\}$ on
$\mathcal{F}(\mathbb{C}^{I})$ (Definition~\ref{def:dark-composed}), with a
dark-count channel $R$ whose entries are rational.
\begin{enumerate}[label=(\roman*),leftmargin=3.0em]
\item Suppose that for some $N\ge0$ every block $\blk{M}{\eqn}_{S|t}$ with
  $n\le N$, read in the occupation basis of $\Symn{n}(\mathbb{C}^{I})$ in an
  order fixed for that block, has a complex enclosure
  $\big(C^{(n)}_{S|t},r^{(n)}_{S|t}\big)$
  (Definition~\ref{def:complex-enclosure}). Then
  \begin{equation*}\tag{\getrefnumber{C-eq:exact-dark}}\digesteqlabel{eq:exact-dark}
    \Big(\sum_{S'\subseteq[K]}R_{\cg^{-1}(v)|S'}\,C^{(n)}_{S'|t},\
         \sum_{S'\subseteq[K]}R_{\cg^{-1}(v)|S'}\,r^{(n)}_{S'|t}\Big)
  \end{equation*}
  is a complex enclosure of $\blk{\widetilde F}{\eqn}_{v|t}$, exact on every
  entry on which every $r^{(n)}_{S'|t}$ with $R_{\cg^{-1}(v)|S'}\neq0$ is
  zero; the same holds for $\blk{\widetilde M}{\eqn}_{S|t}$, with $R_{S|S'}$
  in place of $R_{\cg^{-1}(v)|S'}$.
\item Made Hermitian-centred by
  Lemma~\ref{C-lem:enclosure-arithmetic}\,(v), the enclosures of (i) make
  $\{\widetilde F_{v|t}\}_v$, restricted to the blocks $n\le N$, a
  reference with constant $1$
  (\xref{foundations/def:exact-reference}) at every setting, against every
  sector map whose blocks are the photon numbers $n=0,\dots,N$ in
  increasing order, with any block kind and identifiers, each block read in
  its fixed order; it is an exact reference when every radius is zero.
  The same holds for $\{\widetilde M_{S|t}\}_S$.
\item For a receiver instrument with $\InpS=\varnothing$, the
  hypothesis of (i) holds with $\Tset=\{\off\}$, $M_{S|\off}=M_S$ and the
  enclosures of Theorem~\ref{thm:enclosed-blocks}\,(ii), at every $N$.
\item $R$ has rational entries on the generated dark-count route,
  whose $p_{\dc|k}$ are declared values, and on the declared route, whose
  entries are declared values; a declared value is the dyadic rational it
  denotes (\xref{foundations/def:binary64}\,(iii)).
\end{enumerate}
\end{digestclaim}

% claim receiver/detector-imperfections/prop:noise-floor
\digestlabel{prop:noise-floor}{20.3}\pdfbookmark[2]{Proposition 20.3 Loss and dark counts make every dark-composed element definite}{digest-receiver/detector-imperfections/prop:noise-floor}
\renewcommand*{\RestatedName}{\hyperref{contract.pdf}{}{claim-receiver/detector-imperfections/prop:noise-floor}{Proposition~20.3}}%
\begin{digestclaim}[Loss and dark counts make every dark-composed element definite]
Let $\{\widetilde M_{S|t}\}$ and $\{\widetilde F_{v|t}\}$ be the
dark-composed family (Definition~\ref{def:dark-composed}) of a
setting family $\big(\Tset,\{M_{S|t}\}\big)$ with operators $B^{(t)}$
and prefactors $\gamma^{(t)}$ (Definition~\ref{def:setting-family}), a
dark-count channel $R$ and a declared coarse-graining $\cg$, and write
$r_S\coloneqq R_{S|\varnothing}$.
\begin{enumerate}[label=(\roman*),leftmargin=3.0em]
\item \emph{The floor.} For every $t$, $S$, $v$ and $n\ge0$,
  \begin{equation*}\tag{\getrefnumber{C-eq:noise-floor}}\digesteqlabel{eq:noise-floor}
    \blk{\widetilde M}{\eqn}_{S|t}\succeq r_S\,\gamma^{(t)}
      \big(1-\lmax(B^{(t)})\big)^{n}\,\id,
    \qquad
    \blk{\widetilde F}{\eqn}_{v|t}\succeq R_{\cg^{-1}(v)|\varnothing}\,
      \gamma^{(t)}\big(1-\lmax(B^{(t)})\big)^{n}\,\id .
  \end{equation*}
  Hence, if $r_S>0$ for every $S\subseteq[K]$ and $\lmax(B^{(t)})<1$, every
  $\widetilde M_{S|t}$ and every $\widetilde F_{v|t}$ is positive definite on
  every photon-number block.
\item \emph{Dark counts in product form.} On the generated dark-count route,
  $r_S=\prod_{k\in S}p_{\dc|k}\prod_{k\notin S}(1-p_{\dc|k})$, and $r_S>0$
  for every $S$ if and only if $0<p_{\dc|k}<1$ for every $k$.
\item \emph{Loss through the efficiencies.} If the instrument matrix of a
  receiver instrument is $D_\eta\Gmat'$, with $\Gmat'$ a contraction and
  $D_\eta$ the terminal detector-efficiency layer, then
  $\lmax(A_\Sigma)\le\max_k\eta_k$. For the unconditioned receiver, with
  $\InpS=\varnothing$, $B=A_\Sigma$ and $\gamma=1$, so $\eta_k<1$ and
  $0<p_{\dc|k}<1$ for every $k$ give
  $\lmin\big(\blk{\widetilde M}{\eqn}_S\big)\ge r_S\,(1-\max_k\eta_k)^{n}>0$.
\item \emph{Loss is necessary.} If $\lmax(B^{(t)})=1$, then for every
  dark-count channel and every $n\ge1$ the element
  $\blk{\widetilde M}{\eqn}_{\varnothing|t}=r_\varnothing\,\blk{M}{\eqn}_{\varnothing|t}$
  is singular; if $B^{(t)}=\id$ it is zero.
\item \emph{Dark counts in every port are necessary.} For the unconditioned
  receiver on the generated route, let
  $Z\coloneqq\{k:p_{\dc|k}=0\}$. Then $\blk{\widetilde M}{\eqn}_S=0$ for
  every $S$ with $|S\cap Z|>n$; in particular
  $\blk{\widetilde M}{(=0)}_{\{k\}}=0$ for every $k\in Z$. If $p_{\dc|k}=1$,
  then $\widetilde M_S=0$ for every $S\not\ni k$.
\item \emph{Equivalence.} For the unconditioned receiver on the generated
  route, every $\widetilde M_S$ is positive definite on every
  photon-number block if and only if $\lmax(A_\Sigma)<1$ and
  $0<p_{\dc|k}<1$ for every $k$.
\end{enumerate}
\end{digestclaim}

\section*{21\quad Efficiency normalization}\pdfbookmark[1]{21 Efficiency normalization}{digest-group-receiver/detector-imperfections/def:pure-loss}

% claim receiver/detector-imperfections/def:pure-loss
\digestlabel{def:pure-loss}{21.1}\pdfbookmark[2]{Definition 21.1 Pure-loss channel on the channel modes}{digest-receiver/detector-imperfections/def:pure-loss}
\renewcommand*{\RestatedName}{\hyperref{contract.pdf}{}{claim-receiver/detector-imperfections/def:pure-loss}{Definition~21.1}}%
\begin{digestclaim}[Pure-loss channel on the channel modes]
Let $\InpC$ be the channel modes of a receiver instrument, so
that the channel system is $\Hs{\Chan}=\mathcal{F}(\mathbb{C}^{\InpC})$
(Definition~\ref{def:instrument}), let $\InpC'$ be a disjoint copy of
$\InpC$, with $m\mapsto m'$ the copying bijection, and let $t\in[0,1]$. Put
\begin{equation*}\tag{\getrefnumber{C-eq:lossiso}}\digesteqlabel{eq:lossiso}
  w_t\colon\mathbb{C}^{\InpC}\to\mathbb{C}^{\InpC\sqcup\InpC'},\qquad
  w_t\ket{m}\coloneqq\sqrt{t}\,\ket{m}+\sqrt{1-t}\,\ket{m'} ,
\end{equation*}
and identify $\mathcal{F}(\mathbb{C}^{\InpC\sqcup\InpC'})$ with
$\mathcal{F}(\mathbb{C}^{\InpC})\otimes\mathcal{F}(\mathbb{C}^{\InpC'})$ by
the occupation identification of Lemma~\ref{C-lem:fock-maps}\,(ii),
which is unitary. The
\emph{pure-loss channel}
of transmittance $t$ is
\begin{equation*}\tag{\getrefnumber{C-eq:pureloss}}\digesteqlabel{eq:pureloss}
  \Lambda^{\mathrm{loss}}_t(\rho)\;\coloneqq\;
  \tr_{\mathcal{F}(\mathbb{C}^{\InpC'})}\!\big[\Fock(w_t)\,\rho\,\Fock(w_t)^{\dagger}\big],
  \qquad \rho\in\Tcl{\Chan},
\end{equation*}
and its adjoint is
\begin{equation*}\tag{\getrefnumber{C-eq:purelossadj}}\digesteqlabel{eq:purelossadj}
  (\Lambda^{\mathrm{loss}}_t)^{\dagger}(Y)\;\coloneqq\;
  \Fock(w_t)^{\dagger}\big(Y\otimes\idon{\mathcal{F}(\mathbb{C}^{\InpC'})}\big)\Fock(w_t),
  \qquad Y\in\Lin{\Chan}.
\end{equation*}
Since $w_t^{\dagger}w_t=t\,\id+(1-t)\,\id=\id$, the map $\Fock(w_t)$ is an
isometry (Lemma~\ref{C-lem:fock-maps}\,(i)); so $\Lambda^{\mathrm{loss}}_t$ is completely positive and trace
preserving, $(\Lambda^{\mathrm{loss}}_t)^{\dagger}$ is completely positive and
unital, and
$\tr[Y\,\Lambda^{\mathrm{loss}}_t(\rho)]=\tr[(\Lambda^{\mathrm{loss}}_t)^{\dagger}(Y)\,\rho]$
for every bounded $Y$ and trace-class $\rho$, by the defining property of the
partial trace and cyclicity of the trace. At $t=1$ the channel is the
identity.
\end{digestclaim}

% claim receiver/detector-imperfections/def:normalized-station
\digestlabel{def:normalized-station}{21.2}\pdfbookmark[2]{Definition 21.2 Normalized station and factored loss}{digest-receiver/detector-imperfections/def:normalized-station}
\renewcommand*{\RestatedName}{\hyperref{contract.pdf}{}{claim-receiver/detector-imperfections/def:normalized-station}{Definition~21.2}}%
\begin{digestclaim}[Normalized station and factored loss]
Let $\Recv$ be a receiver instrument without
declared source modes, $\InpS=\varnothing$, with $K$ ports, whose
port Grams are
\begin{equation*}\tag{\getrefnumber{C-eq:efficiencygrams}}\digesteqlabel{eq:efficiencygrams}
  A_k=\eta_k\,A^{\circ}_k,\qquad k=1,\dots,K,
\end{equation*}
for an admissible Gram family $\{A^{\circ}_k\}$ on
$\mathbb{C}^{\InpC}$ and efficiencies $\eta_k\in[0,1]$ with
$\eta_{\max}\coloneqq\max_k\eta_k>0$. This is the case for an instrument
matrix $D_\eta\Gmat_0$, $\Gmat_0$ a contraction, with $A^{\circ}_k$ the port
Grams of $\Gmat_0$ (Theorem~\ref{thm:efficiency}\,(i)), and for a declared
Gram family $\{A^{\circ}_k\}$ with declared efficiencies. For
$t\in[\eta_{\max},1]$, a
\emph{normalized station}
of $\Recv$ at the
\emph{factored loss}
$t$ is a receiver instrument $\Recv'$ without declared source modes, with the
channel modes and the $K$ ports of $\Recv$, whose port Grams are
\begin{equation*}\tag{\getrefnumber{C-eq:normalizedgrams}}\digesteqlabel{eq:normalizedgrams}
  A'_k=\frac{\eta_k}{t}\,A^{\circ}_k ,\qquad\text{so that}\qquad A_k=t\,A'_k .
\end{equation*}
One exists: $\eta_k/t\in[0,1]$, so $0\preceq A'_k\preceq A^{\circ}_k$ and
$\sum_kA'_k\preceq\sum_kA^{\circ}_k\preceq\id$, and an admissible family is
the port-Gram family of some receiver instrument
(Theorem~\ref{thm:representation-independence}\,(ii)); for the instrument
matrix $D_\eta\Gmat_0$ one is $D_{\eta/t}\Gmat_0$, a product of two
contractions (Proposition~\ref{prop:mode-algebra}\,(ii)), with
$D_{\eta/t}$ the terminal layer of Definition~\ref{def:compile} at the
efficiencies $\eta_k/t$. Its click-pattern POVM depends on the port Grams
alone (Theorem~\ref{thm:representation-independence}\,(i)), so it is one
POVM for every choice. The choice $t=\eta_{\max}$ is the \emph{full}
normalization, at which the largest efficiency becomes one, and $t=1$ gives
$\Recv'$ the port Grams of $\Recv$.
\end{digestclaim}

% claim receiver/detector-imperfections/prop:efficiency-normalization
\digestlabel{prop:efficiency-normalization}{21.4}\pdfbookmark[2]{Proposition 21.4 A receiver as uniform loss followed by its normalized station}{digest-receiver/detector-imperfections/prop:efficiency-normalization}
\renewcommand*{\RestatedName}{\hyperref{contract.pdf}{}{claim-receiver/detector-imperfections/prop:efficiency-normalization}{Proposition~21.4}}%
\begin{digestclaim}[A receiver as uniform loss followed by its normalized station]
Let $\Recv$ be a receiver instrument without
declared source modes, with efficiencies $\eta_k$,
$\eta_{\max}>0$, as in Definition~\ref{def:normalized-station}, let
$t\in[\eta_{\max},1]$ and let $\Recv'$ be a normalized station of
$\Recv$ at the factored loss $t$. Let $\{M_S\}$ and $\{M'_S\}$,
$S\subseteq[K]$, be their click-pattern POVMs on
$\Hs{\Chan}=\mathcal{F}(\mathbb{C}^{\InpC})$ (Theorem~\ref{thm:click-povm}).
\begin{enumerate}[label=(\roman*),leftmargin=3.0em]
\item \emph{The click patterns.} For every $S\subseteq[K]$,
  \begin{equation*}\tag{\getrefnumber{C-eq:normalization}}\digesteqlabel{eq:normalization}
    M_S\;=\;(\Lambda^{\mathrm{loss}}_t)^{\dagger}\big(M'_S\big),
    \qquad\text{equivalently}\qquad
    \tr[M_S\,\rho]=\tr\big[M'_S\,\Lambda^{\mathrm{loss}}_t(\rho)\big]
    \ \ \text{for every }\rho\in\Tcl{\Chan}.
  \end{equation*}
\item \emph{Classical post-processing, dark counts included.} For every
  column-stochastic matrix $R$ on $2^{[K]}$, in particular every
  dark-count channel, and every declared coarse-graining $\cg$
  onto $\Vset$,
  \begin{equation*}\tag{\getrefnumber{C-eq:normalizationdark}}\digesteqlabel{eq:normalizationdark}
    \sum_{S'}R_{S|S'}M_{S'}=(\Lambda^{\mathrm{loss}}_t)^{\dagger}
    \Big(\sum_{S'}R_{S|S'}M'_{S'}\Big),
    \qquad
    \sum_{S'}R_{\cg^{-1}(v)|S'}M_{S'}=(\Lambda^{\mathrm{loss}}_t)^{\dagger}
    \Big(\sum_{S'}R_{\cg^{-1}(v)|S'}M'_{S'}\Big),
  \end{equation*}
  so the dark-composed family of $\Recv$ with $R$ and $\cg$
  (Definition~\ref{def:dark-composed}) is the pull-back of that of $\Recv'$
  with the same $R$ and $\cg$: the dark counts act after the factored loss.
\item \emph{The device.} Under the receiver's device assumptions,
  Assumptions~\ref{ass:trust}, \ref{ass:passive}, \ref{ass:vac},
  \ref{ass:modes}, \ref{ass:thresh}, \ref{ass:port}, \ref{ass:exact},
  \ref{ass:nosel} and~\ref{ass:round}, $\{M_S\}$ is the measurement the
  physical receiver makes on the light it receives in its channel modes, and,
  by Assumptions~\ref{ass:modes} and~\ref{ass:port}, every port responds
  alike to every mode through which light reaches it; the physical receiver,
  followed by any dark-count channel, therefore acts on that light exactly as
  $\Lambda^{\mathrm{loss}}_t$ followed by $\Recv'$ and the same dark-count
  channel.
\end{enumerate}
\end{digestclaim}

\section*{22\quad Requirements}\pdfbookmark[1]{22 Requirements}{digest-group-receiver/detector-imperfections/req:efficiency-schema}

% claim receiver/detector-imperfections/req:efficiency-schema
\digestlabel{req:efficiency-schema}{22.1}\pdfbookmark[2]{Requirement 22.1 Detector efficiency and component loss in separate fields}{digest-receiver/detector-imperfections/req:efficiency-schema}
\renewcommand*{\RestatedName}{\hyperref{contract.pdf}{}{claim-receiver/detector-imperfections/req:efficiency-schema}{Requirement~22.1}}%
\begin{digestrequirement}[Detector efficiency and component loss in separate fields]
A C7 element is not a detector efficiency, and the engine's input schema keeps
the two in different fields.
\end{digestrequirement}

% claim receiver/detector-imperfections/req:efficiency-record
\digestlabel{req:efficiency-record}{22.2}\pdfbookmark[2]{Requirement 22.2 The efficiency record lists every declared loss parameter}{digest-receiver/detector-imperfections/req:efficiency-record}
\renewcommand*{\RestatedName}{\hyperref{contract.pdf}{}{claim-receiver/detector-imperfections/req:efficiency-record}{Requirement~22.2}}%
\begin{digestrequirement}[The efficiency record lists every declared loss parameter]
The emitted \iface{efficiency_record} lists every declared loss parameter
$q_i$ that Theorem~\ref{thm:efficiency}\,(iv) names, by catalogue entry and
parameter, with the element that declares it and its position in the compiled
product of Definition~\ref{def:compile}. Every row names its parameter beside
its value, which is read as that parameter: a transmittance, a reflectivity,
or, for $y$, $y_1$ and $y_2$, a loss whose transmittance is $1-q_i$. Two
elements may carry equal numerical values; they are still two losses and
appear as two rows. A polarizer's pass transmittance $t$ is one such row,
under its element; the factored loss $t$ that the same record carries by
Requirement~\ref{req:normalization-record} is not a row and is recorded under
its own name.
\end{digestrequirement}

% claim receiver/detector-imperfections/req:efficiency-convention
\digestlabel{req:efficiency-convention}{22.3}\pdfbookmark[2]{Requirement 22.3 Declared efficiency convention of the port Grams}{digest-receiver/detector-imperfections/req:efficiency-convention}
\renewcommand*{\RestatedName}{\hyperref{contract.pdf}{}{claim-receiver/detector-imperfections/req:efficiency-convention}{Requirement~22.3}}%
\begin{digestrequirement}[Declared efficiency convention of the port Grams]
Which of the two conventions a declaration means, Grams from power data
rescaled by $\eta_k$ or Grams from click data with every $\eta_k$ declared as
$1$, is a declared field, and the \iface{efficiency_record}
(Requirement~\ref{req:efficiency-record}) is where the answer is written down.
\end{digestrequirement}

% claim receiver/detector-imperfections/req:efficiency-normalization
\digestlabel{req:efficiency-normalization}{22.4}\pdfbookmark[2]{Requirement 22.4 Compilation of a probe-free receiver at its normalized station}{digest-receiver/detector-imperfections/req:efficiency-normalization}
\renewcommand*{\RestatedName}{\hyperref{contract.pdf}{}{claim-receiver/detector-imperfections/req:efficiency-normalization}{Requirement~22.4}}%
\begin{digestrequirement}[Compilation of a probe-free receiver at its normalized station]
For a request whose receiver declares no source modes and some efficiency
$\eta_k>0$, the engine compiles, in place of the declared receiver, its
normalized station at the factored loss $t$, with the efficiency of port $k$
equal to $\eta_k/t$, and emits that station's blocks, multiclick weights and
dark-composed family; $t$ is the value of \iface{efficiency_normalization}
where the request declares one; otherwise, with $\eta_{\max}$ the largest
declared efficiency, $t=\eta_{\max}/(99/100)$, formed exactly, where
$\eta_{\max}<99/100$, and $t=1$ where $\eta_{\max}\ge 99/100$. On every route the station compiled is the one at
the exact rational efficiencies $\eta_k/t$: its exact statements, its
reference and its enclosures are formed from those rationals, its
floating-point evaluation evaluates them as the evaluation model
evaluates every declared number, and no value rounded from $\eta_k/t$ is
declared in their place; the quotient is exactly one at a port with
$\eta_k=t$. A declared $t$ outside
$[\eta_{\max},1]$ is refused with
\iface{efficiency_normalization_out_of_range}. A request that declares source
modes, or whose efficiencies are all zero, is compiled at $t=1$, and a
declared $t<1$ in such a request is refused with the same reason.
\end{digestrequirement}

% claim receiver/detector-imperfections/req:normalization-record
\digestlabel{req:normalization-record}{22.5}\pdfbookmark[2]{Requirement 22.5 The efficiency record names the normalization}{digest-receiver/detector-imperfections/req:normalization-record}
\renewcommand*{\RestatedName}{\hyperref{contract.pdf}{}{claim-receiver/detector-imperfections/req:normalization-record}{Requirement~22.5}}%
\begin{digestrequirement}[The efficiency record names the normalization]
The emitted \iface{efficiency_record} carries the factored loss $t$, whether
$t$ was declared in the request or set by
Requirement~\ref{req:efficiency-normalization}, every declared efficiency
$\eta_k$ beside the exact rational $\eta_k/t$ at which the receiver was
compiled, and whether the normalization was applied, which it was exactly
where $t<1$. Where it was not, it carries the reason, the first that holds of:
source modes declared, every efficiency zero, or the factor one, declared or
by default. Where $t<1$ the artefact states that its blocks, multiclick
weights and dark-composed family are those of the normalized station at $t$,
whose pull-back through the pure-loss channel of transmittance $t$ is the
receiver's measurement.
\end{digestrequirement}

% claim receiver/detector-imperfections/req:loss-before-evaluation
\digestlabel{req:loss-before-evaluation}{22.6}\pdfbookmark[2]{Requirement 22.6 The factored loss is applied before the artefact is evaluated}{digest-receiver/detector-imperfections/req:loss-before-evaluation}
\renewcommand*{\RestatedName}{\hyperref{contract.pdf}{}{claim-receiver/detector-imperfections/req:loss-before-evaluation}{Requirement~22.6}}%
\begin{digestrequirement}[The factored loss is applied before the artefact is evaluated]
Where a receiver artefact's \iface{efficiency_record} carries a factored loss
$t<1$, every evaluation of its emitted blocks or dark-composed
family, or of a target formed from them, on a state that stands
for the receiver's statistics on an input $\rho$ of the channel system (an
honest-statistics provider forming the honest law among them) is made on
$\Lambda^{\mathrm{loss}}_t(\rho)$, or on the state the same construction forms
from $\Lambda^{\mathrm{loss}}_t(\rho)$, in place of $\rho$, with the $t$ of that
receiver's record. Where the input is a joint state of which that receiver's
channel system is one factor, such as a state shared with another party or
with a second normalized receiver, $\Lambda^{\mathrm{loss}}_t$ stands here for
$\idmap\otimes\Lambda^{\mathrm{loss}}_t$ acting on that receiver's
channel-system factor, each receiver with the $t$ of its own record. The states the
key-rate program optimises over are not inputs in this sense, and the loss
is not applied to them. A consumer that cannot apply $\Lambda^{\mathrm{loss}}_t$ to
the input it is handed refuses with \iface{factored_loss_not_applied}.
\end{digestrequirement}

% claim receiver/detector-imperfections/req:dark-route-exclusive
\digestlabel{req:dark-route-exclusive}{22.7}\pdfbookmark[2]{Requirement 22.7 Refusal of a dark-count channel supplied on both routes}{digest-receiver/detector-imperfections/req:dark-route-exclusive}
\renewcommand*{\RestatedName}{\hyperref{contract.pdf}{}{claim-receiver/detector-imperfections/req:dark-route-exclusive}{Requirement~22.7}}%
\begin{digestrequirement}[Refusal of a dark-count channel supplied on both routes]
The two routes of Definition~\ref{def:darkadm} are exclusive: a request
declaring $R$ while some $p_{\dc|k}\neq0$ is refused with
\iface{dark_count_channel_declared_twice}.
\end{digestrequirement}

% claim receiver/detector-imperfections/req:dark-channel-record
\digestlabel{req:dark-channel-record}{22.8}\pdfbookmark[2]{Requirement 22.8 The emitted dark-count channel and the scalars of its composed bound}{digest-receiver/detector-imperfections/req:dark-channel-record}
\renewcommand*{\RestatedName}{\hyperref{contract.pdf}{}{claim-receiver/detector-imperfections/req:dark-channel-record}{Requirement~22.8}}%
\begin{digestrequirement}[The emitted dark-count channel and the scalars of its composed bound]
The artifact's \iface{dark_count_channel} entry carries $R$, which of the two
routes of Definition~\ref{def:darkadm} it came from
(Remark~\ref{C-rem:corrdark}), and $r_{\min}$, $\bar r$ and $\delta_{\dc}$ under
distinct names. Nothing is clamped: $R$ is emitted as declared and $r_{\min}$
as \eqref{eq:rmin} computes it. \eqref{eq:darkpmcslack}, \emph{not}
\eqref{eq:darkpmc}, is the inequality a downstream module composes with.
\end{digestrequirement}

% claim receiver/detector-imperfections/req:dark-composed-reference
\digestlabel{req:dark-composed-reference}{22.9}\pdfbookmark[2]{Requirement 22.9 The emitted blocks are dark-composed and carry their reference}{digest-receiver/detector-imperfections/req:dark-composed-reference}
\renewcommand*{\RestatedName}{\hyperref{contract.pdf}{}{claim-receiver/detector-imperfections/req:dark-composed-reference}{Requirement~22.9}}%
\begin{digestrequirement}[The emitted blocks are dark-composed and carry their reference]
\begin{enumerate}[label=(\roman*),leftmargin=2.4em]
\item The \iface{photon_blocks} entry carries, for every probe setting $t$
  and every $n\le N$, the blocks
  $\blk{\widetilde F}{\eqn}_{v|t}=\sum_{S'\subseteq[K]}R_{\cg^{-1}(v)|S'}\,\blk{M}{\eqn}_{S'|t}$
  given by the formula \eqref{eq:dark-composed} for the emitted matrix $R$, in
  floating point. They are the blocks of the dark-composed family
  (Definition~\ref{def:dark-composed}) when $R$ is a dark-count channel,
  satisfying (D1) and (D2) exactly; for a declared $R$ admitted only at the
  receiver tolerances they are the same formula evaluated at that $R$.
\item For a declaration with a model admissible exactly
  (Requirement~\ref{req:exact-route-decision}) whose $R$ satisfies (D1) and
  (D2) exactly, decided in exact arithmetic on its declared entries on the
  declared dark-count route, the entry also carries, for every setting
  and every $n\le N$, a complex enclosure with a Hermitian centre of
  every block $\blk{\widetilde F}{\eqn}_{v|t}$
  (Definition~\ref{def:complex-enclosure}), computed in exact rational
  arithmetic on the exact island
  (\xref{foundations/def:exact-island}\,(i)) and with no binary64 value as an
  input: from the enclosures of the port Grams of
  Requirement~\ref{req:exact-route-decision}\,(iii), by
  Theorem~\ref{thm:enclosed-blocks}\,(ii) for an unconditioned receiver and
  by the enclosure of its blocks that a submodule constructing the setting
  family states for that family, then by Corollary~\ref{cor:exact-dark-composed}\,(i) and
  Lemma~\ref{C-lem:enclosure-arithmetic}\,(v). The
  square-root enclosures and
  angle enclosures used may be any that
  Definition~\ref{def:complex-enclosure} admits. The following entries are
  enclosed exactly, as zero with radius zero:
  \begin{enumerate}[label=(\alph*),leftmargin=2.0em]
  \item every entry of a block $\blk{\widetilde F}{\eqn}_{v|t}$ such that
    every $\blk{\widetilde M}{\eqn}_{S|t}$ with $S\in\cg^{-1}(v)$, whose sum
    it is (Definition~\ref{def:dark-composed}), is shown to vanish: for the
    unconditioned receiver on the generated dark-count route, and
    for no other receiver or route, by
    Proposition~\ref{prop:noise-floor}\,(v), decided on its declared
    per-port dark-count probabilities; or, for $S=\varnothing$ and $n\ge1$, by
    Proposition~\ref{prop:noise-floor}\,(iv) at a setting whose operator
    $B^{(t)}$ the engine has computed in exact arithmetic and found to be the
    identity, which for the unconditioned receiver is $A_\Sigma$ computed
    from the port Grams of
    Requirement~\ref{req:exact-route-decision}\,(i) or (ii);
  \item every entry whose enclosure is exact by
    Theorem~\ref{thm:exact-blocks}\,(ii) or by the exactness of its blocks
    that a submodule constructing the setting family states, carried through
    Corollary~\ref{cor:exact-dark-composed}\,(i), and whose value is zero.
  \end{enumerate}
  No other entry is required to be enclosed as zero.
\item The enclosures of (ii), with the family they enclose, are the
  reference (\xref{foundations/def:exact-reference}) that the receiver
  offers for a snap of its emitted family
  (\xref{foundations/req:snap-reference}), one per setting, with constant $1$
  (Corollary~\ref{cor:exact-dark-composed}\,(ii)). The entry records the
  \emph{offered} reference kind of every block: \emph{exact} where every
  radius of the block is zero, and \emph{enclosed} otherwise, with the
  largest radius of the block and the largest radius of the enclosures of
  irrational inputs (square roots, trigonometric values) from which the
  centres are computed exactly. The record of every snap of the family
  carries besides the kind of the reference the snap ran on
  (\xref{foundations/req:who-snaps}): \emph{exact} on a block whose
  offered kind is exact, the snap running on its centres, and \emph{none}
  on every other block, the snap running on the floating-point blocks of
  (i). The reference is an
  exact reference where every block is exact. It also names the claims that establish the
  enclosure, positive semidefiniteness and completeness of the reference,
  which \xref{foundations/req:snap-reference}\,(i) leaves to the
  constructing module: Corollary~\ref{cor:exact-dark-composed}\,(ii), with
  Theorem~\ref{thm:enclosed-blocks}\,(ii) through
  Corollary~\ref{cor:exact-dark-composed}\,(iii) for an unconditioned
  receiver and the enclosure claim of the constructing submodule for a
  setting family a submodule constructs. On the exact route, with the
  port Grams enclosed exactly, every block $n\le1$ has radius zero
  (Theorem~\ref{thm:exact-blocks}\,(iii) for an unconditioned receiver, the
  exactness claim of the constructing submodule for a setting family a
  submodule constructs, and Corollary~\ref{cor:exact-dark-composed}\,(i)). The reference is read against
  the sector map of the floating-point blocks of (i) in the same entry
  (\xref{foundations/def:sector-map}): the same block kind and block
  identifiers, the photon numbers $n=0,\dots,N$ in increasing order, and
  within each block the occupation vectors in the one order in which the
  entry lists the floating-point blocks, which is also the order of the rows
  and columns of every centre and radius matrix. The entries of the reference
  that a snap treats as known to be exactly zero
  (\xref{foundations/def:snap}) are therefore exactly those enclosed as zero
  with radius zero.
\item Where the declaration has no model admissible exactly, where a declared
  $R$ fails (D1) or (D2) exactly, or where a reciprocal that the enclosure
  of a setting family a submodule constructs forms receives no enclosure, the
  entry carries no reference, records the offered reference kind of every
  block as \emph{none}, and records the cause.
\item $R$ remains emitted in \iface{dark_count_channel}
  (Requirement~\ref{req:dark-channel-record}), where it serves the
  high-photon weight. It is already composed into the blocks of (i)--(iii),
  and a consumer does not compose it into them again.
\item Items (ii) and (iii) bind every run, save one that declares that
  enclosed references are not used
  (\xref{foundations/req:snap-reference}\,(i)). Such a run need form no
  enclosure of (ii) with a radius that is not zero, beyond what deciding that
  a block is not enclosed exactly takes, and carries none: it carries the
  reference of (iii) only on the blocks every entry of which it encloses
  exactly, with radius zero, from port Grams enclosed exactly, each block so
  carried being an exact reference of that block, and records their
  offered kind \emph{exact}; every other block it records as offered with
  the kind \emph{none}, with the cause that enclosed references are not
  used (\xref{foundations/req:snap-reference}\,(iii)). The reference the entry
  carries, its centres and radii as data, is the handed reference
  that the instance hands beside its snap of the emitted family
  (\xref{foundations/def:evaluated-off-snap}\,(i)), whatever the snap itself
  uses of it: the snap uses only the centres of the exact blocks
  (\xref{foundations/req:snap-reference}\,(ii)), and the run's first-order
  bound reads the reference on every block it carries
  (\xref{meat-core/frank-wolfe/req:bracket-on-reference}). Item (iv) binds
  as stated.
\item \emph{Factor enclosures.} Where the entry carries the reference of
  (iii) for a receiver with $\InpS=\varnothing$, whose one setting is $\off$
  (Corollary~\ref{cor:exact-dark-composed}\,(iii)), the engine forms
  besides, on the blocks the count below admits and on no other, a
  factor enclosure (\xref{foundations/def:factor-enclosure}) of the
  block $\blk{\widetilde F}{\eqn}_{v|\off}$: the one of
  Theorem~\ref{thm:factor-blocks}\,(iii) at the weights
  $w_{S'}\coloneqq R_{\cg^{-1}(v)|S'}$, rationals by
  Corollary~\ref{cor:exact-dark-composed}\,(iv), at which the block of that
  theorem is $\blk{\widetilde F}{\eqn}_{v|\off}$
  (Definition~\ref{def:dark-composed}), its rows in the order of (iii), for
  the port dilation (Definition~\ref{def:port-dilation}) formed as follows.
  \begin{enumerate}[label=(\alph*),leftmargin=2.0em]
  \item On the route that declares a component graph, the dilation of
    Lemma~\ref{lem:catalogue-dilations}\,(ii), with the loss factor of
    part (i) of that lemma for every element. A row of a loss factor, and of
    the efficiency layer's, is formed only where the declared value that
    scales it is decided in exact arithmetic to make it non-zero, so that
    whether an element is lossless is read off its declared parameters and
    never off a computed number. The loss factor of a generic block is formed
    from a decomposition
    $\id_\Sigma-\Gmat_\Sigma^{\dagger}\Gmat_\Sigma=E^{\dagger}DE$ formed and
    decided in exact rational arithmetic; where none is formed, no dilation of
    the declaration is.
  \item On a route that declares the port Grams or the instrument matrix, the
    dilation of Theorem~\ref{thm:factor-blocks}\,(iv), where the engine forms
    decompositions $A_k=E_k^{\dagger}D_kE_k$ and $\defic=E_0^{\dagger}D_0E_0$
    of the kind named there, formed and decided in exact rational arithmetic.
  \end{enumerate}
  The entries of the dilation are enclosed by the rules of
  Lemma~\ref{C-lem:enclosure-arithmetic} from square-root enclosures and angle
  enclosures, as Requirement~\ref{req:exact-route-decision}\,(iii) encloses
  the elements' matrices and the $\sqrt{\eta_k}$, and the columns as
  Theorem~\ref{thm:factor-blocks}\,(ii) encloses them, in exact rational
  arithmetic or by a rigorously enclosed computation
  (\xref{foundations/def:tiers}), with no binary64 value as an input; a column
  enclosed exactly as zero is omitted, and the core is exact. The engine forms
  the factor enclosure of a block only where the number of occupation vectors
  $\vec m$ of weight $n$ on the modes of the dilation whose row is not
  enclosed exactly as zero, with $R_{\cg^{-1}(v)|S_{\vec m}}>0$, is at most
  the block's dimension $\binom{|\Inp|+n-1}{n}$, a count it decides before it
  computes any column: a factor enclosure with more columns than that
  dimension fixes no rank
  (\xref{foundations/cor:rank-over-factor-enclosure}\,(ii)). It forms none
  for a receiver with $\InpS\neq\varnothing$, whose blocks are not formed
  through a port dilation; none where no dilation of (a) or (b) is formed; none on a block
  the count does not admit; and none in a run that declares that enclosed
  references are not used (vi). A factor enclosure is carried beside the
  entrywise enclosure of its block, which it leaves unaltered, and is handed
  with the reference, altered by no step. The entry records per block whether
  one is carried, with its number of columns and its largest radius, or the
  cause where none is; and it names Theorem~\ref{thm:factor-blocks} and, on
  the route of (a), Lemma~\ref{lem:catalogue-dilations} as the claims that
  establish that the block is a member of its factor enclosure, which
  \xref{foundations/req:snap-reference}\,(i) leaves to the module that
  constructs the reference.
\end{enumerate}
\end{digestrequirement}

% claim receiver/detector-imperfections/req:snap-fix
\digestlabel{req:snap-fix}{22.10}\pdfbookmark[2]{Requirement 22.10 Cause and fix recorded for a receiver block resting on its snap}{digest-receiver/detector-imperfections/req:snap-fix}
\renewcommand*{\RestatedName}{\hyperref{contract.pdf}{}{claim-receiver/detector-imperfections/req:snap-fix}{Requirement~22.10}}%
\begin{digestrequirement}[Cause and fix recorded for a receiver block resting on its snap]
When a moved block of a snap of the receiver's emitted dark-composed blocks
(Requirement~\ref{req:dark-composed-reference}) rests on
\xref{foundations/ass:snapped-measurement} under
\xref{foundations/req:snap-reference}\,(iii), its record names, beside the
setting, the block with its cause and the group that requirement names for
it, the reference kind of the snap, as \xref{foundations/req:who-snaps}
defines it (\emph{exact}, where the snap ran on the centres of an exact
reference of the block, or \emph{none}), the kind of the reference handed on
the block, if any (\emph{exact}, every radius of the block zero, or
\emph{enclosed}, with the largest radius of the block and the largest radius
of the enclosures of irrational inputs (square roots, trigonometric values)
from which the centres were computed exactly), and the fix below for the
cause, which details the fix of its group. A run that declares the strict
option of \xref{foundations/req:snap-reference}\,(iv) refuses instead,
naming the same; no other run refuses for these causes, save where the last
resort of \xref{foundations/req:support-safe-snap}\,(iv) yields no snap.
\begin{enumerate}[label=(\roman*),leftmargin=2.4em]
\item \emph{No reference.} The cause recorded by
  Requirement~\ref{req:dark-composed-reference}\,(iv), in the group
  \emph{no usable reference}, and the fix for it:
  \begin{enumerate}[label=(\alph*),leftmargin=2.0em]
  \item for declared port Grams that are not Hermitian, a declared family that
    is not an admissible Gram family exactly, or a generic block that is not a
    contraction exactly (Requirement~\ref{req:exact-route-decision}\,(i),
    (ii)), declared data that are Hermitian and admissible exactly, in
    rational arithmetic and not only at the receiver tolerances;
  \item for a declared dark-count matrix that does not satisfy (D1) and (D2)
    exactly, one that is exactly column-stochastic and monotone, or the
    generated dark-count route;
  \item for a reciprocal that the enclosure of a setting family a submodule
    constructs forms and that receives no enclosure, square-root and angle
    enclosures with smaller
    radii (Definition~\ref{def:complex-enclosure}).
  \end{enumerate}
\item \emph{Enclosed references not used.} Where the run declares that
  enclosed references are not used and the entry offers an enclosed one, the
  block is handed no reference, in the group \emph{no usable reference}, and
  the fix names the kind the entry offers for every block
  (Requirement~\ref{req:dark-composed-reference}\,(iii)), with the largest
  radius of every block offered as enclosed, and the alternative of an exact
  reference, which the receiver offers for a declaration on the exact route
  at a cutoff $N\le1$ when its port Grams are enclosed exactly
  (Theorem~\ref{thm:exact-blocks}\,(iii),
  Corollary~\ref{cor:exact-dark-composed}\,(ii)), or a run that uses
  enclosed references.
\end{enumerate}
A block that the snap makes by the last resort of
\xref{foundations/req:support-safe-snap}\,(iv) is flagged
\iface{support_not_preserved} in the record, beside its cause and reference
kind, and where it rests, in the group \emph{last resort}, the fix names the
ports whose declared efficiency is $1$ where $\lmax(A_\Sigma)=1$, and the
click patterns $S$ whose entry $R_{S|\varnothing}$ of the
dark-count channel is $0$, since Proposition~\ref{prop:noise-floor}\,(i)
needs $R_{S|\varnothing}>0$ for every $S$: on the
declared dark-count route it asks for a declared $R$ with
$R_{S|\varnothing}>0$ for every $S$, and on the
generated dark-count route, where every $R_{S|\varnothing}$ is positive
exactly when $0<p_{\dc|k}<1$ for every $k$
(Proposition~\ref{prop:noise-floor}\,(ii)), a dark-count probability in
$(0,1)$ in every port; in either case a model with loss, at the smallest loss
and dark-count rates the user is sure the device has
(Proposition~\ref{prop:noise-floor}); and, where the entry offers no
reference, the fix of (i) as well. Where every element of a block of the
reference is positive definite (Proposition~\ref{prop:noise-floor}\,(i)), no
row of it is structurally zero (\xref{foundations/def:snap}), so every index
lies in the structural support of every element and the uniform lift
of that item meets the rules (i) to (iii) of
\xref{foundations/req:support-safe-snap} on the block. For a block whose
cause is the run's, a handed reference on which the companion's bound is
withdrawn or not evaluated, or counts while the key length is formed from a
number certified on the snapped declaration in its place, the cause
\iface{companion_not_certified}
(\xref{meat-core/frank-wolfe/req:certified-evaluation}\,(i)), the record of
the run names the fix
\xref{foundations/req:snap-reference}\,(iii) names for that cause, and, for
such a block snapped by the last resort, which rests in the group
\emph{last resort}, that fix beside the one above; the run decides that cause
after its search, so the fix is completed there and not at the snap. A
handed block that the run evaluates off its snap rests on nothing and is
named with no fix, whether or not the last resort snapped it. A reference that fails
the decisions of \xref{foundations/req:snap-reference}\,(i) cannot arise for
the receiver: the centres are Hermitian with entries in $\mathbb{Q}(i)$ and
the radii symmetric non-negative rationals by construction
(Lemma~\ref{C-lem:enclosure-arithmetic}\,(v)), and the enclosure, positive
semidefiniteness and completeness that the constructing module establishes
are Corollary~\ref{cor:exact-dark-composed}\,(ii), whose hypothesis
Theorem~\ref{thm:enclosed-blocks}\,(ii), through
Corollary~\ref{cor:exact-dark-composed}\,(iii), and, for a setting family a
submodule constructs, that submodule's enclosure of its blocks supply.
\end{digestrequirement}

\bibliographystyle{plain}
\bibliography{../../../refs/references}

\end{document}
